% Generated mechanically from frozen input canon/ko-kr/integration-20260908-r10-p03-register/inputs/p03/ko/sdga.tex.
% Do not edit this generated file; edit the bound locale source instead.

% OK, start here.
%

\setcounter{chapter}{23}
\chapter{미분 등급 층}



\phantomsection
\label{sdga-section-phantom}


\section{서론}
\label{sdga-section-introduction}

\noindent
이 장에서는 미분 등급 대수, 절 \ref{dga-section-introduction}에서 시작한
논의를 이어간다. 개관 논문으로는 \cite{Keller-survey}가 있다.




\section{규약}
\label{sdga-section-conventions}

\noindent
이 장에서도 {\it 환}은 $1$을 갖는 가환환을 뜻한다는 규약을 따른다.
$R$이 환일 때, {\it $R$-대수 $A$}란 결합적이고 항등원을 갖는 곱셈,
즉 $R$-쌍선형 사상 $A \times A \to A$가 주어진 $R$-가군 $A$를 뜻한다.
달리 말하면, 이는 구조 사상 $R \to A$의 상이 $A$의 중심에 들어가는
단위원적 결합 $R$-대수이다.








\section{등급 대수의 층}
\label{sdga-section-ga}

\noindent
이 절은 건너뛰어도 된다.

\begin{definition}
\label{sdga-definition-ga}
$(\mathcal{C}, \mathcal{O})$를 환 달린 사이트라 하자.
$(\mathcal{C}, \mathcal{O})$ 위의
{\it 등급 $\mathcal{O}$-대수의 층}, 또는 {\it 등급 대수의 층}은
$n \in \mathbf{Z}$를 첨자로 하는 $\mathcal{O}$-가군의 족
$\mathcal{A}^n$과 다음 $\mathcal{O}$-쌍선형 사상들로 주어진다.
$$
\mathcal{A}^n \times \mathcal{A}^m \to \mathcal{A}^{n + m},\quad
(a, b) \longmapsto ab
$$
이 사상들을 곱셈 사상이라 하며 다음 성질을 요구한다.
\begin{enumerate}
\item 곱셈은 결합적이다.
\item $\mathcal{A}^0$에는 곱셈의 양쪽 항등원이 되는 대역 단면 $1$이 있다.
\end{enumerate}
이러한 구조를 흔히 $\mathcal{A}$로 나타낸다.
{\it 등급 $\mathcal{O}$-대수의 준동형}
$f : \mathcal{A} \to \mathcal{B}$란 곱셈 사상과 양립하는
$\mathcal{O}$-가군 사상들의 족
$f^n : \mathcal{A}^n \to \mathcal{B}^n$을 뜻한다.
\end{definition}

\noindent
등급 $\mathcal{O}$-대수 $\mathcal{A}$와 대상
$U \in \Ob(\mathcal{C})$가 주어졌을 때 다음 기호를 사용한다.
$$
\mathcal{A}(U) =
\Gamma(U, \mathcal{A}) =
\bigoplus\nolimits_{n \in \mathbf{Z}} \mathcal{A}^n(U)
$$
이는 등급 $\mathcal{O}(U)$-대수이다.

\begin{remark}
\label{sdga-remark-functoriality-ga}
$(f, f^\sharp) : (\Sh(\mathcal{C}), \mathcal{O}_\mathcal{C})
\to (\Sh(\mathcal{D}), \mathcal{O}_\mathcal{D})$를
환 달린 토포스의 사상이라 하자. 다음이 성립한다.
\begin{enumerate}
\item $\mathcal{A}$를 등급 $\mathcal{O}_\mathcal{C}$-대수라 하자.
$\mathcal{A}$의 곱셈 사상은 다음 곱셈 사상을 유도한다.
$f_*\mathcal{A}^n \times f_*\mathcal{A}^m \to f_*\mathcal{A}^{n + m}$
$f^\sharp$을 통하여 이 사상들을 $\mathcal{O}_\mathcal{D}$-쌍선형
사상으로 볼 수 있다. 이렇게 얻은 등급 $\mathcal{O}_\mathcal{D}$-대수를
$f_*\mathcal{A}$로 나타낸다.
\item $\mathcal{B}$를 등급 $\mathcal{O}_\mathcal{D}$-대수라 하자.
$\mathcal{B}$의 곱셈 사상은 다음 곱셈 사상을 유도한다.
$f^*\mathcal{B}^n \times f^*\mathcal{B}^m \to f^*\mathcal{B}^{n + m}$
$f^\sharp$을 사용하여 이 사상들을 $\mathcal{O}_\mathcal{C}$-쌍선형
사상으로 볼 수 있다. 이렇게 얻은 등급 $\mathcal{O}_\mathcal{C}$-대수를
$f^*\mathcal{B}$로 나타낸다.
\item 등급 $\mathcal{O}_\mathcal{C}$-대수의 준동형
$f^*\mathcal{B} \to \mathcal{A}$들의 집합은 등급
$\mathcal{O}_\mathcal{C}$-대수의 준동형
$\mathcal{B} \to f_*\mathcal{A}$들의 집합과 $1$ 대 $1$로 대응한다.
\end{enumerate}
(3)은 가군의 층에 대한 $f^*$와 $f_*$ 사이의 통상적인 수반에서
곧바로 따른다.
\end{remark}




\section{등급 가군의 층}
\label{sdga-section-graded-modules}

\noindent
이 절은 건너뛰어도 된다.

\begin{definition}
\label{sdga-definition-gm}
$(\mathcal{C}, \mathcal{O})$를 환 달린 사이트라 하고,
$\mathcal{A}$를 $(\mathcal{C}, \mathcal{O})$ 위의 등급 대수의 층이라 하자.
$\mathcal{A}$ 위의 (오른쪽) {\it 등급 $\mathcal{A}$-가군}, 또는
(오른쪽) {\it 등급 가군}은 $n \in \mathbf{Z}$를 첨자로 하는
$\mathcal{O}$-가군의 족 $\mathcal{M}^n$과 다음
$\mathcal{O}$-쌍선형 사상들로 주어진다.
$$
\mathcal{M}^n \times \mathcal{A}^m \to \mathcal{M}^{n + m},\quad
(x, a) \longmapsto xa
$$
이 사상들을 곱셈 사상이라 하며 다음 성질을 요구한다.
\begin{enumerate}
\item 곱셈은 $(xa)a' = x(aa')$를 만족한다.
\item $\mathcal{A}^0$의 항등 단면 $1$은 모든 $n$에 대하여
$\mathcal{M}^n$ 위에서 항등원으로 작용한다.
\end{enumerate}
이 상황을 흔히 ``$\mathcal{M}$을 등급 $\mathcal{A}$-가군이라 하자''라고
표현한다. {\it 등급 $\mathcal{A}$-가군의 준동형}
$f : \mathcal{M} \to \mathcal{N}$이란 곱셈 사상과 양립하는
$\mathcal{O}$-가군 사상들의 족
$f^n : \mathcal{M}^n \to \mathcal{N}^n$을 뜻한다.
(오른쪽) 등급 $\mathcal{A}$-가군의 범주를
$\textit{Mod}(\mathcal{A})$로 나타낸다.
\end{definition}

\noindent
{\it 왼쪽 등급 가군}도 완전히 같은 방식으로 정의할 수 있지만,
이 장에서는 오른쪽 가군을 기본으로 삼는다.

\medskip\noindent
등급 $\mathcal{A}$-가군 $\mathcal{M}$과 대상
$U \in \Ob(\mathcal{C})$가 주어졌을 때 다음 기호를 사용한다.
$$
\mathcal{M}(U) =
\Gamma(U, \mathcal{M}) =
\bigoplus\nolimits_{n \in \mathbf{Z}} \mathcal{M}^n(U)
$$
이는 (오른쪽) 등급 $\mathcal{A}(U)$-가군이다.

\begin{lemma}
\label{sdga-lemma-gm-abelian}
$(\mathcal{C}, \mathcal{O})$를 환 달린 사이트라 하고,
$\mathcal{A}$를 등급 $\mathcal{O}$-대수라 하자.
범주 $\textit{Mod}(\mathcal{A})$는 다음 성질을 갖는 아벨 범주이다.
\begin{enumerate}
\item $\textit{Mod}(\mathcal{A})$에는 임의의 직합이 존재한다.
\item $\textit{Mod}(\mathcal{A})$에는 임의의 여극한이 존재한다.
\item $\textit{Mod}(\mathcal{A})$에서 여과 여극한은 완전하다.
\item $\textit{Mod}(\mathcal{A})$에는 임의의 곱이 존재한다.
\item $\textit{Mod}(\mathcal{A})$에는 임의의 극한이 존재한다.
\end{enumerate}
등급 $\mathcal{A}$-가군을 그 $n$차 항으로 보내는 함자
$$
\textit{Mod}(\mathcal{A}) \longrightarrow \textit{Mod}(\mathcal{O}),\quad
\mathcal{M} \longmapsto \mathcal{M}^n
$$
는 모든 극한 및 여극한과 가환한다.
\end{lemma}

\noindent
이 보조정리는 극한과 여극한을 항별로 취할 수 있음을 말한다.
또한 다음 망각 함자가 모든 극한 및 여극한과 가환한다는 것도
말해 준다(원한다면 이를 따르는 결과로 보아도 된다).
$$
\textit{Mod}(\mathcal{A})
\longrightarrow
\text{등급 }\mathcal{O}\text{-가군}
$$

\begin{proof}
보조정리의 명제에 나오는 함자를
$\text{gr}^n : \textit{Mod}(\mathcal{A}) \to \textit{Mod}(\mathcal{O})$로
나타내자. 등급 $\mathcal{A}$-가군의 준동형
$f : \mathcal{M} \to \mathcal{N}$을 생각하자. 등급
$\mathcal{O}$-가군의 사상으로 본 $f$의 핵과 여핵에는
정의 \ref{sdga-definition-gm}과 같은 곱셈 사상이 추가로 주어진다.
따라서 이들은 $\textit{Mod}(\mathcal{A})$에서도 각각 핵과 여핵이다.
그러므로 $\textit{Mod}(\mathcal{A})$는 아벨 범주이고,
핵과 여핵을 취하는 연산은 $\text{gr}^n$과 가환한다.

\medskip\noindent
극한과 여극한의 존재를 보이려면 곱과 직합의 존재를 보이는 것으로
충분하다. 범주론, 보조정리
\ref{categories-lemma-limits-products-equalizers} 및
\ref{categories-lemma-colimits-coproducts-coequalizers}를 보라.
같은 보조정리들에 따르면, 극한 및 여극한과의 $\text{gr}^n$ 가환성을
보이려면 $\text{gr}^n$이 직합 및 곱과 가환함을 보이면 충분하다.

\medskip\noindent
$\mathcal{M}_t$, $t \in T$를 등급 $\mathcal{A}$-가군들의 집합이라 하자.
그 $n$차 항이 $\bigoplus_{t \in T} \mathcal{M}_t^n$이고 자명한 곱셈
사상을 갖는 등급 $\mathcal{A}$-가군을 생각할 수 있다. 이것이
$\textit{Mod}(\mathcal{A})$에서 직합임은 독자가 쉽게 확인할 수 있다.
곱도 마찬가지이다.

\medskip\noindent
모든 $n$에 대하여 $\text{gr}^n$은 완전 함자이며,
$\textit{Mod}(\mathcal{A})$의 복합체
$\mathcal{M}_1 \to \mathcal{M}_2 \to \mathcal{M}_3$가 완전일 필요충분조건은
$\text{gr}^n\mathcal{M}_1 \to \text{gr}^n\mathcal{M}_2 \to
\text{gr}^n\mathcal{M}_3$가 모든 $n$에 대하여
$\textit{Mod}(\mathcal{O})$에서 완전인 것이다.
따라서 (3)이 성립한다. 실제로 $\textit{Mod}(\mathcal{O})$는
그로텐디크 아벨 범주이므로 여과 여극한이 완전하다. 사이트 위의 코호몰로지,
절 \ref{sites-cohomology-section-unbounded}를 보라.
\end{proof}






\section{등급 가군의 층으로 이루어진 등급 범주}
\label{sdga-section-gm-gr-cat}

\noindent
이 절은 건너뛰어도 된다. 이 절은 미분 등급 대수, 예
\ref{dga-example-gm-gr-cat}에 대응한다. 등급 범주에 관한 규약은
미분 등급 대수, 절 \ref{dga-section-graded}를 보라.

\medskip\noindent
$(\mathcal{C}, \mathcal{O})$를 환 달린 사이트라 하고,
$\mathcal{A}$를 $(\mathcal{C}, \mathcal{O})$ 위의 등급 대수의 층이라 하자.
연관 범주 $(\textit{Mod}^{gr}(\mathcal{A}))^0$가 등급
$\mathcal{A}$-가군의 범주인, $R = \Gamma(\mathcal{C}, \mathcal{O})$ 위의
등급 범주 $\textit{Mod}^{gr}(\mathcal{A})$를 구성하겠다.
$\textit{Mod}^{gr}(\mathcal{A})$의 대상으로는 오른쪽 등급
$\mathcal{A}$-가군을 취한다(절 \ref{sdga-section-graded-modules}를 보라).
등급 $\mathcal{A}$-가군 $\mathcal{L}$과 $\mathcal{M}$이 주어졌을 때
다음과 같이 놓는다.
$$
\Hom_{\textit{Mod}^{gr}(\mathcal{A})}(\mathcal{L}, \mathcal{M}) =
\bigoplus\nolimits_{n \in \mathbf{Z}} \Hom^n(\mathcal{L}, \mathcal{M})
$$
여기서 $\Hom^n(\mathcal{L}, \mathcal{M})$은 차수 $n$인 동차 오른쪽
$\mathcal{A}$-가군 사상 $f : \mathcal{L} \to \mathcal{M}$들의 집합이다.
더 정확히 말하면, $f$는 곱셈 사상과 양립하는 사상들의 족
$f : \mathcal{L}^i \to \mathcal{M}^{i + n}$,
$i \in \mathbf{Z}$로 주어진다. 성분으로 쓰면
$$
\Hom^n(\mathcal{L}, \mathcal{M})
\subset
\prod\nolimits_{p + q = n}
\Hom_\mathcal{O}(\mathcal{L}^{-q}, \mathcal{M}^p)
$$
(첨자의 순서가 뒤집힘에 유의하라) 중 다음을 만족하는
$f = (f_{p, q})$들의 부분집합이다.
$$
f_{p, q}(m a) = f_{p - i, q + i}(m)a
$$
여기서 $a$는 $\mathcal{A}^i$의 국소 단면이고
$m$은 $\mathcal{L}^{-q - i}$의 국소 단면이다. 등급
$\mathcal{A}$-가군 $\mathcal{K}$, $\mathcal{L}$, $\mathcal{M}$에 대하여
$\textit{Mod}^{gr}(\mathcal{A})$의 합성은 다음 사상들로 정의한다.
$$
\Hom^m(\mathcal{L}, \mathcal{M}) \times
\Hom^n(\mathcal{K}, \mathcal{L}) \longrightarrow
\Hom^{n + m}(\mathcal{K}, \mathcal{M})
$$
이는 오른쪽 $\mathcal{A}$-가군 사상의 통상적인 합성,
$(g, f) \mapsto g \circ f$로 주어진다.





\section{등급 가군의 층에 대한 텐서곱}
\label{sdga-section-tensor-product}

\noindent
이 절은 건너뛰어도 된다. 이 절은 미분 등급 대수,
절 \ref{dga-section-tensor-product}의 일부에 대응한다.

\medskip\noindent
$(\mathcal{C}, \mathcal{O})$를 환 달린 사이트라 하고,
$\mathcal{A}$를 $(\mathcal{C}, \mathcal{O})$ 위의 등급 대수의 층이라 하자.
$\mathcal{M}$을 오른쪽 등급 $\mathcal{A}$-가군,
$\mathcal{N}$을 왼쪽 등급 $\mathcal{A}$-가군이라 하자.
그러면 {\it 텐서곱} $\mathcal{M} \otimes_\mathcal{A} \mathcal{N}$을
그 $n$차 항이 다음과 같은 등급 $\mathcal{O}$-가군으로 정의한다.
$$
(\mathcal{M} \otimes_\mathcal{A} \mathcal{N})^n =
\Coker\left(
\bigoplus\nolimits_{r + s + t = n} \mathcal{M}^r \otimes_\mathcal{O}
\mathcal{A}^s \otimes_\mathcal{O} \mathcal{N}^t
\longrightarrow
\bigoplus\nolimits_{p + q = n} \mathcal{M}^p \otimes_\mathcal{O} \mathcal{N}^q
\right)
$$
여기서 사상은
$\mathcal{M}^r \otimes_\mathcal{O} \mathcal{A}^s
\otimes_\mathcal{O} \mathcal{N}^t$의 국소 단면
$x \otimes a \otimes y$를 $xa \otimes y - x \otimes ay$로 보낸다.
이 정의에 따르면 $(\mathcal{M} \otimes_\mathcal{A} \mathcal{N})^n$은
준층 $U \mapsto (\mathcal{M}(U) \otimes_{\mathcal{A}(U)} \mathcal{N}(U))^n$의
층화이다. 여기서 등급 가군의 텐서곱은 미분 등급 대수,
절 \ref{dga-section-tensor-product}에서 정의한 것이다.

\medskip\noindent
왼쪽 등급 $\mathcal{A}$-가군 $\mathcal{N}$을 고정하면 다음 함자를 얻는다.
$$
- \otimes_\mathcal{A} \mathcal{N} :
\textit{Mod}(\mathcal{A})
\longrightarrow
\text{Gr}(\textit{Mod}(\mathcal{O})) =
\text{등급 }\mathcal{O}\text{-가군}
$$
$\text{Gr}(-)$라는 기호는 호몰로지, 정의
\ref{homology-definition-graded}를 보라. 등급 $\mathcal{O}$-가군의
등급 범주는 $\text{Gr}^{gr}(\textit{Mod}(\mathcal{O}))$로 나타낸다.
미분 등급 대수, 예 \ref{dga-example-graded-category-graded-objects}를 보라.
위 함자는 다음과 같은 등급 범주의 함자로 확장할 수 있다.
$$
- \otimes_\mathcal{A} \mathcal{N} :
\textit{Mod}^{gr}(\mathcal{A})
\longrightarrow
\text{Gr}^{gr}(\textit{Mod}(\mathcal{O}))
$$
이는 $\mathcal{M} \to \mathcal{M}'$인 차수 $n$의 준동형을
$\mathcal{M} \otimes_\mathcal{A} \mathcal{N}$에서
$\mathcal{M}' \otimes_\mathcal{A} \mathcal{N}$으로 가는 유도된
차수 $n$의 사상으로 보냄으로써 얻는다.






\section{등급 가군의 층에 대한 내부 Hom}
\label{sdga-section-internal-hom-graded}

\noindent
독자에게 이 절을 건너뛸 것을 권한다.

\medskip\noindent
절 \ref{sdga-section-gm-gr-cat}의 구성을 층화한 판본이 필요하다.
$(\mathcal{C}, \mathcal{O})$, $\mathcal{A}$,
$\mathcal{M}$, $\mathcal{L}$을 절 \ref{sdga-section-gm-gr-cat}에서와 같이 잡자.
다음을 정의한다.
$$
\SheafHom^{gr}_\mathcal{A}(\mathcal{M}, \mathcal{L})
$$
이는 그 $n$차 항이
$$
\SheafHom^n_\mathcal{A}(\mathcal{M}, \mathcal{L})
\subset
\prod\nolimits_{p + q = n}
\SheafHom_\mathcal{O}(\mathcal{L}^{-q}, \mathcal{M}^p)
$$
중 다음을 만족하는 국소 단면 $f = (f_{p, q})$들로 이루어진 부분층인
등급 $\mathcal{O}$-가군이다.
$$
f_{p, q}(m a) = f_{p - i, q + i}(m)a
$$
여기서 $a$는 $\mathcal{A}^i$의 국소 단면이고,
$m$은 $\mathcal{L}^{-q - i}$의 국소 단면이다. 절
\ref{sdga-section-gm-gr-cat}에서와 같이 다음 합성 사상이 있다.
$$
\SheafHom^{gr}_\mathcal{A}(\mathcal{L}, \mathcal{M}) \otimes_\mathcal{O}
\SheafHom^{gr}_\mathcal{A}(\mathcal{K}, \mathcal{L})
\longrightarrow
\SheafHom^{gr}_\mathcal{A}(\mathcal{K}, \mathcal{M})
$$
여기서 왼쪽 항은 절 \ref{sdga-section-tensor-product}에서 정의한 등급
$\mathcal{O}$-가군의 텐서곱이다. 이 사상은 다음 합성 사상으로 주어진다.
$$
\SheafHom^m_\mathcal{A}(\mathcal{L}, \mathcal{M}) \otimes_\mathcal{O}
\SheafHom^n_\mathcal{A}(\mathcal{K}, \mathcal{L}) \longrightarrow
\SheafHom^{n + m}_\mathcal{A}(\mathcal{K}, \mathcal{M})
$$
이 사상은 (국소적으로) 통상적인 합성으로 정의된다.

\medskip\noindent
이 정의들에 따라
$$
\Hom_{\textit{Mod}^{gr}(\mathcal{A})}(\mathcal{L}, \mathcal{M}) =
\Gamma(\mathcal{C}, \SheafHom^{gr}_\mathcal{A}(\mathcal{L}, \mathcal{M}))
$$
가 합성과 양립하는 등급 $R$-가군으로서 성립한다.







\section{등급 쌍가군의 층과 텐서-Hom 수반}
\label{sdga-section-graded-bimodules}

\noindent
이 절은 건너뛰어도 된다.

\begin{definition}
\label{sdga-definition-bimodule}
$(\mathcal{C}, \mathcal{O})$를 환 달린 사이트라 하고, $\mathcal{A}$와
$\mathcal{B}$를 $(\mathcal{C}, \mathcal{O})$ 위의 등급 대수의 층이라 하자.
% P03-STACKS-E1034
{\it 등급 $(\mathcal{A}, \mathcal{B})$-쌍가군}은
$n \in \mathbf{Z}$를 첨자로 하는 $\mathcal{O}$-가군의 족
$\mathcal{M}^n$과 다음 $\mathcal{O}$-쌍선형 사상들로 주어진다.
$$
\mathcal{M}^n \times \mathcal{B}^m \to \mathcal{M}^{n + m},\quad
(x, b) \longmapsto xb
$$
및
$$
\mathcal{A}^n \times \mathcal{M}^m \to \mathcal{M}^{n + m},\quad
(a, x) \longmapsto ax
$$
이 사상들을 곱셈 사상이라 하며 다음 성질을 요구한다.
\begin{enumerate}
\item 곱셈은 $a(a'x) = (aa')x$ 및 $(xb)b' = x(bb')$를 만족한다.
\item $(ax)b = a(xb)$,
\item $\mathcal{A}^0$의 항등 단면 $1$은 곱셈에 의해 항등원으로 작용한다.
\item $\mathcal{B}^0$의 항등 단면 $1$은 곱셈에 의해 항등원으로 작용한다.
\end{enumerate}
이러한 구조를 흔히 $\mathcal{M}$으로 나타낸다.
{\it 등급 $(\mathcal{A}, \mathcal{B})$-쌍가군의 준동형}
$f : \mathcal{M} \to \mathcal{N}$이란 곱셈 사상과 양립하는
$\mathcal{O}$-가군 사상들의 족
$f^n : \mathcal{M}^n \to \mathcal{N}^n$을 뜻한다.
\end{definition}

\noindent
등급 $(\mathcal{A}, \mathcal{B})$-쌍가군 $\mathcal{M}$과 대상
$U \in \Ob(\mathcal{C})$가 주어졌을 때 다음 기호를 사용한다.
$$
\mathcal{M}(U) =
\Gamma(U, \mathcal{M}) =
\bigoplus\nolimits_{n \in \mathbf{Z}} \mathcal{M}^n(U)
$$
이는 등급 $(\mathcal{A}(U), \mathcal{B}(U))$-쌍가군이다.

\medskip\noindent
$(\mathcal{C}, \mathcal{O})$를 환 달린 사이트라 하고, $\mathcal{A}$와
$\mathcal{B}$를 $(\mathcal{C}, \mathcal{O})$ 위의 등급 대수의 층이라 하자.
$\mathcal{M}$을 오른쪽 등급 $\mathcal{A}$-가군,
$\mathcal{N}$을 등급 $(\mathcal{A}, \mathcal{B})$-쌍가군이라 하자.
이 경우 절 \ref{sdga-section-tensor-product}에서 정의한 등급 텐서곱
$$
\mathcal{M} \otimes_\mathcal{A} \mathcal{N}
$$
에는 자명한 곱셈 사상이 주어져 오른쪽 등급 $\mathcal{B}$-가군이 된다.
이 구성은 다음 함자 및 등급 범주의 함자를 정의한다.
$$
\otimes_\mathcal{A} \mathcal{N} :
\textit{Mod}(\mathcal{A})
\longrightarrow
\textit{Mod}(\mathcal{B})
\quad\text{및}\quad
\otimes_\mathcal{A} \mathcal{N} :
\textit{Mod}^{gr}(\mathcal{A})
\longrightarrow
\textit{Mod}^{gr}(\mathcal{B})
$$
이는 $\mathcal{M} \to \mathcal{M}'$인 차수 $n$의 준동형을
$\mathcal{M} \otimes_\mathcal{A} \mathcal{N}$에서
$\mathcal{M}' \otimes_\mathcal{A} \mathcal{N}$으로 가는 유도된
차수 $n$의 사상으로 보냄으로써 얻는다.

\medskip\noindent
$(\mathcal{C}, \mathcal{O})$를 환 달린 사이트라 하고, $\mathcal{A}$와
$\mathcal{B}$를 $(\mathcal{C}, \mathcal{O})$ 위의 등급 대수의 층이라 하자.
$\mathcal{N}$을 등급 $(\mathcal{A}, \mathcal{B})$-쌍가군이라 하고,
$\mathcal{L}$을 오른쪽 등급 $\mathcal{B}$-가군이라 하자.
이 경우 절 \ref{sdga-section-internal-hom-graded}에서 정의한 등급 내부 Hom
$$
\SheafHom_\mathcal{B}^{gr}(\mathcal{N}, \mathcal{L})
$$
은 다음 곱셈 사상을 갖는 오른쪽 등급 $\mathcal{A}$-가군이다.\footnote{여기서
우리는 어떠한 부호도 들어가지 않는다는 규약을 사용한다.}
$$
\SheafHom^n_\mathcal{B}(\mathcal{N}, \mathcal{L})
\times \mathcal{A}^m
\longrightarrow
\SheafHom^{n + m}_\mathcal{B}(\mathcal{N}, \mathcal{L})
$$
이는 $U$ 위의 $\SheafHom^n_\mathcal{B}(\mathcal{N}, \mathcal{L})$의 단면
$f = (f_{p,q})$와 $U$ 위의 $\mathcal{A}^m$의 단면 $a$를,
$U$ 위의 $\SheafHom^{n + m}_\mathcal{B}(\mathcal{N}, \mathcal{L})$의
단면 $f a$로 보낸다. 이 단면은 다음 사상들의 족으로 정의된다.
% P03-STACKS-E1035
$$
\mathcal{N}^{-q - m}|_U \xrightarrow{a \cdot -}
\mathcal{N}^{-q}|_U \xrightarrow{f_{p, q}}
\mathcal{M}^p|_U
$$
% P03-STACKS-E1036
이것이 잘 정의된다는 확인은 생략한다. 이 구성은 다음 함자 및
등급 범주의 함자를 정의한다.
$$
\SheafHom_\mathcal{B}^{gr}(\mathcal{N}, -) :
\textit{Mod}(\mathcal{B})
\longrightarrow
\textit{Mod}(\mathcal{A})
\quad\text{및}\quad
\SheafHom_\mathcal{B}^{gr}(\mathcal{N}, -) :
\textit{Mod}^{gr}(\mathcal{B})
\longrightarrow
\textit{Mod}^{gr}(\mathcal{A})
$$
이는 $\mathcal{L} \to \mathcal{L}'$인 차수 $n$의 준동형을
$\SheafHom_\mathcal{B}^{gr}(\mathcal{N}, \mathcal{L})$에서
$\SheafHom_\mathcal{B}^{gr}(\mathcal{N}, \mathcal{L}')$으로 가는
유도된 차수 $n$의 사상으로 보냄으로써 얻는다.

\begin{lemma}
\label{sdga-lemma-tensor-hom-adjunction-gr}
$(\mathcal{C}, \mathcal{O})$를 환 달린 사이트라 하고, $\mathcal{A}$와
$\mathcal{B}$를 $(\mathcal{C}, \mathcal{O})$ 위의 등급 대수의 층이라 하자.
$\mathcal{M}$을 오른쪽 등급 $\mathcal{A}$-가군,
$\mathcal{N}$을 등급 $(\mathcal{A}, \mathcal{B})$-쌍가군,
$\mathcal{L}$을 오른쪽 등급 $\mathcal{B}$-가군이라 하자.
위의 규약 아래 다음이 성립한다.
$$
\Hom_{\textit{Mod}^{gr}(\mathcal{B})}(
\mathcal{M} \otimes_\mathcal{A} \mathcal{N}, \mathcal{L}) =
\Hom_{\textit{Mod}^{gr}(\mathcal{A})}(
\mathcal{M}, \SheafHom_\mathcal{B}^{gr}(\mathcal{N}, \mathcal{L}))
$$
및
$$
\SheafHom_\mathcal{B}^{gr}(
\mathcal{M} \otimes_\mathcal{A} \mathcal{N}, \mathcal{L}) =
\SheafHom_\mathcal{A}^{gr}(
\mathcal{M}, \SheafHom_\mathcal{B}^{gr}(\mathcal{N}, \mathcal{L}))
$$
$\mathcal{M}$, $\mathcal{N}$, $\mathcal{L}$에 대하여 함자적으로 성립한다.
\end{lemma}

\begin{proof}
생략한다. 힌트: 양변을 오른쪽에서 $\mathcal{B}$-선형인
$\mathcal{A}$-쌍선형 등급 사상
$\psi : \mathcal{M} \times \mathcal{N} \to \mathcal{L}$으로
해석하면 된다.
\end{proof}

\noindent
$(\mathcal{C}, \mathcal{O})$를 환 달린 사이트라 하고, $\mathcal{A}$와
$\mathcal{B}$를 $(\mathcal{C}, \mathcal{O})$ 위의 등급 대수의 층이라 하자.
위의 특수한 경우로서 등급 $\mathcal{O}$-대수의 준동형
$\varphi : \mathcal{A} \to \mathcal{B}$가 주어졌다고 하자.
그러면 다음 함자 및 등급 범주의 함자를 얻는다.
$$
\otimes_{\mathcal{A}, \varphi} \mathcal{B} :
\textit{Mod}(\mathcal{A})
\longrightarrow
\textit{Mod}(\mathcal{B})
\quad\text{및}\quad
\otimes_{\mathcal{A}, \varphi} \mathcal{B} :
\textit{Mod}^{gr}(\mathcal{A})
\longrightarrow
\textit{Mod}^{gr}(\mathcal{B})
$$
한편 다음 제한 함자들이 있다.
$$
res_\varphi :
\textit{Mod}(\mathcal{B})
\longrightarrow
\textit{Mod}(\mathcal{A})
\quad\text{및}\quad
res_\varphi :
\textit{Mod}^{gr}(\mathcal{B})
\longrightarrow
\textit{Mod}^{gr}(\mathcal{A})
$$
위의 보조정리를 사용하면 이 함자들이 서로 수반임을 보일 수 있다
(제한과 밑바꿈의 통상적인 경우와 같다). 구체적으로,
$\mathcal{B}$를 등급 $(\mathcal{A}, \mathcal{B})$-쌍가군으로 보았을 때
${}_\mathcal{A}\mathcal{B}_\mathcal{B}$라고 쓰자. 그러면 임의의
오른쪽 등급 $\mathcal{B}$-가군 $\mathcal{L}$에 대하여
$$
\SheafHom_\mathcal{B}^{gr}({}_\mathcal{A}\mathcal{B}_\mathcal{B}, \mathcal{L})
= res_\varphi(\mathcal{L})
$$
가 오른쪽 등급 $\mathcal{A}$-가군으로서 성립한다. 따라서 보조정리
\ref{sdga-lemma-tensor-hom-adjunction-gr}에 의해 다음 함자적 동형이 있다.
$$
\Hom_{\textit{Mod}^{gr}(\mathcal{B})}(
\mathcal{M} \otimes_{\mathcal{A}, \varphi} \mathcal{B}, \mathcal{L}) =
\Hom_{\textit{Mod}^{gr}(\mathcal{A})}(
\mathcal{M}, res_\varphi(\mathcal{L}))
$$
문맥에서 분명할 때에는 이 식에서 $\varphi$에 대한 의존성을 흔히
표기하지 않는다. 같은 방식으로 다음 등식을 얻는다.
$$
\SheafHom^{gr}_\mathcal{B}(
\mathcal{M} \otimes_\mathcal{A} \mathcal{B}, \mathcal{L}) =
\SheafHom_\mathcal{A}^{gr}(\mathcal{M}, \mathcal{L})
$$
이는 등급 $\mathcal{O}$-가군의 등식이다.




\section{등급 가군의 층에 대한 당김과 밂}
\label{sdga-section-functoriality-graded}

\noindent
독자에게 이 절을 건너뛸 것을 권한다.

\medskip\noindent
$(f, f^\sharp) : (\Sh(\mathcal{C}), \mathcal{O}_\mathcal{C})
\to (\Sh(\mathcal{D}), \mathcal{O}_\mathcal{D})$를 환 달린 토포스의
사상이라 하자. $\mathcal{A}$를 등급 $\mathcal{O}_\mathcal{C}$-대수,
$\mathcal{B}$를 등급 $\mathcal{O}_\mathcal{D}$-대수라 하자.
다음 사상이 주어졌다고 하자.
$$
\varphi : f^{-1}\mathcal{B} \to \mathcal{A}
$$
이는 등급 $f^{-1}\mathcal{O}_\mathcal{D}$-대수의 사상이다.
스칼라 제한과 스칼라 확장의 수반에 의해, 이는 등급
$\mathcal{O}_\mathcal{C}$-대수의 사상
$\varphi : f^*\mathcal{B} \to \mathcal{A}$와 같은 자료이다.
동치로, $\varphi$를 다음 사상으로 볼 수 있다.
$$
\varphi : \mathcal{B} \to f_*\mathcal{A}
$$
이는 등급 $\mathcal{O}_\mathcal{D}$-대수의 사상이다.
주의 \ref{sdga-remark-functoriality-ga}를 보라.

\medskip\noindent
다음 함자를 정의하자.
$$
f_* :
\textit{Mod}(\mathcal{A})
\longrightarrow
\textit{Mod}(\mathcal{B})
$$
등급 $\mathcal{A}$-가군 $\mathcal{M}$이 주어졌을 때,
$f_*\mathcal{M}$을 그 $n$차 항이 $f_*\mathcal{M}^n$인 등급
$\mathcal{B}$-가군으로 정의한다. 곱셈으로는 다음을 사용한다.
$$
f_*\mathcal{M}^n \times \mathcal{B}^m
\xrightarrow{(\text{id}, \varphi^m)}
f_*\mathcal{M}^n \times f_*\mathcal{A}^m
\xrightarrow{f_*\mu_{n, m}}
f_*\mathcal{M}^{n + m}
$$
여기서 $\mu_{n, m} : \mathcal{M}^n \times \mathcal{A}^m
\to \mathcal{M}^{n + m}$은 $\mathcal{A}$ 위의 $\mathcal{M}$에 대한
곱셈 사상이다. 여기서는 $f_*$가 곱과 가환한다는 사실을 사용한다.
이 구성은 분명히 $\mathcal{M}$에 대하여 함자적이므로 원하는 함자를 얻는다.

\medskip\noindent
다음 함자를 정의하자.
$$
f^* :
\textit{Mod}(\mathcal{B})
\longrightarrow
\textit{Mod}(\mathcal{A})
$$
이 함자를 다음 함자들의 합성으로 정의한다.
$$
\textit{Mod}(\mathcal{B})
\xrightarrow{f^{-1}}
\textit{Mod}(f^{-1}\mathcal{B})
\xrightarrow{ - \otimes_{f^{-1}\mathcal{B}} \mathcal{A}}
\textit{Mod}(\mathcal{A})
$$
먼저 등급 $\mathcal{B}$-가군 $\mathcal{N}$이 주어졌을 때,
$f^{-1}\mathcal{N}$을 그 $n$차 항이 $f^{-1}\mathcal{N}^n$인 등급
$f^{-1}\mathcal{B}$-가군으로 정의한다. 곱셈으로는 다음을 사용한다.
$$
f^{-1}\nu_{n, m} :
f^{-1}\mathcal{N}^n \times f^{-1}\mathcal{B}^m
\longrightarrow
f^{-1}\mathcal{N}^{n + m}
$$
여기서 $\nu_{n, m} : \mathcal{N}^n \times \mathcal{B}^m
\to \mathcal{N}^{n + m}$은 $\mathcal{B}$ 위의 $\mathcal{N}$에 대한
곱셈 사상이다. 여기서는 $f^{-1}$이 곱과 가환한다는 사실을 사용한다.
이 구성은 분명히 $\mathcal{N}$에 대하여 함자적이므로 함자 $f^{-1}$을 얻는다.
이제 절 \ref{sdga-section-graded-bimodules}의 텐서곱 논의를 사용하여
다음 함자를 정의할 수 있다.
$$
- \otimes_{f^{-1}\mathcal{B}} \mathcal{A} :
\textit{Mod}(f^{-1}\mathcal{B})
\longrightarrow
\textit{Mod}(\mathcal{A})
$$
마지막으로, 위에서 이미 예고한 대로 다음과 같이 놓는다.
$$
f^*\mathcal{N} =
f^{-1}\mathcal{N} \otimes_{f^{-1}\mathcal{B}, \varphi} \mathcal{A}
$$

\medskip\noindent
함자 $f_*$와 $f^*$는 다음 등급 범주의 함자로 쉽게 확장된다.
$$
f_* :
\textit{Mod}^{gr}(\mathcal{A})
\longrightarrow
\textit{Mod}^{gr}(\mathcal{B})
\quad\text{및}\quad
f^* :
\textit{Mod}^{gr}(\mathcal{B})
\longrightarrow
\textit{Mod}^{gr}(\mathcal{A})
$$
이들은 바탕 대상에 대해서는 같은 작용을 하며, 차수 $n$의 동차 사상에
대한 위 구성들의 함자성으로 정의된다.

\begin{lemma}
\label{sdga-lemma-adjunction-push-pull-gr}
위의 상황에서 다음이 성립한다.
$$
\Hom_{\textit{Mod}^{gr}(\mathcal{B})}(
\mathcal{N}, f_*\mathcal{M}) =
\Hom_{\textit{Mod}^{gr}(\mathcal{A})}(
f^*\mathcal{N}, \mathcal{M})
$$
\end{lemma}

\begin{proof}
생략한다. 힌트: 먼저 아벨 군의 층에 대한 $f^{-1}$과 $f_*$의 수반을
사용하여, $\textit{Mod}(\mathcal{B})$와
$\textit{Mod}(f^{-1}\mathcal{B})$ 사이의 함자로서 $f^{-1}$과 $f_*$가
수반임을 보인다. 다음으로 절 \ref{sdga-section-graded-bimodules}에서 주어진
밑바꿈과 제한 사이의 수반을 사용한다.
\end{proof}





\section{국소화와 등급 가군의 층}
\label{sdga-section-localize-graded}

\noindent
독자에게 이 절을 건너뛸 것을 권한다.

\medskip\noindent
$(\mathcal{C}, \mathcal{O})$를 환 달린 사이트라 하자.
$U \in \Ob(\mathcal{C})$라 하고 다음으로 나타내자.
$$
j :
(\Sh(\mathcal{C}/U), \mathcal{O}_U)
\longrightarrow
(\Sh(\mathcal{C}), \mathcal{O})
$$
이는 대응하는 국소화 사상이다(사이트 위의 가군,
절 \ref{sites-modules-section-localize}). 아래에서는 다음 사실을 사용한다.
$\mathcal{O}_U$-가군 $\mathcal{M}_i$, $i = 1, 2$와
$\mathcal{O}$-가군 $\mathcal{A}$에 대하여 다음 표준 사상이 있다.
$$
j_! :
\Hom_{\mathcal{O}_U}(
\mathcal{M}_1 \otimes_{\mathcal{O}_U} \mathcal{A}|_U, \mathcal{M}_2)
\longrightarrow
\Hom_\mathcal{O}(
j_!\mathcal{M}_1 \otimes_\mathcal{O} \mathcal{A}, j_!\mathcal{M}_2)
$$
실제로
$j_!(\mathcal{M}_1 \otimes_{\mathcal{O}_U} \mathcal{A}|_U) =
j_!\mathcal{M}_1 \otimes_\mathcal{O} \mathcal{A}$가 사이트 위의 가군,
보조정리 \ref{sites-modules-lemma-j-shriek-and-tensor}에 의해 성립한다.

\medskip\noindent
$\mathcal{A}$를 등급 $\mathcal{O}$-대수라 하자. $\mathcal{A}$의
$\mathcal{C}/U$로의 제한을 $\mathcal{A}_U$로 나타낸다. 즉,
$\mathcal{A}_U = j^*\mathcal{A} = j^{-1}\mathcal{A}$이다.
절 \ref{sdga-section-functoriality-graded}에서 다음 수반 함자들을 구성하였다.
$$
j_* :
\textit{Mod}^{gr}(\mathcal{A}_U)
\longrightarrow
\textit{Mod}^{gr}(\mathcal{A})
\quad\text{및}\quad
j^* :
\textit{Mod}^{gr}(\mathcal{A})
\longrightarrow
\textit{Mod}^{gr}(\mathcal{A}_U)
$$
여기서 $j^*$는 $j_*$의 왼쪽 수반이다. 이에 더하여 $j^*$의 왼쪽 수반인
다음 완전 함자가 있다고 주장한다.
$$
j_! :
\textit{Mod}^{gr}(\mathcal{A}_U)
\longrightarrow
\textit{Mod}^{gr}(\mathcal{A})
$$
구체적으로, 등급 $\mathcal{A}_U$-가군 $\mathcal{M}$이 주어졌을 때
$j_!\mathcal{M}$을 그 $n$차 항이 $j_!\mathcal{M}^n$인 등급
$\mathcal{A}$-가군으로 정의한다. 곱셈 사상으로는 다음을 사용한다.
$$
j_!\mu_{n, m} :
j_!\mathcal{M}^n \times \mathcal{A}^m \to
j_!\mathcal{M}^{n + m}
$$
여기서 $\mu_{m, n} : \mathcal{M}^n \times \mathcal{A}^m \to \mathcal{M}^{n + m}$은
주어진 곱셈 사상이다. 등급 $\mathcal{A}_U$-가군의 차수 $n$인 동차 사상
$f : \mathcal{M} \to \mathcal{M}'$이 주어지면 차수 $n$인 동차 사상
$j_!f : j_!\mathcal{M} \to j_!\mathcal{M}'$을 얻는다. 따라서 원하는
함자를 얻는다.
% P03-STACKS-E1037

\begin{lemma}
\label{sdga-lemma-extension-by-zero-graded}
위의 상황에서 다음이 성립한다.
$$
\Hom_{\textit{Mod}^{gr}(\mathcal{A})}(
j_!\mathcal{M}, \mathcal{N}) =
\Hom_{\textit{Mod}^{gr}(\mathcal{A}_U)}(
\mathcal{M}, j^*\mathcal{N})
$$
\end{lemma}

\begin{proof}
사이트 위의 가군, 절 \ref{sites-modules-section-localize}의 논의에 의해
$\mathcal{O}$-가군에 대한 함자 $j_!$와 $j^*$는 수반이다.
따라서 $\mathcal{O}$-가군 구조만 보면 다음을 안다.
$$
\Hom_{\text{Gr}^{gr}(\textit{Mod}(\mathcal{O}))}(
j_!\mathcal{M}, \mathcal{N}) =
\Hom_{\text{Gr}^{gr}(\textit{Mod}(\mathcal{O}_U))}(
\mathcal{M}, j^*\mathcal{N})
$$
($\text{Gr}^{gr}(\textit{Mod}(\mathcal{O}))$가 등급
$\mathcal{O}$-가군의 등급 범주를 나타냄을 상기하라.) 이제 이 동일시가
왼쪽 항의 $\mathcal{A}$-가군 사상을 오른쪽 항의 $\mathcal{A}_U$-가군
사상으로 보내는지 확인해야 한다. 이를 확인하기 위해,
$\mathcal{O}_U$-선형 사상
$f^n : \mathcal{M}^n \to j^*\mathcal{N}^{n + d}$
에 대응하는 $\mathcal{O}$-선형 사상
$g^n : j_!\mathcal{M}^n \to \mathcal{N}^{n + d}$
이 주어졌을 때 다음이 가환하는 것과
$$
\xymatrix{
\mathcal{M}^n \otimes_{\mathcal{O}_U} \mathcal{A}_U^m
\ar[r]_{f^n \otimes 1} \ar[d] &
j^*\mathcal{N}^{n + d} \otimes_{\mathcal{O}_U} \mathcal{A}_U^m \ar[d] \\
\mathcal{M}^{n + m} \ar[r]^{f^{n + m}} &
j^*\mathcal{N}^{n + m + d}
}
$$
다음이 가환하는 것이 동치임을 보이면 충분하다.
$$
\xymatrix{
j_!\mathcal{M}^n \otimes_\mathcal{O} \mathcal{A}^m
\ar[r]_{g^n \otimes 1} \ar[d] &
\mathcal{N}^{n + d} \otimes_\mathcal{O} \mathcal{A}_U^m \ar[d] \\
j_!\mathcal{M}^{n + m} \ar[r]^{g^{n + m}} &
\mathcal{N}^{n + m + d}
}
$$
% P03-STACKS-E1038
그런데 다음을 안다.
\begin{align*}
\Hom_{\mathcal{O}_U}(\mathcal{M}^n \otimes_{\mathcal{O}_U} \mathcal{A}_U^m,
j^*\mathcal{N}^{n + d + m})
& =
\Hom_\mathcal{O}(j_!(\mathcal{M}^n \otimes_{\mathcal{O}_U} \mathcal{A}_U^m),
\mathcal{N}^{n + d + m}) \\
& =
\Hom_\mathcal{O}(j_!\mathcal{M}^n \otimes_\mathcal{O} \mathcal{A}^m,
\mathcal{N}^{n + d + m})
\end{align*}
이는 앞서 사용한 사이트 위의 가군, 보조정리
\ref{sites-modules-lemma-j-shriek-and-tensor}에 따른다.
첫 번째 군에서 첫 번째 도표의 가환성에 대한 장애물이 마지막 군에서
두 번째 도표의 가환성에 대한 장애물로 사상된다는 확인은 생략한다.
\end{proof}

\begin{lemma}
\label{sdga-lemma-tensor-with-extension-by-zero}
위의 상황에서 $\mathcal{M}$을 오른쪽 등급 $\mathcal{A}_U$-가군,
$\mathcal{N}$을 왼쪽 등급 $\mathcal{A}$-가군이라 하자. 그러면
$$
j_!\mathcal{M} \otimes_\mathcal{A} \mathcal{N} =
j_!(\mathcal{M} \otimes_{\mathcal{A}_U} \mathcal{N}|_U)
$$
가 $\mathcal{M}$과 $\mathcal{N}$에 대하여 함자적으로, 등급
$\mathcal{O}$-가군으로서 성립한다.
\end{lemma}

\begin{proof}
$j_!\mathcal{M} \otimes_\mathcal{A} \mathcal{N}$의 차수 $n$ 성분은
다음 표준 사상의 여핵임을 상기하라.
$$
\bigoplus\nolimits_{r + s + t = n}
j_!\mathcal{M}^r \otimes_\mathcal{O}
\mathcal{A}^s \otimes_\mathcal{O}
\mathcal{N}^t
\longrightarrow
\bigoplus\nolimits_{p + q = n}
j_!\mathcal{M}^p \otimes_\mathcal{O} \mathcal{N}^q
$$
절 \ref{sdga-section-tensor-product}를 보라. 사이트 위의 가군, 보조정리
\ref{sites-modules-lemma-j-shriek-and-tensor}에 의해 이는 다음 사상의
여핵과 같다.
$$
\bigoplus\nolimits_{r + s + t = n}
j_!(\mathcal{M}^r \otimes_{\mathcal{O}_U}
\mathcal{A}^s|_U \otimes_{\mathcal{O}_U}
\mathcal{N}^t|_U)
\longrightarrow
\bigoplus\nolimits_{p + q = n}
j_!(\mathcal{M}^p \otimes_{\mathcal{O}_U} \mathcal{N}^q|_U)
$$
이로써 증명이 끝난다. 다른 증명으로는 위에서 인용한 보조정리의 증명에
나오는 요네다 논증을 되풀이할 수 있다.
\end{proof}








\section{등급 가군의 층에 대한 이동 함자}
\label{sdga-section-shift}

\noindent
독자에게 이 절을 건너뛸 것을 권한다. 등급 대수 위의 등급 가군의 층은
미분 등급 대수, 주의 \ref{dga-remark-graded-shift-functors}에서 논의한
현상의 한 예가 된다.

\medskip\noindent
$(\mathcal{C}, \mathcal{O})$를 환 달린 사이트라 하고,
$\mathcal{A}$를 $(\mathcal{C}, \mathcal{O})$ 위의 등급 대수의 층이라 하자.
$\mathcal{M}$을 등급 $\mathcal{A}$-가군이라 하고 $k \in \mathbf{Z}$라 하자.
$\mathcal{M}$의 {\it $k$차 이동} $\mathcal{M}[k]$를, 그 $n$차 부분이
다음과 같은 등급 $\mathcal{A}$-가군으로 정의한다.
$$
(\mathcal{M}[k])^n = \mathcal{M}^{n + k}
$$
이는 $\mathcal{M}$의 $(n + k)$차 부분이다. 곱셈 사상으로는
% P03-STACKS-E1039
$$
(\mathcal{M}[k])^n \times \mathcal{A}^m
\longrightarrow
(\mathcal{M}[k])^{n + m}
$$
단순히 $\mathcal{M}$의 다음 곱셈 사상을 사용한다.
$$
\mathcal{M}^{n + k} \times \mathcal{A}^m
\longrightarrow
\mathcal{M}^{n + m + k}
$$
함자 $[k]$를 정의하였고 $[k + l] = [k] \circ [l]$이며 다음이 성립함은
분명하다.
$$
\Hom_{\textit{Mod}^{gr}(\mathcal{A})}(\mathcal{L}, \mathcal{M}[k]) =
\Hom_{\textit{Mod}^{gr}(\mathcal{A})}(\mathcal{L}, \mathcal{M})[k]
$$
이는 (부호가 개입하지 않고) $\mathcal{M}$과 $\mathcal{L}$에 대하여
함자적으로 성립한다. 따라서 미분 등급 대수, 주의
\ref{dga-remark-graded-shift-functors}에서 논의한 바와 같이,
등급 $\mathcal{A}$-가군의 등급 범주는 등급 $\mathcal{A}$-가군의 통상적인
범주와 이동 함자들로부터 복원될 수 있음을 실제로 알 수 있다.

\begin{lemma}
\label{sdga-lemma-gm-grothendieck-abelian}
$(\mathcal{C}, \mathcal{O})$를 환 달린 사이트라 하고,
$\mathcal{A}$를 등급 $\mathcal{O}$-대수라 하자.
범주 $\textit{Mod}(\mathcal{A})$는 그로텐디크 아벨 범주이다.
\end{lemma}

\begin{proof}
보조정리 \ref{sdga-lemma-gm-abelian}와 그로텐디크 아벨 범주의 정의
(단사 대상, 정의 \ref{injectives-definition-grothendieck-conditions})에 의해
$\textit{Mod}(\mathcal{A})$가 생성자를 가짐을 보이면 충분하다. 다음이
생성자라고 주장한다.
$$
\mathcal{G} = \bigoplus\nolimits_{k, U} j_{U!}\mathcal{A}_U[k]
$$
여기서 합은 $\mathcal{C}$의 모든 대상 $U$와 $k \in \mathbf{Z}$에 걸친다.
실제로 등급 $\mathcal{A}$-가군 $\mathcal{M}$이 주어졌을 때
$\mathcal{G}$에서 $\mathcal{M}$으로 가는 영이 아닌 사상이 없다면,
모든 $k$와 $U$에 대하여
$$
\Hom_{\textit{Mod}(\mathcal{A})}(j_{U!}\mathcal{A}_U[k], \mathcal{M}) =
\Hom_{\textit{Mod}(\mathcal{A}_U)}(\mathcal{A}_U[k], \mathcal{M}|_U) =
\Gamma(U, \mathcal{M}^{-k})
$$
가 영이다. 따라서 $\mathcal{M}$은 영이다.
\end{proof}









\section{미분 등급 대수의 층}
\label{sdga-section-dga}

\noindent
이 절은 미분 등급 대수, 절 \ref{dga-section-dga}에 대응한다.

\begin{definition}
\label{sdga-definition-dga}
$(\mathcal{C}, \mathcal{O})$를 환 달린 사이트라 하자.
$(\mathcal{C}, \mathcal{O})$ 위의
{\it 미분 등급 $\mathcal{O}$-대수의 층}, 또는
{\it 미분 등급 대수의 층}은 다음 $\mathcal{O}$-쌍선형 사상들이
주어진 $\mathcal{O}$-가군의 공사슬 복합체 $\mathcal{A}^\bullet$이다.
$$
\mathcal{A}^n \times \mathcal{A}^m \to \mathcal{A}^{n + m},\quad
(a, b) \longmapsto ab
$$
이 사상들을 곱셈 사상이라 하며 다음 성질을 요구한다.
\begin{enumerate}
\item 곱셈은 결합적이다.
\item $\mathcal{A}^0$에는 곱셈의 양쪽 항등원이 되는 대역 단면 $1$이 있다.
\item $U \in \Ob(\mathcal{C})$, $a \in \mathcal{A}^n(U)$,
$b \in \mathcal{A}^m(U)$에 대하여 다음이 성립한다.
$$
\text{d}^{n + m}(ab) = \text{d}^n(a)b + (-1)^n a\text{d}^m(b)
$$
\end{enumerate}
이러한 구조를 흔히 $(\mathcal{A}, \text{d})$로 나타낸다.
$(\mathcal{A}, \text{d})$에서 $(\mathcal{B}, \text{d})$로 가는
{\it 미분 등급 $\mathcal{O}$-대수의 준동형}이란 곱셈 사상과 양립하는
$\mathcal{O}$-가군 복합체의 사상
$f : \mathcal{A}^\bullet \to \mathcal{B}^\bullet$을 뜻한다.
\end{definition}

\noindent
미분 등급 $\mathcal{O}$-대수 $(\mathcal{A}, \text{d})$와 대상
$U \in \Ob(\mathcal{C})$가 주어졌을 때 다음 기호를 사용한다.
$$
\mathcal{A}(U) =
\Gamma(U, \mathcal{A}) =
\bigoplus\nolimits_{n \in \mathbf{Z}} \mathcal{A}^n(U)
$$
이는 미분 등급 $\mathcal{O}(U)$-대수이다.

\medskip\noindent
가능한 한 미분 등급 $\mathcal{O}$-대수 $(\mathcal{A}, \text{d})$를,
정의에 주어진 라이프니츠 법칙을 만족하는 차수 $1$의 연산자
$\text{d} : \mathcal{A} \to \mathcal{A}$가 주어진 등급
$\mathcal{O}$-대수 $\mathcal{A}$로 생각하겠다
($\mathcal{A}$는 등급 $\mathcal{O}$-가군으로 본다).

\begin{remark}
\label{sdga-remark-functoriality-dga}
$(f, f^\sharp) : (\Sh(\mathcal{C}), \mathcal{O}_\mathcal{C})
\to (\Sh(\mathcal{D}), \mathcal{O}_\mathcal{D})$를 환 달린 토포스의
사상이라 하자.
\begin{enumerate}
\item $(\mathcal{A}, \text{d})$를 미분 등급
$\mathcal{O}_\mathcal{C}$-대수라 하자. 밂은 미분 등급
$\mathcal{O}_\mathcal{D}$-대수 $(f_*\mathcal{A}, \text{d})$이다.
여기서 $f_*\mathcal{A}$는 주의 \ref{sdga-remark-functoriality-ga}에서와 같고,
$f_*\mathcal{A}^n \to f_*\mathcal{A}^{n + 1}$인 사상으로서
$\text{d} = f_*\text{d}$이다. 라이프니츠 법칙이 성립한다는 확인은 생략한다.
\item $\mathcal{B}$를 미분 등급 $\mathcal{O}_\mathcal{D}$-대수라 하자.
당김은 미분 등급 $\mathcal{O}_\mathcal{C}$-대수
$(f^*\mathcal{B}, \text{d})$이다. 여기서 $f^*\mathcal{B}$는 주의
\ref{sdga-remark-functoriality-ga}에서와 같고,
$f^*\mathcal{B}^n \to f^*\mathcal{B}^{n + 1}$인 사상으로서
$\text{d} = f^*\text{d}$이다. 라이프니츠 법칙이 성립한다는 확인은 생략한다.
\item 미분 등급 $\mathcal{O}_\mathcal{C}$-대수의 준동형
$f^*\mathcal{B} \to \mathcal{A}$들의 집합은 미분 등급
$\mathcal{O}_\mathcal{D}$-대수의 준동형
$\mathcal{B} \to f_*\mathcal{A}$들의 집합과 $1$ 대 $1$로 대응한다.
\end{enumerate}
(3)은 가군의 층에 대한 $f^*$와 $f_*$ 사이의 통상적인 수반에서
곧바로 따른다.
\end{remark}
















\section{미분 등급 가군의 층}
\label{sdga-section-modules}

\noindent
이 절은 미분 등급 대수, 절 \ref{dga-section-modules}에 대응한다.

\begin{definition}
\label{sdga-definition-dgm}
$(\mathcal{C}, \mathcal{O})$를 환 달린 사이트라 하고,
$(\mathcal{A}, \text{d})$를 $(\mathcal{C}, \mathcal{O})$ 위의
미분 등급 대수의 층이라 하자. $\mathcal{A}$ 위의 (오른쪽)
{\it 미분 등급 $\mathcal{A}$-가군}, 또는 (오른쪽)
{\it 미분 등급 가군}은 다음 $\mathcal{O}$-쌍선형 사상들이 주어진
공사슬 복합체 $\mathcal{M}^\bullet$이다.
$$
\mathcal{M}^n \times \mathcal{A}^m \to \mathcal{M}^{n + m},\quad
(x, a) \longmapsto xa
$$
이 사상들을 곱셈 사상이라 하며 다음 성질을 요구한다.
\begin{enumerate}
\item 곱셈은 $(xa)a' = x(aa')$를 만족한다.
\item $\mathcal{A}^0$의 항등 단면 $1$은 모든 $n$에 대하여
$\mathcal{M}^n$ 위에서 항등원으로 작용한다.
\item $U \in \Ob(\mathcal{C})$, $x \in \mathcal{M}^n(U)$,
$a \in \mathcal{A}^m(U)$에 대하여 다음이 성립한다.
$$
\text{d}^{n + m}(xa) = \text{d}^n(x)a + (-1)^n x\text{d}^m(a)
$$
\end{enumerate}
이 상황을 흔히 ``$\mathcal{M}$을 미분 등급 $\mathcal{A}$-가군이라 하자''라고
표현한다. $\mathcal{M}$에서 $\mathcal{N}$으로 가는
{\it 미분 등급 $\mathcal{A}$-가군의 준동형}이란 곱셈 사상과 양립하는
$\mathcal{O}$-가군 복합체의 사상
$f : \mathcal{M}^\bullet \to \mathcal{N}^\bullet$을 뜻한다.
(오른쪽) 미분 등급 $\mathcal{A}$-가군의 범주를
$\textit{Mod}(\mathcal{A}, \text{d})$로 나타낸다.
\end{definition}

\noindent
{\it 왼쪽 미분 등급 가군}도 완전히 같은 방식으로 정의할 수 있지만,
이 장에서는 오른쪽 가군을 기본으로 삼는다.

\medskip\noindent
미분 등급 $\mathcal{A}$-가군 $\mathcal{M}$과 대상
$U \in \Ob(\mathcal{C})$가 주어졌을 때 다음 기호를 사용한다.
$$
\mathcal{M}(U) =
\Gamma(U, \mathcal{M}) =
\bigoplus\nolimits_{n \in \mathbf{Z}} \mathcal{M}^n(U)
$$
이는 (오른쪽) 미분 등급 $\mathcal{A}(U)$-가군이다.

\begin{lemma}
\label{sdga-lemma-dgm-abelian}
$(\mathcal{C}, \mathcal{O})$를 환 달린 사이트라 하고,
$(\mathcal{A}, \text{d})$를 미분 등급 $\mathcal{O}$-대수라 하자.
범주 $\textit{Mod}(\mathcal{A}, \text{d})$는 다음 성질을 갖는 아벨 범주이다.
\begin{enumerate}
\item $\textit{Mod}(\mathcal{A}, \text{d})$에는 임의의 직합이 존재한다.
\item $\textit{Mod}(\mathcal{A}, \text{d})$에는 임의의 여극한이 존재한다.
\item $\textit{Mod}(\mathcal{A}, \text{d})$에서 여과 여극한은 완전하다.
% P03-STACKS-E1040
\item $\textit{Mod}(\mathcal{A}, \text{d})$에는 임의의 곱이 존재한다.
\item $\textit{Mod}(\mathcal{A}, \text{d})$에는 임의의 극한이 존재한다.
\end{enumerate}
다음 망각 함자는
$$
\textit{Mod}(\mathcal{A}, \text{d})
\longrightarrow
\textit{Mod}(\mathcal{A})
$$
미분 등급 $\mathcal{A}$-가군을 그 바탕 등급 가군으로 보내며,
모든 극한 및 여극한과 가환한다.
\end{lemma}

\begin{proof}
보조정리의 명제에 나오는 함자를
$F : \textit{Mod}(\mathcal{A}, \text{d}) \to \textit{Mod}(\mathcal{A})$로
나타내자. 범주 $\textit{Mod}(\mathcal{A})$가 성질 (1) -- (5)를 가짐에
유의하라. 보조정리 \ref{sdga-lemma-gm-abelian}을 보라.

\medskip\noindent
등급 $\mathcal{A}$-가군의 준동형
$f : \mathcal{M} \to \mathcal{N}$을 생각하자.
% P03-STACKS-E1041
등급 $\mathcal{A}$-가군의 사상으로 본 $f$의 핵과 여핵에는
정의 \ref{sdga-definition-dgm}과 같은 미분이 추가로 주어진다. 따라서 이들은
$\textit{Mod}(\mathcal{A}, \text{d})$에서도 각각 핵과 여핵이다.
그러므로 $\textit{Mod}(\mathcal{A}, \text{d})$는 아벨 범주이고,
핵과 여핵을 취하는 연산은 $F$와 가환한다.

\medskip\noindent
극한과 여극한의 존재를 보이려면 곱과 직합의 존재를 보이는 것으로
충분하다. 범주론, 보조정리
\ref{categories-lemma-limits-products-equalizers} 및
\ref{categories-lemma-colimits-coproducts-coequalizers}를 보라.
같은 보조정리들에 따르면, 극한 및 여극한과의 $F$ 가환성을 보이려면
$F$가 직합 및 곱과 가환함을 보이면 충분하다.

\medskip\noindent
$\mathcal{M}_t$, $t \in T$를 미분 등급 $\mathcal{A}$-가군들의 집합이라 하자.
직합 $\bigoplus \mathcal{M}_t$를 등급 $\mathcal{A}$-가군으로 생각할 수 있다.
등급 가군의 직합은 항별로 취하므로 $\bigoplus \mathcal{M}_t$에 미분이
주어짐은 분명하다. 이것이 $\textit{Mod}(\mathcal{A}, \text{d})$에서
직합임은 독자가 쉽게 확인할 수 있다. 곱도 마찬가지이다.

\medskip\noindent
$F$가 완전 함자이고 $\textit{Mod}(\mathcal{A}, \text{d})$의 복합체
$\mathcal{M}_1 \to \mathcal{M}_2 \to \mathcal{M}_3$가 완전일 필요충분조건은
$F(\mathcal{M}_1) \to F(\mathcal{M}_2) \to F(\mathcal{M}_3)$가
$\textit{Mod}(\mathcal{A})$에서 완전인 것임에 유의하라. 따라서
$\textit{Mod}(\mathcal{A})$에서 여과 여극한이 완전하므로 (3)이 성립한다.
\end{proof}

\noindent
보조정리 \ref{sdga-lemma-dgm-abelian}과 \ref{sdga-lemma-gm-abelian}을 결합하면
다음과 같은 완전하고 충실한 함자가 있음을 알 수 있다.
$$
\textit{Mod}(\mathcal{A}, \text{d}) \longrightarrow \text{Comp}(\mathcal{O})
$$
이는 아벨 범주 사이의 함자이다. 미분 등급 $\mathcal{A}$-가군
$\mathcal{M}$에 대하여 $H^i(\mathcal{M})$로 나타내는 코호몰로지
$\mathcal{O}$-가군은 $\mathcal{M}$에 대응하는 $\mathcal{O}$-가군
복합체의 코호몰로지로 정의한다. 따라서 미분 등급 $\mathcal{A}$-가군의
짧은 완전열 $0 \to \mathcal{K} \to \mathcal{L} \to \mathcal{M} \to 0$은
다음 긴 완전열을 유도한다.
\begin{equation}
\label{sdga-equation-les}
H^n(\mathcal{K}) \to H^n(\mathcal{L}) \to H^n(\mathcal{M}) \to
H^{n + 1}(\mathcal{K})
\end{equation}
이는 코호몰로지 가군의 열이다. 호몰로지, 보조정리
\ref{homology-lemma-long-exact-sequence-cochain}을 보라.

\medskip\noindent
더 나아가 이제부터 가군 복합체에 사용하는 모든 용어를 빌려 쓴다.
예를 들어, 모든 $k \in \mathbf{Z}$에 대하여
$H^k(\mathcal{M}) = 0$이면 미분 등급 $\mathcal{A}$-가군
$\mathcal{M}$을 {\it 비순환}이라고 한다. 미분 등급
$\mathcal{A}$-가군의 준동형 $\mathcal{M} \to \mathcal{N}$이 모든
$k \in \mathbf{Z}$에 대하여 동형
$H^k(\mathcal{M}) \to H^k(\mathcal{N})$을 유도하면 이를
{\it 유사동형}이라고 한다. 이와 같은 방식으로 용어를 사용한다.


















\section{가군의 미분 등급 범주}
\label{sdga-section-dgm-dg-cat}

\noindent
이 절은 미분 등급 대수, 예 \ref{dga-example-dgm-dg-cat}에 대응한다.
미분 등급 범주에 관한 규약은 미분 등급 대수,
절 \ref{dga-section-dga-categories}를 보라.

\medskip\noindent
$(\mathcal{C}, \mathcal{O})$를 환 달린 사이트라 하고,
$(\mathcal{A}, \text{d})$를 $(\mathcal{C}, \mathcal{O})$ 위의
미분 등급 대수의 층이라 하자. 다음 미분 등급 범주를 구성하겠다.
$$
\textit{Mod}^{dg}(\mathcal{A}, \text{d})
$$
이는 $R = \Gamma(\mathcal{C}, \mathcal{O})$ 위의 범주이고, 그에 연관된
복합체 범주는 미분 등급 $\mathcal{A}$-가군의 범주이다.
$$
\textit{Mod}(\mathcal{A}, \text{d}) =
\text{Comp}(\textit{Mod}^{dg}(\mathcal{A}, \text{d}))
$$
$\textit{Mod}^{dg}(\mathcal{A}, \text{d})$의 대상으로 오른쪽 미분 등급
$\mathcal{A}$-가군을 취한다. 절 \ref{sdga-section-modules}를 보라.
미분 등급 $\mathcal{A}$-가군 $\mathcal{L}$과 $\mathcal{M}$이 주어졌을 때
다음과 같이 놓는다.
$$
\Hom_{\textit{Mod}^{dg}(\mathcal{A}, \text{d})}(\mathcal{L}, \mathcal{M}) =
\Hom_{\textit{Mod}^{gr}(\mathcal{A})}(\mathcal{L}, \mathcal{M}) =
\bigoplus\nolimits_{n \in \mathbf{Z}} \Hom^n(\mathcal{L}, \mathcal{M})
$$
이는 등급 $R$-가군의 등식이다. 절 \ref{sdga-section-gm-gr-cat}를 보라.
달리 말하면, $n$차 등급 성분 $\Hom^n(\mathcal{L}, \mathcal{M})$은
차수 $n$인 동차 오른쪽 $\mathcal{A}$-가군 사상들의 $R$-가군이다.
원소 $f \in \Hom^n(\mathcal{L}, \mathcal{M})$에 대하여 다음과 같이 놓는다.
$$
\text{d}(f) =
\text{d}_\mathcal{M} \circ f - (-1)^n f \circ \text{d}_\mathcal{L}
$$
이를 해석하기 위해 $\text{d}_\mathcal{M}$과 $\text{d}_\mathcal{L}$을
등급 $\mathcal{O}$-가군 사상으로 보고, 등급 $\mathcal{O}$-가군 사상의
합성을 사용한다. $\text{d}(f)$가 등급 $\mathcal{O}$-가군 사상으로서
차수 $n + 1$인 동차 사상임은 분명하다. 또한 $\mathcal{M}$과
$\mathcal{A}$의 동차 국소 단면 $x$와 $a$에 대하여 다음이 성립하므로
$\mathcal{A}$-선형이다.
% P03-STACKS-E1042
\begin{align*}
\text{d}(f)(xa)
& =
\text{d}_\mathcal{M}(f(x) a) - (-1)^n f (\text{d}_\mathcal{L}(xa)) \\
& =
\text{d}_\mathcal{M}(f(x)) a + (-1)^{\deg(x) + n} f(x) \text{d}(a) 
- (-1)^n f(\text{d}_\mathcal{L}(x)) a - (-1)^{n + \deg(x)} f(x) \text{d}(a) \\
& = \text{d}(f)(x) a
\end{align*}
이는 원하는 결과이다($\Hom$ 위의 미분 정의에 있는 부호가 없으면
이 계산이 성립하지 않음에 유의하라).

\medskip\noindent
미분 등급 $\mathcal{A}$-가군
$\mathcal{K}$, $\mathcal{L}$, $\mathcal{M}$에 대해서는 이미 다음 합성을
정의하였다.
$$
\Hom^m(\mathcal{L}, \mathcal{M}) \times
\Hom^n(\mathcal{K}, \mathcal{L}) \longrightarrow
\Hom^{n + m}(\mathcal{K}, \mathcal{M})
$$
이는 절 \ref{sdga-section-gm-gr-cat}에서 층 사상의 통상적인 합성으로 정의하였다.
이 합성은 미분 등급 가군의 다음 사상을 정의한다.
$$
\Hom_{\textit{Mod}^{dg}(\mathcal{A}, \text{d})}(\mathcal{L}, \mathcal{M})
\otimes_R
\Hom_{\textit{Mod}^{dg}(\mathcal{A}, \text{d})}(\mathcal{K}, \mathcal{L})
\longrightarrow
\Hom_{\textit{Mod}^{dg}(\mathcal{A}, \text{d})}(\mathcal{K}, \mathcal{M})
$$
이는 미분 등급 대수, 정의 \ref{dga-definition-dga-category}에서
요구되는 사상이다. 실제로
\begin{align*}
\text{d}(g \circ f) & =
\text{d}_\mathcal{M} \circ g \circ f
- (-1)^{n + m} g \circ f \circ \text{d}_\mathcal{K} \\
& =
\left(\text{d}_\mathcal{M} \circ g - (-1)^m g \circ \text{d}_L\right) \circ f
+ (-1)^m g \circ \left(\text{d}_\mathcal{L} \circ f
- (-1)^n f \circ \text{d}_\mathcal{K}\right) \\
& =
\text{d}(g) \circ f + (-1)^m g \circ \text{d}(f)
\end{align*}
$f$의 차수가 $n$이고 $g$의 차수가 $m$이면 원하는 대로 성립한다.
% P03-STACKS-E1043






\section{미분 등급 가군의 층에 대한 텐서곱}
\label{sdga-section-tensor-product-dg}

\noindent
이 절은 미분 등급 대수, 절 \ref{dga-section-tensor-product}의 일부에
대응한다.

\medskip\noindent
$(\mathcal{C}, \mathcal{O})$를 환 달린 사이트라 하고,
$(\mathcal{A}, \text{d})$를 $(\mathcal{C}, \mathcal{O})$ 위의 미분 등급
대수의 층이라 하자. $\mathcal{M}$을 오른쪽 미분 등급
$\mathcal{A}$-가군, $\mathcal{N}$을 왼쪽 미분 등급
$\mathcal{A}$-가군이라 하자. 이 상황에서 {\it 텐서곱}
$\mathcal{M} \otimes_\mathcal{A} \mathcal{N}$을 다음과 같이 정의한다.
등급 $\mathcal{O}$-가군으로서는 절 \ref{sdga-section-tensor-product}의 구성으로
주어지며, 다음 미분이 주어진다.
$$
\text{d}_{\mathcal{M} \otimes_\mathcal{A} \mathcal{N}} :
(\mathcal{M} \otimes_\mathcal{A} \mathcal{N})^n
\longrightarrow
(\mathcal{M} \otimes_\mathcal{A} \mathcal{N})^{n + 1}
$$
이 미분은 다음 규칙으로 정의된다.
$$
\text{d}_{\mathcal{M} \otimes_\mathcal{A} \mathcal{N}}(x \otimes y) =
\text{d}_\mathcal{M}(x) \otimes y +
(-1)^{\deg(x)}x \otimes \text{d}_\mathcal{N}(y)
$$
여기서 $x$와 $y$는 $\mathcal{M}$과 $\mathcal{N}$의 동차 국소 단면이다.
이것이 잘 정의됨을 보이려면 $\mathcal{M}$, $\mathcal{A}$,
$\mathcal{N}$의 동차 국소 단면 $x$, $a$, $y$에 대하여
$\text{d}_{\mathcal{M} \otimes_\mathcal{A} \mathcal{N}}$가
$xa \otimes y - x \otimes ay$ 꼴의 원소를 영으로 보냄을 보여야 한다.
계산하면
\begin{align*}
&
\text{d}_{\mathcal{M} \otimes_\mathcal{A} \mathcal{N}}(
xa \otimes y - x \otimes ay) \\
& =
\text{d}_\mathcal{M}(xa) \otimes y + (-1)^{\deg(x) + \deg(a)}
xa \otimes \text{d}_\mathcal{N}(y)
-\text{d}_\mathcal{M}(x) \otimes ay - (-1)^{\deg(x)}
x \otimes \text{d}_\mathcal{N}(ay) \\
& =
\text{d}_\mathcal{M}(x)a \otimes y + (-1)^{\deg(x)}x\text{d}(a) \otimes y
+ (-1)^{\deg(x) + \deg(a)} xa \otimes \text{d}_\mathcal{N}(y) \\
&
-\text{d}_\mathcal{M}(x) \otimes ay - (-1)^{\deg(x)}
x \otimes \text{d}(a)y - (-1)^{\deg(x) +
\deg(a)} x\otimes a\text{d}_\mathcal{N}(y)
\end{align*}
이다. 이제 다음 원소들이
$$
\text{d}_\mathcal{M}(x)a \otimes y - \text{d}_\mathcal{M}(x) \otimes ay,\quad
x\text{d}(a) \otimes y - x \otimes \text{d}(a)y,\quad
\text{및}\quad
xa \otimes \text{d}_\mathcal{N}(y) - x\otimes a\text{d}_\mathcal{N}(y)
$$
모두 $\mathcal{M} \otimes_\mathcal{A} \mathcal{N}$에서 영으로 가므로
결론을 얻는다. 다음 등식의 확인은 생략한다.
$\text{d}_{\mathcal{M} \otimes_\mathcal{A} \mathcal{N}} \circ
\text{d}_{\mathcal{M} \otimes_\mathcal{A} \mathcal{N}} = 0$.

\medskip\noindent
왼쪽 미분 등급 $\mathcal{A}$-가군 $\mathcal{N}$을 고정하면 다음 함자를
얻는다.
$$
- \otimes_\mathcal{A} \mathcal{N} :
\textit{Mod}(\mathcal{A}, \text{d})
\longrightarrow
\text{Comp}(\mathcal{O})
$$
오른쪽 항은 $\mathcal{O}$-가군 복합체의 범주이다. 이 함자는 다음
미분 등급 범주의 함자로 확장할 수 있다.
$$
- \otimes_\mathcal{A} \mathcal{N} :
\textit{Mod}^{dg}(\mathcal{A}, \text{d})
\longrightarrow
\text{Comp}^{dg}(\mathcal{O})
$$
바탕 등급 대상에서는 차수 $n$의 준동형
$f : \mathcal{M} \to \mathcal{M}'$을 차수 $n$의 사상
$f \otimes \text{id}_\mathcal{N} :
\mathcal{M} \otimes_\mathcal{A} \mathcal{N} \to
\mathcal{M}' \otimes_\mathcal{A} \mathcal{N}$으로 보낸다. 이는
절 \ref{sdga-section-tensor-product}에서 한 구성이다. 이 구성이 성립함을
보이려면 다음 사상이 미분과 양립함을 확인해야 한다.
$$
\Hom_{\textit{Mod}^{dg}(\mathcal{A}, \text{d})}(\mathcal{M}, \mathcal{M}')
\longrightarrow
\Hom_{\text{Comp}^{dg}(\mathcal{O})}(
\mathcal{M} \otimes_\mathcal{A} \mathcal{N},
\mathcal{M}' \otimes_\mathcal{A} \mathcal{N})
$$
위와 같은 $f$에 대하여 이를 보이려면
$$
(\text{d}_{\mathcal{M}'} \circ f - (-1)^n f \circ \text{d}_\mathcal{M})
\otimes \text{id}_\mathcal{N}
$$
가 다음과 같음을 보여야 한다.
$$
\text{d}_{\mathcal{M}' \otimes_\mathcal{A} \mathcal{N}}
\circ (f \otimes \text{id}_\mathcal{N})
- (-1)^n (f \otimes \text{id}_\mathcal{N}) \circ
\text{d}_{\mathcal{M} \otimes_\mathcal{A} \mathcal{N}}
$$
$x$와 $y$가 $\mathcal{M}$과 $\mathcal{N}$의 동차 국소 단면일 때,
$x \otimes y$ 꼴의 국소 단면에 대한 이 연산자들의 작용을 계산하자.
첫 번째 연산자에서는
$$
(\text{d}_{\mathcal{M}'}(f(x)) - (-1)^n f(\text{d}_\mathcal{M}(x))) \otimes y
$$
를 얻고, 두 번째 연산자에서는
\begin{align*}
&\text{d}_{\mathcal{M}' \otimes_\mathcal{A} \mathcal{N}}(f(x) \otimes y)
- (-1)^n (f \otimes \text{id}_\mathcal{N})(
\text{d}_{\mathcal{M} \otimes_\mathcal{A} \mathcal{N}}(x \otimes y) \\
& =
\text{d}_{\mathcal{M}'}(f(x)) \otimes y +
(-1)^{\deg(x) + n}f(x) \otimes \text{d}_\mathcal{N}(y) \\
&
-(-1)^n f(\text{d}_\mathcal{M}(x)) \otimes y
-(-1)^n (-1)^{\deg(x)}f(x) \otimes \text{d}_\mathcal{N}(y)
\end{align*}
를 얻는다. 이는 실제로 같은 국소 단면이다.
% P03-STACKS-E1044
% P03-STACKS-E1045










\section{미분 등급 가군의 층에 대한 내부 Hom}
\label{sdga-section-internal-hom-dg}

\noindent
절 \ref{sdga-section-dgm-dg-cat}의 구성을 층화한 판본이 필요하다.
$(\mathcal{C}, \mathcal{O})$, $\mathcal{A}$,
$\mathcal{M}$, $\mathcal{L}$을 절 \ref{sdga-section-dgm-dg-cat}에서와 같이 잡자.
다음을 정의한다.
$$
\SheafHom^{dg}_\mathcal{A}(\mathcal{M}, \mathcal{L}) =
\SheafHom^{gr}_\mathcal{A}(\mathcal{M}, \mathcal{L}) =
\bigoplus\nolimits_{n \in \mathbf{Z}}
\SheafHom^n_\mathcal{A}(\mathcal{M}, \mathcal{L})
$$
이는 등급 $\mathcal{O}$-가군의 등식이다. 절
\ref{sdga-section-internal-hom-graded}를 보라. 달리 말하면, $U$ 위의
$n$차 등급 성분 $\SheafHom^n_\mathcal{A}(\mathcal{L}, \mathcal{M})$의
단면 $f$는 차수 $n$인 동차 오른쪽 $\mathcal{A}_U$-가군 사상
$\mathcal{L}|_U \to \mathcal{M}|_U$이다. 이러한 $f$에 대하여
다음과 같이 놓는다.
% P03-STACKS-E1046
$$
\text{d}(f) =
\text{d}_\mathcal{M}|_U \circ f - (-1)^n f \circ \text{d}_\mathcal{L}|_U
$$
이를 해석하기 위해 $\text{d}_\mathcal{M}|_U$와
$\text{d}_\mathcal{L}|_U$를 등급 $\mathcal{O}_U$-가군 사상으로 보고,
등급 $\mathcal{O}_U$-가군 사상의 합성을 사용한다. $\text{d}(f)$가 등급
$\mathcal{O}_U$-가군 사상으로서 차수 $n + 1$인 동차 사상임은 분명하다.
절 \ref{sdga-section-dgm-dg-cat}와 완전히 같은 계산을 사용하면
$\text{d}(f)$가 $\mathcal{A}_U$-선형임을 알 수 있다.

\medskip\noindent
절 \ref{sdga-section-dgm-dg-cat}에서와 같이 다음 합성 사상이 있다.
$$
\SheafHom^{dg}_\mathcal{A}(\mathcal{L}, \mathcal{M}) \otimes_\mathcal{O}
\SheafHom^{dg}_\mathcal{A}(\mathcal{K}, \mathcal{L})
\longrightarrow
\SheafHom^{dg}_\mathcal{A}(\mathcal{K}, \mathcal{M})
$$
여기서 왼쪽 항은 절 \ref{sdga-section-tensor-product-dg}에서 정의한
미분 등급 $\mathcal{O}$-가군의 텐서곱이다. 이 사상은 다음 합성 사상으로
주어진다.
$$
\SheafHom^m(\mathcal{L}, \mathcal{M}) \otimes_\mathcal{O}
\SheafHom^n(\mathcal{K}, \mathcal{L}) \longrightarrow
\SheafHom^{n + m}(\mathcal{K}, \mathcal{M})
$$
이는 (국소적으로) 통상적인 합성으로 정의된다. 국소 단면에 대하여
절 \ref{sdga-section-dgm-dg-cat}와 완전히 같은 계산을 사용하면 이 합성 사상이
미분 등급 $\mathcal{O}$-가군의 사상임을 알 수 있다.

\medskip\noindent
이 정의들에 따라
$$
\Hom_{\textit{Mod}^{dg}(\mathcal{A})}(\mathcal{L}, \mathcal{M}) =
\Gamma(\mathcal{C}, \SheafHom^{dg}_\mathcal{A}(\mathcal{L}, \mathcal{M}))
$$
가 합성과 양립하는 등급 $R$-가군으로서 성립한다.









\section{미분 등급 쌍가군의 층과 텐서--Hom 수반}
\label{sdga-section-dg-bimodules}

\noindent
이 절은 Differential Graded Algebra, Section
\ref{dga-section-tensor-product}의 일부에 대응한다.

\begin{definition}
\label{sdga-definition-dg-bimodule}
$(\mathcal{C}, \mathcal{O})$를 환 달린 사이트라 하자. $\mathcal{A}$와
$\mathcal{B}$를 $(\mathcal{C}, \mathcal{O})$ 위의 미분 등급 대수의 층이라 하자.
% P03-STACKS-E1047: six declarations say "a sheaves" in the frozen source.
{\it 미분 등급 $(\mathcal{A}, \mathcal{B})$-쌍가군}이란
$\mathcal{O}$-쌍선형 사상
$$
\mathcal{M}^n \times \mathcal{B}^m \to \mathcal{M}^{n + m},\quad
(x, b) \longmapsto xb
$$
과
$$
\mathcal{A}^n \times \mathcal{M}^m \to \mathcal{M}^{n + m},\quad
(a, x) \longmapsto ax
$$
가 주어진 $\mathcal{O}$-가군의 복합체 $\mathcal{M}^\bullet$을 말한다.
이 사상들을 곱셈 사상이라 부르며, 다음 성질을 만족해야 한다.
\begin{enumerate}
\item 곱셈은 $a(a'x) = (aa')x$ 및 $(xb)b' = x(bb')$을 만족한다.
\item $(ax)b = a(xb)$,
\item $\text{d}(ax) = \text{d}(a) x + (-1)^{\deg(a)}a \text{d}(x)$ 및
$\text{d}(xb) = \text{d}(x) b + (-1)^{\deg(x)}x \text{d}(b)$,
\item $\mathcal{A}^0$의 항등 단면 $1$은 곱셈으로 항등적으로 작용한다.
\item $\mathcal{B}^0$의 항등 단면 $1$은 곱셈으로 항등적으로 작용한다.
\end{enumerate}
이러한 구조를 흔히 $\mathcal{M}$으로 나타내고, 때로는
${}_\mathcal{A}\mathcal{M}_\mathcal{B}$로 쓴다.
미분 등급 $(\mathcal{A}, \mathcal{B})$-쌍가군의 {\it 준동형}
$f : \mathcal{M} \to \mathcal{N}$은 곱셈 사상과 양립하는
$\mathcal{O}$-가군 복합체의 사상
$f : \mathcal{M}^\bullet \to \mathcal{N}^\bullet$이다.
\end{definition}

\noindent
미분 등급 $(\mathcal{A}, \mathcal{B})$-쌍가군 $\mathcal{M}$과 대상
$U \in \Ob(\mathcal{C})$가 주어지면 다음 표기를 사용한다.
$$
\mathcal{M}(U) =
\Gamma(U, \mathcal{M}) =
\bigoplus\nolimits_{n \in \mathbf{Z}} \mathcal{M}^n(U)
$$
이는 미분 등급 $(\mathcal{A}(U), \mathcal{B}(U))$-쌍가군이다.

\medskip\noindent
미분 등급 $(\mathcal{A}, \mathcal{B})$-쌍가군 $\mathcal{M}$은 등급과 미분이
일치하고 $\mathcal{A}$-가군 구조와 $\mathcal{B}$-가군 구조가 서로 교환하는,
왼쪽 미분 등급 $\mathcal{A}$-가군이기도 한 오른쪽 미분 등급
$\mathcal{B}$-가군과 같은 자료임을 관찰하자. 이를 정확히 서술하면 다음과 같다.

\begin{lemma}
\label{sdga-lemma-what-makes-a-bimodule-dg}
$(\mathcal{C}, \mathcal{O})$를 환 달린 사이트라 하자. $\mathcal{A}$와
$\mathcal{B}$를 $(\mathcal{C}, \mathcal{O})$ 위의 미분 등급 대수의 층이라 하자.
$\mathcal{N}$을 오른쪽 미분 등급 $\mathcal{B}$-가군이라 하자. 주어진 미분 등급
$\mathcal{B}$-가군 구조와 양립하는 $\mathcal{N}$ 위의
$(\mathcal{A}, \mathcal{B})$-쌍가군 구조와 미분 등급
$\mathcal{O}$-대수의 준동형
$$
\mathcal{A}
\longrightarrow
\SheafHom^{dg}_\mathcal{B}(\mathcal{N}, \mathcal{N})
$$
사이에는 $1$ 대 $1$ 대응이 있다.
\end{lemma}

\begin{proof}
생략한다.
\end{proof}

\noindent
$(\mathcal{C}, \mathcal{O})$를 환 달린 사이트라 하자. $\mathcal{A}$와
$\mathcal{B}$를 $(\mathcal{C}, \mathcal{O})$ 위의 미분 등급 대수의 층이라 하자.
$\mathcal{M}$을 오른쪽 미분 등급 $\mathcal{A}$-가군이라 하고,
$\mathcal{N}$을 미분 등급 $(\mathcal{A}, \mathcal{B})$-쌍가군이라 하자.
이 경우 Section \ref{sdga-section-tensor-product-dg}에서 정의한 미분 등급 텐서곱
$$
\mathcal{M} \otimes_\mathcal{A} \mathcal{N}
$$
은 Section \ref{sdga-section-graded-bimodules}에서와 같은 곱셈 사상을 갖는
오른쪽 미분 등급 $\mathcal{B}$-가군이다. 이 구성은 함자 및 등급 범주들의 함자
$$
\otimes_\mathcal{A} \mathcal{N} :
\textit{Mod}(\mathcal{A}, \text{d})
\longrightarrow
\textit{Mod}(\mathcal{B}, \text{d})
\quad\text{and}\quad
\otimes_\mathcal{A} \mathcal{N} :
\textit{Mod}^{dg}(\mathcal{A}, \text{d})
\longrightarrow
\textit{Mod}^{dg}(\mathcal{B}, \text{d})
$$
를 정의한다. 즉 $\mathcal{M} \to \mathcal{M}'$인 차수 $n$의 준동형을
$\mathcal{M} \otimes_\mathcal{A} \mathcal{N}$에서
$\mathcal{M}' \otimes_\mathcal{A} \mathcal{N}$으로 가는 유도된 차수 $n$의
사상으로 보낸다.

\medskip\noindent
$(\mathcal{C}, \mathcal{O})$를 환 달린 사이트라 하자. $\mathcal{A}$와
$\mathcal{B}$를 $(\mathcal{C}, \mathcal{O})$ 위의 미분 등급 대수의 층이라 하자.
$\mathcal{N}$을 미분 등급 $(\mathcal{A}, \mathcal{B})$-쌍가군이라 하고,
$\mathcal{L}$을 오른쪽 미분 등급 $\mathcal{B}$-가군이라 하자. 이 경우
Section \ref{sdga-section-internal-hom-dg}에서 정의한 미분 등급 내부 Hom
$$
\SheafHom_\mathcal{B}^{dg}(\mathcal{N}, \mathcal{L})
$$
은 오른쪽 미분 등급 $\mathcal{A}$-가군이며, 그 오른쪽 등급
$\mathcal{A}$-가군 구조는 Section \ref{sdga-section-graded-bimodules}에서
정의한 것이다. 곱셈은 다음 합성을 사용하여 정의할 수도 있다.
% P03-STACKS-E1048: frozen source says "is the use the composition".
$$
\SheafHom_\mathcal{B}^{dg}(\mathcal{N}, \mathcal{L})
\otimes_\mathcal{O} \mathcal{A}
\to
\SheafHom_\mathcal{B}^{dg}(\mathcal{N}, \mathcal{L})
\otimes_\mathcal{O} \SheafHom_\mathcal{B}^{dg}(\mathcal{N}, \mathcal{N})
\to
\SheafHom_\mathcal{B}^{dg}(\mathcal{N}, \mathcal{L})
$$
여기서 첫 번째 화살표는 Lemma \ref{sdga-lemma-what-makes-a-bimodule-dg}에서
오고, 두 번째 화살표는 Section \ref{sdga-section-internal-hom-dg}의 합성이다.
두 화살표 모두 미분과 양립하므로 실제로 미분 등급 $\mathcal{A}$-가군을
얻는다. 이 구성은 함자 및 미분 등급 범주들의 함자
$$
\SheafHom_\mathcal{B}^{dg}(\mathcal{N}, -) :
\textit{Mod}(\mathcal{B}, \text{d})
\longrightarrow
\textit{Mod}(\mathcal{A})
% P03-STACKS-E1049: frozen first target omits the differential datum.
\quad\text{and}\quad
\SheafHom_\mathcal{B}^{dg}(\mathcal{N}, -) :
\textit{Mod}^{dg}(\mathcal{B}, \text{d})
\longrightarrow
\textit{Mod}^{dg}(\mathcal{A}, \text{d})
$$
를 정의한다. 즉 $\mathcal{L} \to \mathcal{L}'$인 차수 $n$의 준동형을
$\SheafHom_\mathcal{B}^{dg}(\mathcal{N}, \mathcal{L})$에서
$\SheafHom_\mathcal{B}^{dg}(\mathcal{N}, \mathcal{L}')$로 가는 유도된
차수 $n$의 사상으로 보낸다.

\begin{lemma}
\label{sdga-lemma-tensor-hom-adjunction-dg}
$(\mathcal{C}, \mathcal{O})$를 환 달린 사이트라 하자. $\mathcal{A}$와
$\mathcal{B}$를 $(\mathcal{C}, \mathcal{O})$ 위의 미분 등급 대수의 층이라 하자.
$\mathcal{M}$을 오른쪽 미분 등급 $\mathcal{A}$-가군,
$\mathcal{N}$을 미분 등급 $(\mathcal{A}, \mathcal{B})$-쌍가군,
$\mathcal{L}$을 오른쪽 미분 등급 $\mathcal{B}$-가군이라 하자. 위의 규약에 따라
$$
\Hom_{\textit{Mod}^{dg}(\mathcal{B}, \text{d})}(
\mathcal{M} \otimes_\mathcal{A} \mathcal{N}, \mathcal{L}) =
\Hom_{\textit{Mod}^{dg}(\mathcal{A}, \text{d})}(
\mathcal{M}, \SheafHom_\mathcal{B}^{dg}(\mathcal{N}, \mathcal{L}))
$$
그리고
$$
\SheafHom_\mathcal{B}^{dg}(
\mathcal{M} \otimes_\mathcal{A} \mathcal{N}, \mathcal{L}) =
\SheafHom_\mathcal{A}^{dg}(
\mathcal{M}, \SheafHom_\mathcal{B}^{dg}(\mathcal{N}, \mathcal{L}))
$$
가 성립하며, 이는 $\mathcal{M}$, $\mathcal{N}$, $\mathcal{L}$에 관하여
함자적이다.
\end{lemma}

\begin{proof}
생략한다. 힌트: 등급 수준에서는 Lemma
\ref{sdga-lemma-tensor-hom-adjunction-gr}에서 이미 이를 보았다. 따라서 이 동형들이
미분과 양립함을 확인하면 충분하고, 이는 국소 단면 수준의 계산으로 할 수 있다.
\end{proof}

\noindent
$(\mathcal{C}, \mathcal{O})$를 환 달린 사이트라 하자. $\mathcal{A}$와
$\mathcal{B}$를 $(\mathcal{C}, \mathcal{O})$ 위의 미분 등급 대수의 층이라 하자.
위의 특수한 경우로, 미분 등급 $\mathcal{O}$-대수의 준동형
$\varphi : \mathcal{A} \to \mathcal{B}$가 주어졌다고 하자. 그러면 함자 및
미분 등급 범주들의 함자
$$
\otimes_{\mathcal{A}, \varphi} \mathcal{B} :
\textit{Mod}(\mathcal{A}, \text{d})
\longrightarrow
\textit{Mod}(\mathcal{B}, \text{d})
\quad\text{and}\quad
\otimes_{\mathcal{A}, \varphi} \mathcal{B} :
\textit{Mod}^{dg}(\mathcal{A}, \text{d})
\longrightarrow
\textit{Mod}^{dg}(\mathcal{B}, \text{d})
$$
를 얻는다. 한편 제한 함자
$$
res_\varphi :
\textit{Mod}(\mathcal{B}, \text{d})
\longrightarrow
\textit{Mod}(\mathcal{A}, \text{d})
\quad\text{and}\quad
res_\varphi :
\textit{Mod}^{dg}(\mathcal{B}, \text{d})
\longrightarrow
\textit{Mod}^{dg}(\mathcal{A}, \text{d})
$$
도 있다. 위의 보조정리를 사용하면 이 함자들이 서로 수반임을 보일 수 있다
(제한과 밑변환에서 통상적인 대로이다). 구체적으로, $\mathcal{B}$를 미분 등급
$(\mathcal{A}, \mathcal{B})$-쌍가군으로 본 것을
${}_\mathcal{A}\mathcal{B}_\mathcal{B}$로 쓰자. 그러면 임의의 오른쪽 미분 등급
$\mathcal{B}$-가군 $\mathcal{L}$에 대하여
$$
\SheafHom_\mathcal{B}^{dg}({}_\mathcal{A}\mathcal{B}_\mathcal{B}, \mathcal{L})
= res_\varphi(\mathcal{L})
$$
인 오른쪽 미분 등급 $\mathcal{A}$-가군의 등식이 성립한다. 따라서
Lemma \ref{sdga-lemma-tensor-hom-adjunction-gr}에 의해 함자적 동형
% P03-STACKS-E1050: frozen source invokes the graded rather than DG adjunction lemma.
$$
\Hom_{\textit{Mod}^{dg}(\mathcal{B}, \text{d})}(
\mathcal{M} \otimes_{\mathcal{A}, \varphi} \mathcal{B}, \mathcal{L}) =
\Hom_{\textit{Mod}^{dg}(\mathcal{A}, \text{d})}(
\mathcal{M}, res_\varphi(\mathcal{L}))
$$
을 얻는다. 문맥상 명확하면 이 식에서 $\varphi$에 대한 의존성을 보통 생략한다.
같은 방식으로 등급 $\mathcal{O}$-가군의 등식
$$
\SheafHom^{dg}_\mathcal{B}(
\mathcal{M} \otimes_\mathcal{A} \mathcal{B}, \mathcal{L}) =
\SheafHom_\mathcal{A}^{dg}(\mathcal{M}, \mathcal{L})
$$
을 얻는다.












\section{미분 등급 가군의 층에 대한 당김과 밂}
\label{sdga-section-functoriality-dg}

\noindent
$(f, f^\sharp) : (\Sh(\mathcal{C}), \mathcal{O}_\mathcal{C})
\to (\Sh(\mathcal{D}), \mathcal{O}_\mathcal{D})$가 환 달린 토포스의 사상이라 하자.
$\mathcal{A}$를 미분 등급 $\mathcal{O}_\mathcal{C}$-대수라 하고,
$\mathcal{B}$를 미분 등급 $\mathcal{O}_\mathcal{D}$-대수라 하자.
다음 사상이 주어졌다고 하자.
$$
\varphi : f^{-1}\mathcal{B} \to \mathcal{A}
$$
는 미분 등급 $f^{-1}\mathcal{O}_\mathcal{D}$-대수의 사상이다.
스칼라 제한과 스칼라 확장의 수반에 의해 이는 미분 등급
$\mathcal{O}_\mathcal{C}$-대수의 사상
$\varphi : f^*\mathcal{B} \to \mathcal{A}$와 같은 것이며, 동치로
$\varphi$를 다음 사상으로 볼 수 있다.
$$
\varphi : \mathcal{B} \to f_*\mathcal{A}
$$
미분 등급 $\mathcal{O}_\mathcal{D}$-대수의 사상으로 볼 수 있다.
Remark \ref{sdga-remark-functoriality-dga}를 보라.

\medskip\noindent
함자를 정의하자.
$$
f_* :
\textit{Mod}(\mathcal{A}, \text{d})
\longrightarrow
\textit{Mod}(\mathcal{B}, \text{d})
$$
미분 등급 $\mathcal{A}$-가군 $\mathcal{M}$에 대하여
$f_*\mathcal{M}$을 Section \ref{sdga-section-functoriality-graded}에서 구성한
등급 $\mathcal{B}$-가군으로 정의하고, 미분은
$f_*d : f_*\mathcal{M}^n \to f_*\mathcal{M}^{n + 1}$로 주어진다.
이 구성은 $\mathcal{M}$에 관하여 명백히 함자적이므로 원하는 함자를 얻는다.

\medskip\noindent
함자를 정의하자.
$$
f^* :
\textit{Mod}(\mathcal{B}, \text{d})
\longrightarrow
\textit{Mod}(\mathcal{A}, \text{d})
$$
미분 등급 $\mathcal{B}$-가군 $\mathcal{N}$에 대하여
$f^*\mathcal{N}$을 Section \ref{sdga-section-functoriality-graded}에서 구성한
등급 $\mathcal{A}$-가군으로 정의한다. 다음을 상기하자.
$$
f^*\mathcal{N} = f^{-1}\mathcal{N} \otimes_{f^{-1}\mathcal{B}} \mathcal{A}
$$
 $f^{-1}\mathcal{N}$에는 미분
$f^{-1}\text{d} : f^{-1}\mathcal{N}^n \to f^{-1}\mathcal{N}^{n + 1}$
가 주어져 있으므로, 이 텐서곱을 미분을 제공하는
Section \ref{sdga-section-dg-bimodules}의 텐서곱의 한 예로 볼 수 있다.
이 구성은 $\mathcal{N}$에 관하여 명백히 함자적이므로 함자 $f^*$를 얻는다.

\medskip\noindent
$f_*$와 $f^*$는 쉽게 미분 등급 범주들의 함자로 확장된다.
$$
f_* :
\textit{Mod}^{dg}(\mathcal{A}, \text{d})
\longrightarrow
\textit{Mod}^{dg}(\mathcal{B}, \text{d})
\quad\text{and}\quad
f^* :
\textit{Mod}^{dg}(\mathcal{B}, \text{d})
\longrightarrow
\textit{Mod}^{dg}(\mathcal{A}, \text{d})
$$
이들은 바탕 대상에 대해서는 같은 작용을 하며, 차수 $n$의 동차 사상에
대한 구성의 함자성으로 정의된다.

\begin{lemma}
\label{sdga-lemma-adjunction-push-pull-dg}
위 상황에서 다음이 성립한다.
$$
\Hom_{\textit{Mod}^{dg}(\mathcal{B}, \text{d})}(
\mathcal{N}, f_*\mathcal{M}) =
\Hom_{\textit{Mod}^{dg}(\mathcal{A}, \text{d})}(
f^*\mathcal{N}, \mathcal{M})
$$
\end{lemma}

\begin{proof}
생략한다. 힌트: 바탕 등급 범주에 대해서는 Lemma
\ref{sdga-lemma-adjunction-push-pull-gr}에 의해 참이다. 계산하면 이 동형들이
미분과 양립함을 알 수 있다.
\end{proof}















\section{국소화와 미분 등급 가군의 층}
\label{sdga-section-localize-dg}

\noindent
$(\mathcal{C}, \mathcal{O})$를 환 달린 사이트라 하자.
$U \in \Ob(\mathcal{C})$라 하고 다음을 나타내자.
$$
j :
(\Sh(\mathcal{C}/U), \mathcal{O}_U)
\longrightarrow
(\Sh(\mathcal{C}), \mathcal{O})
$$
를 대응하는 국소화 사상이라 하자
(Modules on Sites, Section \ref{sites-modules-section-localize}).
이하에서 다음 사실을 사용한다. $\mathcal{O}_U$-가군
$\mathcal{M}_i$ ($i = 1, 2$)와 $\mathcal{O}$-가군 $\mathcal{A}$에 대하여
표준 사상
$$
j_! :
\Hom_{\mathcal{O}_U}(
\mathcal{M}_1 \otimes_{\mathcal{O}_U} \mathcal{A}|_U, \mathcal{M}_2)
\longrightarrow
\Hom_\mathcal{O}(
j_!\mathcal{M}_1 \otimes_\mathcal{O} \mathcal{A}, j_!\mathcal{M}_2)
$$
이 존재한다. 실제로
$j_!(\mathcal{M}_1 \otimes_{\mathcal{O}_U} \mathcal{A}|_U) =
j_!\mathcal{M}_1 \otimes_\mathcal{O} \mathcal{A}$이고, 이는
Modules on Sites, 보조정리 \ref{sites-modules-lemma-j-shriek-and-tensor}에 의한다.

\medskip\noindent
 $\mathcal{A}$를 미분 등급 $\mathcal{O}$-대수라 하자.
$\mathcal{A}$를 $\mathcal{C}/U$로 제한한 것을 $\mathcal{A}_U$로 나타내며,
즉
$\mathcal{A}_U = j^*\mathcal{A} = j^{-1}\mathcal{A}$.
Section \ref{sdga-section-functoriality-dg}에서 수반 함자
$$
j_* :
\textit{Mod}^{dg}(\mathcal{A}_U, \text{d})
\longrightarrow
\textit{Mod}^{dg}(\mathcal{A}, \text{d})
\quad\text{and}\quad
j^* :
\textit{Mod}^{dg}(\mathcal{A}, \text{d})
\longrightarrow
\textit{Mod}^{dg}(\mathcal{A}_U, \text{d})
$$
$j^*$가 $j_*$의 왼쪽 수반이다. 여기에 더하여 다음 완전 함자가 존재한다.
$$
j_! :
\textit{Mod}^{dg}(\mathcal{A}_U, \text{d})
\longrightarrow
\textit{Mod}^{dg}(\mathcal{A}, \text{d})
$$
이는 $j_*$의 오른쪽 수반이다. 즉 미분 등급
$\mathcal{A}_U$-가군 $\mathcal{M}$이 주어지면 $j_!\mathcal{M}$을
Section \ref{sdga-section-localize-graded}에서 구성한 등급 $\mathcal{A}$-가군으로
정의하고, 미분은
$j_!\text{d} : j_!\mathcal{M}^n \to j_!\mathcal{M}^{n + 1}$.
로 둔다. 미분 등급 $\mathcal{A}_U$-가군의 차수 $n$ 동차 사상
$f : \mathcal{M} \to \mathcal{M}'$이 주어지면, 미분 등급
$\mathcal{A}$-가군의 차수 $n$ 동차 사상
$j_!f : j_!\mathcal{M} \to j_!\mathcal{M}'$을 얻는다.
이 구성이 미분과 양립한다는 직접적인 확인은 생략한다.
따라서 원하는 함자를 얻는다.

\begin{lemma}
\label{sdga-lemma-extension-by-zero-dg}
위 상황에서 다음이 성립한다.
$$
\Hom_{\textit{Mod}^{dg}(\mathcal{A}, \text{d})}(
j_!\mathcal{M}, \mathcal{N}) =
\Hom_{\textit{Mod}^{dg}(\mathcal{A}_U, \text{d})}(
\mathcal{M}, j^*\mathcal{N})
$$
\end{lemma}

\begin{proof}
생략한다. 힌트: Lemma \ref{sdga-lemma-extension-by-zero-graded}에서 이 보조정리가
등급 수준에서 참임을 보았다. 그러므로 결과 동형이 미분과 양립한다는 것만
확인하면 된다.
\end{proof}

\begin{lemma}
\label{sdga-lemma-tensor-with-extension-by-zero-dg}
위 상황에서 $\mathcal{M}$을 오른쪽 미분 등급
$\mathcal{A}_U$-가군이라 하고 $\mathcal{N}$을 왼쪽 미분 등급
$\mathcal{A}$-가군이라 하자. 그러면
$$
j_!\mathcal{M} \otimes_\mathcal{A} \mathcal{N} =
j_!(\mathcal{M} \otimes_{\mathcal{A}_U} \mathcal{N}|_U)
$$
가 성립하고, 이는 $\mathcal{O}$-가군 복합체의 등식이며
$\mathcal{M}$과 $\mathcal{N}$에 관하여 함자적이다.
\end{lemma}

\begin{proof}
등급 가군으로서는 Lemma \ref{sdga-lemma-tensor-with-extension-by-zero}에서 따른다.
이 동형이 미분과 양립한다는 확인은 생략한다.
\end{proof}








\section{미분 등급 가군의 층의 이동 함자}
\label{sdga-section-shift-dg}

\noindent
$(\mathcal{C}, \mathcal{O})$를 환 달린 사이트라 하자.
$\mathcal{A}$를 $(\mathcal{C}, \mathcal{O})$ 위의 미분 등급 대수의 층이라 하자.
$\mathcal{M}$을 미분 등급 $\mathcal{A}$-가군이라 하자.
$k \in \mathbf{Z}$라 하자. $\mathcal{M}$의 {\it $k$번째 이동}을
$\mathcal{M}[k]$으로 나타내고 다음과 같이 정의한다.
\begin{enumerate}
 \item 등급 $\mathcal{A}$-가군으로서 $\mathcal{M}[k]$를 Section
\ref{sdga-section-shift}에서 정의한 것으로 둔다.
 \item 미분
$d_{\mathcal{M}[k]} : (\mathcal{M}[k])^n \to (\mathcal{M}[k])^{n + 1}$
은 다음으로 정의한다.
$(-1)^k\text{d}_\mathcal{M} : \mathcal{M}^{n + k} \to \mathcal{M}^{n + k + 1}$.
\end{enumerate}
$\mathcal{A}$-가군의 차수 $n$ 동차 준동형
$f : \mathcal{L} \to \mathcal{M}$에 대하여 $f[k] :
\mathcal{L}[k] \to \mathcal{M}[k]$를 $f$와 같은 성분 사상으로 정의한다.
그러면 $f[k]$는 차수 $n$의 동차 $\mathcal{A}$-가군 사상이다.
따라서 다음 사상을 얻는다.
$$
\Hom_{\textit{Mod}^{dg}(\mathcal{A}, \text{d})}(\mathcal{L}, \mathcal{M})
\longrightarrow
\Hom_{\textit{Mod}^{dg}(\mathcal{A}, \text{d})}
(\mathcal{L}[k], \mathcal{M}[k])
$$
이며, 이는 미분과 양립한다 (이는 $\mathcal{L}$과 $\mathcal{M}$의 미분 부호가
같은 만큼 변한다는 사실에서 따른다). 이 선택은 Differential Graded Algebra,
Definition \ref{dga-definition-shift}의 선택과 양립한다.
따라서 다음 함자를 정의했음이 명백하다.
$$
[k] :
\textit{Mod}^{dg}(\mathcal{A}, \text{d})
\longrightarrow
\textit{Mod}^{dg}(\mathcal{A}, \text{d})
$$
미분 등급 범주들의 함자이며,
$[k + l] = [k] \circ [l]$.

\medskip\noindent
등급 가군 위에서 Section \ref{sdga-section-shift}에서 정의한 동형
$$
\Hom_{\textit{Mod}^{dg}(\mathcal{A}, \text{d})}(\mathcal{L}, \mathcal{M}[k])
=
\Hom_{\textit{Mod}^{dg}(\mathcal{A}, \text{d})}(\mathcal{L}, \mathcal{M})[k]
$$
은 미분과 양립함을 주장한다. 이를 보이기 위해 차수 $n$ 동차 오른쪽
$\mathcal{A}$-가군 사상 $f : \mathcal{L} \to \mathcal{M}[k]$를 생각하자.
이는 좌변의 차수 $n$ 원소이다. $f$와 {\bf 같은} 성분 사상을 갖는
차수 $n+k$의 동차 $\mathcal{A}$-가군 사상
$f' : \mathcal{L} \to \mathcal{M}$을 택하자.
우리의 규약에 따르면 이는 우변의 대응하는 차수 $n$ 원소이다.
좌변의 미분 정의에 의해
$$
\text{d}_{LHS}(f) =
\text{d}_{\mathcal{M}[k]} \circ f - (-1)^n f \circ \text{d}_\mathcal{L} =
(-1)^k\text{d}_\mathcal{M} \circ f - (-1)^n f \circ \text{d}_\mathcal{L}
$$
그리고 우변의 미분에 대해서는 다음을 얻는다.
$$
\text{d}_{RHS}(f') =
(-1)^k\left(
\text{d}_\mathcal{M} \circ f' - (-1)^{n + k} f' \circ \text{d}_\mathcal{L}
\right) =
(-1)^k\text{d}_\mathcal{M} \circ f' - (-1)^n f' \circ \text{d}_\mathcal{L}
$$
이고, 이 사상들은 같은 성분 사상을 가지므로 증명이 끝난다.










\section{호모토피 범주}
\label{sdga-section-homotopy}

\noindent
이 절은 Differential Graded Algebra, Section
\ref{dga-section-homotopy}에 대응한다.

\begin{definition}
\label{sdga-definition-homotopy}
$(\mathcal{C}, \mathcal{O})$를 환 달린 사이트라 하자.
$\mathcal{A}$를 $(\mathcal{C}, \mathcal{O})$ 위의 미분 등급 대수의 층이라 하자.
$f, g : \mathcal{M} \to \mathcal{N}$을 미분 등급 $\mathcal{A}$-가군의
준동형이라 하자. $f$와 $g$ 사이의 {\it 호모토피}란 차수 $-1$의
동차 등급 $\mathcal{A}$-가군 사상 $h : \mathcal{M} \to \mathcal{N}$으로서
$$
f - g = \text{d}_\mathcal{N} \circ h + h \circ \text{d}_\mathcal{M}
$$
를 만족하는 것이다. 호모토피가 존재하면 $f$와 $g$가 {\it 호모토픽}이라고
말한다.
\end{definition}

\noindent
정의의 상황에서 사상 $a : \mathcal{K} \to \mathcal{M}$과
$c : \mathcal{N} \to \mathcal{L}$이 주어지면 다음을 알 수 있다.
$$
\begin{matrix}
h\text{는 호모토피이다} \\
f\text{와 }g\text{ 사이의}
\end{matrix}
\quad
\Rightarrow
\quad
\begin{matrix}
c \circ h \circ a\text{는 호모토피이다}\\
c \circ f \circ a\text{와 }c\circ g\circ a\text{ 사이의}
\end{matrix}
$$
따라서 $\textit{Mod}(\mathcal{A}, \text{d})$에서 사상의 호모토피류의
합성을 정의할 수 있다.

\begin{definition}
\label{sdga-definition-complexes-notation}
$(\mathcal{C}, \mathcal{O})$를 환 달린 사이트라 하자.
$\mathcal{A}$를 $(\mathcal{C}, \mathcal{O})$ 위의 미분 등급 대수의 층이라 하자.
$K(\textit{Mod}(\mathcal{A}, \text{d}))$로 나타내는 {\it 호모토피 범주}는
대상이 $\textit{Mod}(\mathcal{A}, \text{d})$의 대상이고 사상이 미분 등급
$\mathcal{A}$-가군 준동형의 호모토피류인 범주이다.
\end{definition}

\noindent
표기 $K(\textit{Mod}(\mathcal{A}, \text{d}))$는 표준적이지는 않지만
Stacks project의 다른 곳에서 $K(-)$를 사용하는 것과는 일관된다.

\medskip\noindent
Differential Graded Algebra, Definition
\ref{dga-definition-homotopy-category-of-dga-category}에서 복합체의 범주
$\text{Comp}(\mathcal{S})$
와 미분 등급 범주 $\mathcal{S}$의 호모토피 범주 $K(\mathcal{S})$를 정의했다.
이를 미분 등급 범주 $\textit{Mod}^{dg}(\mathcal{A}, \text{d})$에 적용하면
$$
\textit{Mod}(\mathcal{A}, \text{d}) =
\text{Comp}(\textit{Mod}^{dg}(\mathcal{A}, \text{d}))
$$
(Section \ref{sdga-section-dgm-dg-cat}의 논의를 보라). 또한
$$
K(\textit{Mod}(\mathcal{A}, \text{d})) =
K(\textit{Mod}^{dg}(\mathcal{A}, \text{d}))
$$
를 얻는다. 마지막 등식이 참임을 보이기 위해 다음 등식에 주목하자.
$$
\text{d}_{\Hom_{\textit{Mod}^{dg}(\mathcal{A}, \text{d})}
(\mathcal{M}, \mathcal{N})}(h) =
\text{d}_\mathcal{N} \circ h + h \circ \text{d}_\mathcal{M}
$$
$h$가 Definition \ref{sdga-definition-homotopy}에서와 같을 때 성립한다.
세부사항은 생략한다.

\begin{lemma}
\label{sdga-lemma-homotopy-direct-sums}
$(\mathcal{C}, \mathcal{O})$를 환 달린 사이트라 하자.
$\mathcal{A}$를 $(\mathcal{C}, \mathcal{O})$ 위의 미분 등급 대수의 층이라 하자.
호모토피 범주 $K(\textit{Mod}(\mathcal{A}, \text{d}))$는 직합과 곱을 갖는다.
\end{lemma}

\begin{proof}
생략한다. 힌트: Lemma \ref{sdga-lemma-dgm-abelian}에서와 같이 직합과 곱을
사용하면 된다. 이는 이 함자들이 등급 $\mathcal{A}$-가군의 범주로 가는
망각 함자와 가환하고, 아벨 군의 족의 범주에서 $\prod$와 $\bigoplus$가
완전 함자이기 때문이다.
\end{proof}










\section{원뿔과 삼각형}
\label{sdga-section-conclude-triangulated}

\noindent
이 절에서는 Differential Graded Algebra, Section
\ref{dga-section-review}의 내용을 사용하여 미분 등급
$\mathcal{A}$-가군 범주의 호모토피 범주가 삼각 범주임을 보인다.

\begin{lemma}
\label{sdga-lemma-axioms-AB}
$(\mathcal{C}, \mathcal{O})$를 환 달린 사이트라 하자.
$\mathcal{A}$를 $(\mathcal{C}, \mathcal{O})$ 위의 미분 등급 대수의 층이라 하자.
미분 등급 범주
$\textit{Mod}^{dg}(\mathcal{A}, \text{d})$
는 Differential Graded Algebra, Section \ref{dga-section-review}의
공리 (A)와 (B)를 만족한다.
\end{lemma}

\begin{proof}
미분 등급 $\mathcal{A}$-가군 $\mathcal{M}$과 $\mathcal{N}$이 주어졌다고 하자.
자명한 방식으로 정의한 미분 등급 $\mathcal{A}$-가군
$\mathcal{M} \oplus \mathcal{N}$을 생각하자. 그러면 공사상
$i : \mathcal{M} \to \mathcal{M} \oplus \mathcal{N}$ 및
$j : \mathcal{N} \to \mathcal{M} \oplus \mathcal{N}$ 및
및 사영
$p : \mathcal{M} \oplus \mathcal{N} \to \mathcal{N}$ 및
$q : \mathcal{M} \oplus \mathcal{N} \to \mathcal{M}$
는 미분 등급 $\mathcal{A}$-가군의 사상이다.
따라서 $i,j,p,q$는 차수 $0$의 동차 닫힌 사상, 즉
$\text{d}(i)=0$ 등이다. 그러므로 이 직합은 Differential Graded Algebra,
Definition \ref{dga-definition-dg-direct-sum}의 의미에서 미분 등급 직합이다.
이로써 공리 (A)가 증명된다.

\medskip\noindent
공리 (B)는 Section \ref{sdga-section-shift-dg}에서 보였다.
\end{proof}

\noindent
$(\mathcal{C}, \mathcal{O})$를 환 달린 사이트라 하자.
$\mathcal{A}$를 $(\mathcal{C}, \mathcal{O})$ 위의 미분 등급 대수의 층이라 하자.
다음 열을 생각하자.
$$
0 \to \mathcal{K} \to \mathcal{L} \to \mathcal{N} \to 0
$$
$\textit{Mod}(\mathcal{A}, \text{d})$에서 이를 {\it 허용 짧은 완전열}이라
부른다 (Differential Graded Algebra, Section \ref{dga-section-review}).
이는 $\textit{Mod}(\mathcal{A})$에서 분할되는 경우, 즉 등급
$\mathcal{A}$-가군의 열로서 분할되는 경우이다. 다음 등급
$\mathcal{A}$-가군 분할을 나타내자.
$s : \mathcal{N} \to \mathcal{L}$ 및
$\pi : \mathcal{L} \to \mathcal{K}$
Lemma \ref{sdga-lemma-axioms-AB}와 Differential Graded Algebra, Lemma
\ref{dga-lemma-get-triangle}를 결합하면 다음 삼각형을 얻는다.
$$
\mathcal{K} \to \mathcal{L} \to \mathcal{N} \to \mathcal{K}[1]
$$
여기서 Differential Graded Algebra, Lemma
\ref{dga-lemma-get-triangle}의 증명에서 화살표
$\mathcal{N} \to \mathcal{K}[1]$는 다음과 같이 구성된다.
$$
\delta = \pi \circ
\text{d}_{\Hom_{\textit{Mod}^{dg}(\mathcal{A}, \text{d})}
(\mathcal{L}, \mathcal{M})}(s) =
\pi \circ \text{d}_\mathcal{L} \circ s -
\pi \circ s \circ \text{d}_\mathcal{N} =
\pi \circ \text{d}_\mathcal{L} \circ s
$$
끔찍한 표기에 대해서는 양해를 구한다. 어쨌든 우리의 상황에서
Differential Graded Algebra, Lemma \ref{dga-lemma-get-triangle}에서 구성한
경계 사상 $\delta$는 $\mathcal{O}$-가군의 바탕 복합체에서
항별로 분할된 복합체의 짧은 완전열에 대해 Stacks project 전체에서
사용하는 통상적인 경계 사상과 일치한다. Derived Categories, Definition
\ref{derived-definition-split-ses}를 보라.

\begin{definition}
\label{sdga-definition-cone}
$(\mathcal{C}, \mathcal{O})$를 환 달린 사이트라 하자.
$\mathcal{A}$를 $(\mathcal{C}, \mathcal{O})$ 위의 미분 등급 대수의 층이라 하자.
$f : \mathcal{K} \to \mathcal{L}$을 미분 등급 $\mathcal{A}$-가군의 준동형이라 하자.
$f$의 {\it 원뿔}은 다음과 같이 정의되는 미분 등급 $\mathcal{A}$-가군
$C(f)$이다.
\begin{enumerate}
 \item $\mathcal{O}$-가군의 바탕 복합체는 $\mathcal{A}$-가군 복합체의
사상 $f : \mathcal{K}^\bullet \to \mathcal{L}^\bullet$에 대응하는 원뿔이다.
즉
$C(f)^n = \mathcal{L}^n \oplus \mathcal{K}^{n + 1}$ 및
미분은
$$
d_{C(f)} =
\left(
\begin{matrix}
\text{d}_\mathcal{L} & f \\
0 & -\text{d}_\mathcal{K}
\end{matrix}
\right)
$$
\item 곱셈 사상
$$
C(f)^n \times \mathcal{A}^m \to C(f)^{n + m}
$$
은 다음 두 곱셈 사상의 직합이다.
$\mathcal{L}^n \times \mathcal{A}^m \to \mathcal{L}^{n + m}$
및
$\mathcal{K}^{n + 1} \times \mathcal{A}^m \to \mathcal{K}^{n + 1 + m}$.
\end{enumerate}
이는 명백한 사상에서 유도되는 미분 등급 $\mathcal{A}$-가군의 표준
준동형 $i : \mathcal{L} \to C(f)$ 및 $p : C(f) \to \mathcal{K}[1]$을 갖는다.
% P03-STACKS-E1051: frozen source spells homomorphisms as "hommorphisms".
\end{definition}

\noindent
정의의 상황에서 다음 열은
$$
0 \to \mathcal{L} \to C(f) \to \mathcal{K}[1] \to 0
$$
허용 짧은 완전열임을 관찰하자.
% P03-STACKS-E1052: frozen source spells admissible as "addmissible".

\begin{lemma}
\label{sdga-lemma-axiom-C}
$(\mathcal{C}, \mathcal{O})$를 환 달린 사이트라 하자.
$\mathcal{A}$를 $(\mathcal{C}, \mathcal{O})$ 위의 미분 등급 대수의 층이라 하자.
미분 등급 범주
$\textit{Mod}^{dg}(\mathcal{A}, \text{d})$
는 Differential Graded Algebra, Situation \ref{dga-situation-ABC}에서
정식화한 공리 (C)를 만족한다.
\end{lemma}

\begin{proof}
 $f : \mathcal{K} \to \mathcal{L}$을 미분 등급 $\mathcal{A}$-가군의
준동형이라 하자. 위에서 허용 짧은 완전열
$$
0 \to \mathcal{L} \to C(f) \to \mathcal{K}[1] \to 0
$$
을 얻었다. 증명을 마치려면 이 열에 대응하는 경계 사상
$$
\delta : \mathcal{K}[1] \to \mathcal{L}[1]
$$
가 $f[1]$과 같음을 보여야 한다. 사상
$s : \mathcal{K}[1] \to C(f)$에 대해서는 차수 $n$에서 매입
$\mathcal{K}^{n + 1} \to C(f)^n$을 사용한다. 그러면 차수 $n$에서
$\pi$는 사영 $C(f)^n \to \mathcal{L}^n$으로 주어진다. 따라서
계산해야 할 것은
$$
\delta = \pi \circ \text{d}_{C(f)} \circ s
$$
(위의 논의를 보라). 행렬 표기로 이는 다음과 같다.
$$
\left(
\begin{matrix}
1 & 0
\end{matrix}
\right)
\left(
\begin{matrix}
\text{d}_\mathcal{L} & f \\
0 & -\text{d}_\mathcal{K}
\end{matrix}
\right)
\left(
\begin{matrix}
0 \\
1
\end{matrix}
\right) = f
$$
이고, 원하는 결과를 얻는다.
\end{proof}

\noindent
이제 Differential Graded Algebra, Section \ref{dga-section-review}에서 증명한
모든 보조정리가 미분 등급 범주
$\textit{Mod}^{dg}(\mathcal{A}, \text{d})$에도 성립함을 안다.
특히 다음을 얻는다.

\begin{proposition}
\label{sdga-proposition-homotopy-category-triangulated}
$(\mathcal{C}, \mathcal{O})$를 환 달린 사이트라 하자.
$\mathcal{A}$를 $(\mathcal{C}, \mathcal{O})$ 위의 미분 등급 대수의 층이라 하자.
호모토피 범주 $K(\textit{Mod}(\mathcal{A}, \text{d}))$는 삼각 범주이며,
\begin{enumerate}
\item 이동 함자는 Section \ref{sdga-section-shift-dg}에서 구성한 것이다.
\item 특별 삼각형은 $K(\textit{Mod}(\mathcal{A}, \text{d}))$ 안의 삼각형 중
다음 삼각형과 삼각형으로서 동형인 것들이다.
$$
\mathcal{K} \to \mathcal{L} \to \mathcal{N}
\xrightarrow{\delta} \mathcal{K}[1],\quad\quad
\delta = \pi \circ \text{d}_\mathcal{L} \circ s
$$
이는 위의 $\textit{Mod}(\mathcal{A}, \text{d})$에서의 허용 짧은 완전열
$0 \to \mathcal{K} \to \mathcal{L} \to \mathcal{N} \to 0$으로부터 구성된다.
\end{enumerate}
\end{proposition}

\begin{proof}
 $K(\textit{Mod}(\mathcal{A}, \text{d})) =
K(\textit{Mod}^{dg}(\mathcal{A}, \text{d}))$, Section
\ref{sdga-section-homotopy}임을 상기하자. 그러므로 이 명제는
Lemmas \ref{sdga-lemma-axioms-AB}와 \ref{sdga-lemma-axiom-C}, 그리고
Differential Graded Algebra, Proposition
\ref{dga-proposition-ABC-homotopy-category-triangulated}에서 따른다.
\end{proof}

\begin{remark}
\label{sdga-remark-cone-identity}
$(\mathcal{C}, \mathcal{O})$를 환 달린 사이트라 하자.
$\mathcal{A}$를 $(\mathcal{C}, \mathcal{O})$ 위의 미분 등급 대수의 층이라 하자.
$C = C(\text{id}_\mathcal{A})$를 미분 등급 $\mathcal{A}$-가군의 사상으로
본 항등 사상 $\mathcal{A} \to \mathcal{A}$의 원뿔이라 하자. 그러면
$$
\Hom_{\textit{Mod}(\mathcal{A}, \text{d})}(C, \mathcal{M}) =
\{(x, y) \in
\Gamma(\mathcal{C}, \mathcal{M}^0) \times
\Gamma(\mathcal{C}, \mathcal{M}^{-1}) \mid
x = \text{d}(y)\}
$$
좌변에서 우변으로 가는 사상은 $f$를 다음 쌍 $(x,y)$로 보낸다.
$x$는 $C^{-1} = \mathcal{A}^{-1} \oplus \mathcal{A}^0$의 전역 단면
$(0,1)$의 상이고, $y$는 $C^0 = \mathcal{A}^0 \oplus \mathcal{A}^1$의
전역 단면 $(1,0)$의 상이다.
\end{remark}

\begin{lemma}
\label{sdga-lemma-dgm-grothendieck-abelian}
$(\mathcal{C}, \mathcal{O})$를 환 달린 사이트라 하자.
$(\mathcal{A}, \text{d})$를 미분 등급 $\mathcal{O}$-대수라 하자.
범주 $\textit{Mod}(\mathcal{A}, \text{d})$는 그로텐디크 아벨 범주이다.
\end{lemma}

\begin{proof}
Lemma \ref{sdga-lemma-dgm-abelian}과 그로텐디크 아벨 범주의 정의
(Injectives, Definition \ref{injectives-definition-grothendieck-conditions})에
따르면 $\textit{Mod}(\mathcal{A}, \text{d})$가 생성자를 갖는지만 보이면 충분하다.
$\mathcal{C}$의 모든 대상 $U$에 대하여 Remark
\ref{sdga-remark-cone-identity}와 같이 항등 사상 $\mathcal{A}_U \to \mathcal{A}_U$의
원뿔을 $C_U$로 나타내자. 다음이 생성자임을 주장한다.
$$
\mathcal{G} = \bigoplus\nolimits_{k, U} j_{U!}C_U[k]
$$
여기서 합은 $\mathcal{C}$의 모든 대상 $U$와 $k \in \mathbf{Z}$에 걸친다.
실제로 미분 등급 $\mathcal{A}$-가군 $\mathcal{M}$에 대해
$\mathcal{G}$에서 $\mathcal{M}$으로 가는 0이 아닌 사상이 없다면,
모든 $k$와 $U$에 대해
\begin{align*}
\Hom_{\textit{Mod}(\mathcal{A})}(j_{U!}C_U[k], \mathcal{M}) \\
& =
\Hom_{\textit{Mod}(\mathcal{A}_U)}(C_U[k], \mathcal{M}|_U) \\
& =
\{(x, y) \in \mathcal{M}^{-k}(U) \times \mathcal{M}^{-k - 1}(U) \mid
x = \text{d}(y)\}
\end{align*}
가 0이다. 따라서 $\mathcal{M}$은 0이다.
\end{proof}











\section{평탄 해소}
\label{sdga-section-P-resolutions}

\noindent
이 절은 Differential Graded Algebra, Section
\ref{dga-section-P-resolutions}에 대응한다.

\medskip\noindent
$(\mathcal{C}, \mathcal{O})$를 환 달린 사이트라 하자.
$\mathcal{A}$를 $(\mathcal{C}, \mathcal{O})$ 위의 미분 등급 대수의
층이라 하자. 오른쪽 미분 등급 $\mathcal{A}$-가군
$\mathcal{P}$를 다음 조건을 만족하면 {\it good}이라고 부르자.
\begin{enumerate}
\item 함자
$\mathcal{N} \mapsto \mathcal{P} \otimes_\mathcal{A} \mathcal{N}$
가 등급 왼쪽 $\mathcal{A}$-가군의 범주에서 완전하다.
\item $\mathcal{N}$이 비순환 미분 등급 왼쪽
$\mathcal{A}$-가군이면
$\mathcal{P} \otimes_\mathcal{A} \mathcal{N}$는 비순환이다.
\item 임의의 환 달린 토포스의 사상
$(f, f^\sharp) : (\Sh(\mathcal{C}'), \mathcal{O}')
\to (\Sh(\mathcal{C}), \mathcal{O})$
과 임의의 미분 등급 $\mathcal{O}'$-대수 $\mathcal{A}'$ 및
미분 등급 $f^{-1}\mathcal{O}_\mathcal{D}$-대수의 임의의 사상
$\varphi : f^{-1}\mathcal{A} \to \mathcal{A}'$가 주어졌을 때,
당김 $f^*\mathcal{P}$에 대해서도 (1)과 (2)가 성립한다.
이를 미분 등급 $\mathcal{A}'$-가군으로 보고
(Section \ref{sdga-section-functoriality-dg}) 해석한다.
\end{enumerate}
첫 번째 조건은 $\mathcal{P}$가 오른쪽 등급
$\mathcal{A}$-가군으로서 평탄하다는 뜻이고, 두 번째 조건은
Spaltenstein의 의미에서 $\mathcal{P}$가 K-평탄하다는 뜻이다
(Cohomology on Sites, Section
\ref{sites-cohomology-section-flat} 참조). 세 번째 조건은 임의의
기저변경 뒤에도 이것이 성립한다는 뜻이다.

\medskip\noindent
다소 놀랍게도 good 가군은 많이 존재한다.

\begin{lemma}
\label{sdga-lemma-supply-good}
$(\mathcal{C}, \mathcal{O})$를 환 달린 사이트라 하자.
$\mathcal{A}$를 $(\mathcal{C}, \mathcal{O})$ 위의 미분 등급 대수의
층이라 하자. $U \in \Ob(\mathcal{C})$라 하자.
그러면 $j_!\mathcal{A}_U$는 좋은 미분 등급
$\mathcal{A}$-가군이다.
\end{lemma}

\begin{proof}
$\mathcal{N}$을 왼쪽 등급 $\mathcal{A}$-가군이라 하자.
Lemma \ref{sdga-lemma-tensor-with-extension-by-zero}에 의해
다음이 성립한다.
$$
j_!\mathcal{A}_U \otimes_\mathcal{A} \mathcal{N} =
j_!(\mathcal{A}_U \otimes_{\mathcal{A}_U} \mathcal{N}|_U) =
j_!(\mathcal{N}_U)
$$
등급 가군으로서 성립한다.
$U$로의 제한과 $j_!$가 모두 완전이므로 이는 조건 (1)을 증명한다.
Lemma \ref{sdga-lemma-tensor-with-extension-by-zero-dg}를 사용하면
같은 논법으로 조건 (2)도 얻는다.

\medskip\noindent
환 달린 토포스의 사상
$(f, f^\sharp) : (\Sh(\mathcal{C}'), \mathcal{O}')
\to (\Sh(\mathcal{C}), \mathcal{O})$,
미분 등급 $\mathcal{O}'$-대수 $\mathcal{A}'$,
그리고 미분 등급 $f^{-1}\mathcal{O}$-대수의 사상
$\varphi : f^{-1}\mathcal{A} \to \mathcal{A}'$를 생각하자.
다음을 보여야 한다.
$$
f^*j_!\mathcal{A}_U =
f^{-1}j_!\mathcal{A}_U \otimes_{f^{-1}\mathcal{A}} \mathcal{A}'
$$
이 식이 미분 등급 $\mathcal{O}'$-대수의 층
$\mathcal{A}'$를 갖춘 환 달린 토포스
$(\Sh(\mathcal{C}'), \mathcal{O}')$에서 (1)과 (2)를
만족함을 보여야 한다. 이를 보이기 위해 다음을
동치인 환 달린 토포스로 바꾸어도 된다.
$(\Sh(\mathcal{C}), \mathcal{O})$와
$(\Sh(\mathcal{C}'), \mathcal{O}')$를 동치인 환 달린 토포스로
바꾼다. 따라서
Modules on Sites, Lemma
\ref{sites-modules-lemma-morphism-ringed-topoi-comes-from-morphism-ringed-sites}
$f$가 연속 함자 $u : \mathcal{C} \to \mathcal{C}'$에 의해 주어지는
사이트의 사상
$f : \mathcal{C} \to \mathcal{C}'$에서 온다고 가정할 수 있다.
이 경우 $U' = u(U)$로 놓고, 이에 대응하는 국소화 사상
$j' : \Sh(\mathcal{C}'/U') \to \Sh(\mathcal{C}')$를
나타내자. 그러면 환 달린 토포스 사상들의 다음 가환 정사각형을 얻는다.
$$
\xymatrix{
(\Sh(\mathcal{C}'/U'), \mathcal{O}'_{U'})
\ar[rr]_{(j', (j')^\sharp)} \ar[d]_{(f', (f')^\sharp)} & &
(\Sh(\mathcal{C}'), \mathcal{O}')
\ar[d]^{(f, f^\sharp)} \\
(\Sh(\mathcal{C}/U), \mathcal{O}_U)
\ar[rr]^{(j, j^\sharp)} & &
(\Sh(\mathcal{C}), \mathcal{O}).
}
$$
그리고 $f'_*(j')^{-1} = j^{-1}f_*$가 성립한다. 이는
Modules on Sites, Lemma
\ref{sites-modules-lemma-localize-morphism-ringed-sites}.
수반의 유일성에 의해
$f^{-1}j_! = j'_!(f')^{-1}$을 얻는다. 따라서 다음을 얻는다.
\begin{align*}
f^*j_!\mathcal{A}_U
& =
f^{-1}j_!\mathcal{A}_U \otimes_{f^{-1}\mathcal{A}} \mathcal{A}' \\
& =
j'_!(f')^{-1}\mathcal{A}_U \otimes_{f^{-1}\mathcal{A}} \mathcal{A}' \\
& =
j'_!\left(
(f')^{-1}\mathcal{A}_U \otimes_{f^{-1}\mathcal{A}|_{U'}}
\mathcal{A}'|_{U'}\right) \\
& =
j'_!\mathcal{A}'_{U'}
\end{align*}
첫 번째 식은 $j_!\mathcal{A}_U$를 $\mathcal{A}'$ 위의
미분 등급 가군으로 당기는 것의 정의이다.
두 번째 식은 $f^{-1}j_! = j'_!(f')^{-1}$이기 때문이다.
세 번째 식은 환 달린 사이트 $(\mathcal{C}', f^{-1}\mathcal{O})$에
미분 등급 대수의 층 $f^{-1}\mathcal{A}$를 주고,
$\mathcal{C}'/U'$ 위의 미분 등급 가군 $(f')^{-1}\mathcal{A}_U$와
$\mathcal{C}'$ 위의 $\mathcal{A}'$에 적용한
Lemma \ref{sdga-lemma-tensor-with-extension-by-zero-dg}에 의한 것이다.
네 번째 식은 물론
$(f')^{-1}\mathcal{A}_U = f^{-1}\mathcal{A}|_{U'}$이기 때문이다.
따라서 당김도 같은 종류의 가군이며, 위에서 그 가군에 대해
조건 (1)과 (2)를 증명했다.
\end{proof}

\begin{lemma}
\label{sdga-lemma-good-admissible-ses}
$(\mathcal{C}, \mathcal{O})$를 환 달린 사이트라 하자.
$\mathcal{A}$를 $(\mathcal{C}, \mathcal{O})$ 위의 미분 등급 대수의
층이라 하자. 다음을 생각하자.
$0 \to \mathcal{P} \to \mathcal{P}' \to \mathcal{P}'' \to 0$
이것이 미분 등급 $\mathcal{A}$-가군의 허용 짧은 완전열이라고 하자.
이 가군들 중 셋 가운데 둘이 good이면 나머지 하나도 good이다.
\end{lemma}

\begin{proof}
조건 (1)은 이 열이 등급 수준에서 직합이므로 즉시 성립한다.
조건 (2)를 위해 임의의 왼쪽 미분 등급 $\mathcal{A}$-가군에 대해
다음 열을 생각하자.
$$
0 \to 
\mathcal{P} \otimes_\mathcal{A} \mathcal{N} \to
\mathcal{P}' \otimes_\mathcal{A} \mathcal{N} \to
\mathcal{P}'' \otimes_\mathcal{A} \mathcal{N} \to 0
$$
이는 미분을 잊으면 텐서곱이 등급 가군의 범주에서 취해진다는 점에 의해
미분 등급 $\mathcal{O}$-가군의 허용 짧은 완전열이다.
따라서 셋 가운데 둘이 $\mathcal{O}$-가군의 복합체로서 완전하면
나머지 하나도 그러하다. 마지막으로 같은 논법에 의해
환 달린 토포스의 사상 $(f, f^\sharp) : (\Sh(\mathcal{C}'), \mathcal{O}')
\to (\Sh(\mathcal{C}), \mathcal{O})$
과 미분 등급 $\mathcal{O}'$-대수
$\mathcal{A}'$, 그리고 미분 등급 $f^{-1}\mathcal{O}$-대수의 사상
$\varphi : f^{-1}\mathcal{A} \to \mathcal{A}'$가 주어지면
다음이 성립한다.
$$
0 \to f^*\mathcal{P} \to f^*\mathcal{P}' \to f^*\mathcal{P}'' \to 0
$$
이는 미분 등급 $\mathcal{A}'$-가군의 허용 짧은 완전열이며,
위와 같은 논법을 적용할 수 있다.
% P03-STACKS-E1054: frozen source line 2767 begins "et" instead of "Let".
\end{proof}

\begin{lemma}
\label{sdga-lemma-good-direct-sum}
$(\mathcal{C}, \mathcal{O})$를 환 달린 사이트라 하자.
$\mathcal{A}$를 $(\mathcal{C}, \mathcal{O})$ 위의 미분 등급 대수의
층이라 하자. 좋은 미분 등급 $\mathcal{A}$-가군의 임의의
직합은 good이다. 좋은 미분 등급 $\mathcal{A}$-가군의
여과 여극한도 good이다.
\end{lemma}

\begin{proof}
생략한다. 힌트: 직합과 여과 여극한은 텐서곱 및 당김과
가환한다.
\end{proof}

\begin{lemma}
\label{sdga-lemma-good-quotient}
$(\mathcal{C}, \mathcal{O})$를 환 달린 사이트라 하자.
$\mathcal{A}$를 $(\mathcal{C}, \mathcal{O})$ 위의 미분 등급 대수의
층이라 하자. $\mathcal{M}$을 미분 등급 $\mathcal{A}$-가군이라 하자.
다음 성질을 갖는 미분 등급 $\mathcal{A}$-가군의 준동형
$\mathcal{P} \to \mathcal{M}$이 존재한다.
\begin{enumerate}
\item $\mathcal{P} \to \mathcal{M}$은 전사이다.
\item $\Ker(\text{d}_\mathcal{P}) \to \Ker(\text{d}_\mathcal{M})$
은 전사이고,
\item $\mathcal{P}$는 good이다.
\end{enumerate}
\end{lemma}

\begin{proof}
$(U,k,x)$인 세 쌍을 생각하자. 여기서 $U$는 $\mathcal{C}$의 대상이고,
$k \in \mathbf{Z}$이며, $x$는 $\text{d}_\mathcal{M}(x)=0$을 만족하는
$U$ 위의 $\mathcal{M}^k$의 단면이다. 그러면 $1$을 $x$로 보내는
미분 등급 $\mathcal{A}_U$-가군의 유일한 사상
$\varphi_x : \mathcal{A}_U[-k] \to \mathcal{M}|_U$를 얻는다.
이는 다음 사상에 수반된다.
$\psi_x : j_{U!}\mathcal{A}_U[-k] \to \mathcal{M}$.
$1 \in \mathcal{A}_U(U)$는 미분이 0이고 $\psi_x$에 의해 $x$로
보내지는, 차수 $k$의 단면
$1 \in j_{U!}\mathcal{A}_U[-k](U)$에 대응함에 주의하자.
따라서 다음 사상을 생각하면
$$
\bigoplus\nolimits_{(U, k, x)} j_{U!}\mathcal{A}_U[-k]
\longrightarrow
\mathcal{M}
$$
조건 (2)와 (3)을 얻는다. 즉 대상
$j_{U!}\mathcal{A}_U[-k]$는 good이다
(Lemma \ref{sdga-lemma-supply-good}). 또한 good 대상의 직합은 good이다
(Lemma \ref{sdga-lemma-good-direct-sum}).

\medskip\noindent
다음으로 $U$가 $\mathcal{C}$의 대상이고 $k\in\mathbf{Z}$이며
$x$가 $\mathcal{M}^k$의 단면인 세 쌍 $(U,k,x)$를 생각하자
(미분에 의해 소멸할 필요는 없다). Remark
\ref{sdga-remark-cone-identity}와 같이 항등 사상
$\mathcal{A}_U \to \mathcal{A}_U$의 원뿔 $C_U$를 생각할 수 있다.
원소 $x$는 다음 사상을 정한다.
$\varphi_x : C_U[-k - 1] \to \mathcal{A}_U$.
% P03-STACKS-E1055: frozen source line 2862 maps x into A_U although x is a section of M^k.
Remark \ref{sdga-remark-cone-identity}를 보라. 이제 다음 허용 짧은
완전열이 있으므로
$$
0 \to \mathcal{A}_U \to C_U \to \mathcal{A}_U[1] \to 0
$$
Lemma \ref{sdga-lemma-good-admissible-ses}와 이미 사용한 Lemma
\ref{sdga-lemma-supply-good}에 의해 $j_{U!}C_U$가 good 가군임을 얻는다.
앞과 같이 $\varphi_x$에 수반되는 사상
$\psi_x : j_{U!}C_U \to \mathcal{M}$들의 직합
$$
\bigoplus\nolimits_{(U, k, x)} j_{U!}C_U \longrightarrow \mathcal{M}
$$
은 전사이다. 첫 번째 문단에서 얻은 사상과 직합하면 결론을 얻는다.
\end{proof}

\begin{remark}
\label{sdga-remark-sheaf-graded-sets}
$(\mathcal{C}, \mathcal{O})$를 환 달린 사이트라 하자.
$\mathcal{C}$ 위의 {\it 등급 집합의 층}이란
다음 사상을 갖춘 집합의 층 $\mathcal{S}$를 말한다.
$\deg : \mathcal{S} \to \underline{\mathbf{Z}}$
집합의 층의 사상이다. 등급 집합의 층 $\mathcal{S}$ 위의
자유 $\mathcal{O}$-가군인 등급 $\mathcal{O}$-가군을
$\mathcal{O}[\mathcal{S}]$로 나타내자. 더 정확히 말해
$\mathcal{O}[\mathcal{S}]$의 $n$번째 등급 부분은 다음 규칙의
층화이다.
$$
U \longmapsto
\bigoplus\nolimits_{s \in \mathcal{S}(U),\ \deg(s) = n} s \cdot \mathcal{O}(U)
$$
미분을 0으로 두면 이를 미분 등급 $\mathcal{O}$-가군으로도
생각할 수 있다. $\mathcal{A}$를 등급 $\mathcal{O}$-대수의 층이라
하자. 마찬가지로 $\mathcal{A}[\mathcal{S}]$를 등급
$\mathcal{A}$-가군으로 정의하자. 그 $n$번째 등급 부분은 다음
규칙의 층화이다.
$$
U \longmapsto
\bigoplus\nolimits_{s \in \mathcal{S}(U)} s \cdot \mathcal{A}^{n - \deg(s)}(U)
$$
$\mathcal{A}$가 미분 등급 $\mathcal{O}$-대수의 층이면,
모든 $s \in \mathcal{S}(U)$에 대해 $\text{d}(s)=0$으로 두고
층화하여 이를 미분 등급 $\mathcal{O}$-가군으로 만든다.
\end{remark}

\begin{lemma}
\label{sdga-lemma-free-graded-module-good}
$(\mathcal{C}, \mathcal{O})$를 환 달린 사이트라 하자.
$\mathcal{A}$를 미분 등급 $\mathcal{A}$-대수라 하자.
$\mathcal{S}$를 $\mathcal{C}$ 위의 등급 집합의 층이라 하자.
그러면 Remark \ref{sdga-remark-sheaf-graded-sets}에서와 같이 미분을
부여한 $\mathcal{S}$ 위의 자유 등급 가군 $\mathcal{A}[\mathcal{S}]$는
좋은 미분 등급 $\mathcal{A}$-가군이다.
\end{lemma}

\begin{proof}
$\mathcal{N}$을 왼쪽 등급 $\mathcal{A}$-가군이라 하자.
그러면
$$
\mathcal{A}[\mathcal{S}] \otimes_\mathcal{A} \mathcal{N} =
\mathcal{O}[\mathcal{S}] \otimes_\mathcal{O} \mathcal{N} =
\mathcal{N}[\mathcal{S}]
$$
여기서 $\mathcal{N}[\mathcal{S}$는 차수 $n$ 부분이 다음
준층에 대응하는 층인 등급 $\mathcal{O}$-가군이다.
% P03-STACKS-E1056: frozen source line 2930 omits the closing bracket in \mathcal{N}[\mathcal{S}.
$$
U \longmapsto
\bigoplus\nolimits_{s \in \mathcal{S}(U)} s \cdot \mathcal{N}^{n - \deg(s)}(U)
$$
$\mathcal{N} \to \mathcal{N}[\mathcal{S}]$가 완전 함자임은
명백하므로 $\mathcal{A}[\mathcal{S}$는 등급
$\mathcal{A}$-가군으로서 평탄하다. 이제 $\mathcal{N}$이
미분 등급 왼쪽 $\mathcal{A}$-가군이라고 하자. 그러면
% P03-STACKS-E1057: frozen source line 2938 omits the closing bracket in \mathcal{A}[\mathcal{S}.
$$
H^*(\mathcal{A}[\mathcal{S}] \otimes_\mathcal{A} \mathcal{N}) =
H^*(\mathcal{O}[\mathcal{S}] \otimes_\mathcal{O} \mathcal{N})
$$
등급 $\mathcal{O}$-가군의 층으로서 성립하며, 이는
($\mathcal{O})$ 위의 평탄성에 의해 다음과 같다.
% P03-STACKS-E1058: frozen source line 2946 has an unmatched parenthesis in "(over \mathcal{O})".
$$
H^*(\mathcal{N})[\mathcal{S}]
$$
등급 $\mathcal{O}$-가군으로서 성립한다. 따라서 $\mathcal{N}$이
비순환이면 $\mathcal{A}[\mathcal{S}] \otimes_\mathcal{A} \mathcal{N}$도
비순환이다.

\medskip\noindent
마지막으로 다음 사상을 생각하자.
$(f, f^\sharp) : (\Sh(\mathcal{C}'), \mathcal{O}')
\to (\Sh(\mathcal{C}), \mathcal{O})$
환 달린 토포스의 사상, 미분 등급 $\mathcal{O}'$-대수
$\mathcal{A}'$, 그리고 미분 등급 $f^{-1}\mathcal{O}$-대수의 사상
$\varphi : f^{-1}\mathcal{A} \to \mathcal{A}'$에 대해
다음이 자명하게 성립한다.
$$
f^*\mathcal{A}[\mathcal{S}] = \mathcal{A}'[f^{-1}\mathcal{S}]
$$
이로써 이 가군이 good임이 증명된다.
\end{proof}

\begin{lemma}
\label{sdga-lemma-resolve}
$(\mathcal{C}, \mathcal{O})$를 환 달린 사이트라 하자.
$\mathcal{A}$를 $(\mathcal{C}, \mathcal{O})$ 위의 미분 등급 대수의
층이라 하자. $\mathcal{M}$을 미분 등급 $\mathcal{A}$-가군이라 하자.
다음 성질을 갖는 미분 등급 $\mathcal{A}$-가군의 준동형
$\mathcal{P} \to \mathcal{M}$이 존재한다.
\begin{enumerate}
\item $\mathcal{P} \to \mathcal{M}$은 유사동형이다.
\item $\mathcal{P}$는 good이다.
\end{enumerate}
\end{lemma}

\begin{proof}[첫 번째 증명]
$\mathcal{S}_0$를 (Remark \ref{sdga-remark-sheaf-graded-sets})에서와 같이
차수 $n$ 부분이 $\Ker(\text{d}_\mathcal{M}^n)$인
등급 집합의 층이라 하자. 다음 미분 등급 가군의 준동형을 생각하자.
$$
\mathcal{P}_0 = \mathcal{A}[\mathcal{S}_0] \longrightarrow \mathcal{M}
$$
좌변은 Remark \ref{sdga-remark-sheaf-graded-sets}에서와 같고,
이 사상은 $\mathcal{S}_0$의 국소 단면 $s$를 이에 대응하는
$\mathcal{M}^{\deg(s)}$의 국소 단면으로 보낸다.
이는 미분의 핵에 있으므로 실제로 미분 등급 가군의 사상이다.
구성에 의해 코호몰로지 층 사이의 사상
$H^n(\mathcal{P}_0) \to H^n(\mathcal{M})$은 전사이다.
다음 사상들을 귀납적으로 구성하자.
$$
\mathcal{P}_0 \to \mathcal{P}_1 \to \mathcal{P}_2 \to \ldots \to \mathcal{M}
$$
분명히 모든 $i$에 대해 $H^*(\mathcal{P}_i) \to H^*(\mathcal{M})$은
전사이다. $\mathcal{P}_i \to \mathcal{M}$이 주어졌다고 하자.
차수 $n$ 부분이 다음인 등급 집합의 층을 $\mathcal{S}_{i+1}$로
나타내자.
$$
\Ker(\text{d}_{\mathcal{P}_i}^{n + 1})
\times_{\mathcal{M}^{n + 1}, \text{d}}
\mathcal{M}^n
$$
그리고
$$
\mathcal{P}_{i + 1} = \mathcal{P}_i \oplus \mathcal{A}[\mathcal{S}_{i + 1}]
$$
를 미분과 $\mathcal{M}$으로의 사상을 다음과 같이 정의한
등급 $\mathcal{A}$-가군으로 둔다.
\begin{enumerate}
\item $\mathcal{P}_i$의 국소 단면에는 $\mathcal{P}_i$의 미분과
$\mathcal{M}$으로의 주어진 사상을 사용한다.
\item $\mathcal{S}_{i+1}$의 국소 단면 $s=(p,m)$에 대해
$\text{d}(s)$는 차수 $\deg(s)+1$인 $\mathcal{P}_i$의 단면으로 본
$p$로 정하고, $s$를 $\mathcal{M}$의 $m$으로 보낸다.
\item 라이프니츠 법칙이 성립하도록 미분을 유일하게 확장한다.
\end{enumerate}
이는 $\text{d}(m)$이 $p$의 상이고 $\text{d}(p)=0$이므로
타당하다. 마지막으로 유도된 $\mathcal{M}$으로의 사상과 함께
$\mathcal{P}=\colim \mathcal{P}_i$로 둔다.

\medskip\noindent
사상 $\mathcal{P} \to \mathcal{M}$은 유사동형이다.
$H^n(\mathcal{P})=\colim H^n(\mathcal{P}_i)$이고, 각 $i$에 대해
$H^n(\mathcal{P}_i) \to H^n(\mathcal{M})$은 전사이며 그 핵은
구성에 의해 $H^n(\mathcal{P}_i) \to H^n(\mathcal{P}_{i+1})$에 의해
소멸한다. 각 $\mathcal{P}_i$는 good이다. 왜냐하면
$\mathcal{P}_0$는 Lemma \ref{sdga-lemma-free-graded-module-good}에 의해
good이고, 각 $\mathcal{P}_{i+1}$은 다음 허용 짧은 완전열의 가운데 항이기
때문이다.
$0 \to \mathcal{P}_i \to \mathcal{P}_{i + 1} \to
\mathcal{A}[\mathcal{S}_{i + 1}] \to 0$
양 끝 항은 귀납 가정에 의해 good이다. 따라서
$\mathcal{P}_{i+1}$은 Lemma \ref{sdga-lemma-good-admissible-ses}에 의해
good이다. 마지막으로 Lemma \ref{sdga-lemma-good-direct-sum}에 의해
$\mathcal{P}$도 good임을 얻는다.
\end{proof}

\begin{proof}[두 번째 증명]
이 증명을 읽기 전에 Differential Graded Algebra,
Lemma \ref{dga-lemma-resolve}의 증명을 읽기를 권한다.
$\mathcal{M}=\mathcal{M}_0$로 둔다.
다음 짧은 완전열을 귀납적으로 택한다.
$$
0 \to \mathcal{M}_{i + 1} \to \mathcal{P}_i \to \mathcal{M}_i \to 0
$$
여기서 $\mathcal{P}_i \to \mathcal{M}_i$는 Lemma
\ref{sdga-lemma-good-quotient}에서와 같이 택한다.
그러면 다음 ‘해소’를 얻는다.
$$
\ldots \to \mathcal{P}_2 \xrightarrow{f_2}
\mathcal{P}_1 \xrightarrow{f_1} \mathcal{P}_0 \to \mathcal{M} \to 0
$$
이제 $\mathcal{P}$를 다음과 같이 정의한 미분 등급
$\mathcal{A}$-가군으로 둔다.
\begin{enumerate}
\item 등급 $\mathcal{A}$-가군으로서 다음과 같이 둔다.
$\mathcal{P} = \bigoplus_{a \leq 0} \mathcal{P}_{-a}[-a]$, 즉
차수 $n$ 부분은 다음과 같다.
$\mathcal{P}^n = \bigoplus\nolimits_{a + b = n} \mathcal{P}_{-a}^b$,
\item $\mathcal{P}$의 미분은 다음 식으로 주어지는 이중 복합체에
대응하는 전체 복합체의 구성에서와 같다.
$$
\text{d}_\mathcal{P}(x) = f_{-a}(x) + (-1)^a \text{d}_{\mathcal{P}_{-a}}(x)
$$
여기서 $x$는 $\mathcal{P}_{-a}^b$의 국소 단면이다.
\end{enumerate}
이러한 약속 아래 $\mathcal{P}$는 실제로 미분 등급
$\mathcal{A}$-가군이다. 세부사항은 생략한다.
$a<0$인 합성분 $\mathcal{P}_{-a}[-a]$에서는 0이고
$a=0$인 합성분에서는 주어진 $\mathcal{P}_0\to\mathcal{M}$인
미분 등급 $\mathcal{A}$-가군 사상 $\mathcal{P}\to\mathcal{M}$이 있다.
다음을 주목하자.
$$
\mathcal{P} = \colim_i F_i\mathcal{P}
$$
여기서 $F_i\mathcal{P}\subset\mathcal{P}$는 바탕 등급
$\mathcal{A}$-가군이 다음인 미분 등급 $\mathcal{A}$-부분가군이다.
$$
F_i\mathcal{P} =
\bigoplus\nolimits_{i \geq -a \geq 0} \mathcal{P}_{-a}[-a]
$$
다음 사상들이
$$
0 \to F_1\mathcal{P} \to F_2\mathcal{P} \to F_3\mathcal{P} \to
\ldots \to \mathcal{P}
$$
모두 허용 단사사상이고 다음 허용 짧은 완전열을 갖는다는 것은 자명하다.
$$
0 \to F_i\mathcal{P} \to
F_{i + 1}\mathcal{P} \to
\mathcal{P}_{i + 1}[i + 1] \to 0
$$
귀납과 Lemma \ref{sdga-lemma-good-admissible-ses}에 의해
$F_i\mathcal{P}$가 good 미분 등급 $\mathcal{A}$-가군임을 얻는다.
$\mathcal{P}=\colim F_i\mathcal{P}$이므로 Lemma
\ref{sdga-lemma-good-direct-sum}에 의해 $\mathcal{P}$는 good이다.

\medskip\noindent
마지막으로 $\mathcal{P}\to\mathcal{M}$이 유사동형임을 보여야 한다.
$\mathcal{C}$가 충분한 점을 가지면, 이는 stalk에서 확인하여
초등적인 Homology, Lemma
\ref{homology-lemma-good-resolution-gives-qis}에서 따른다.
일반적으로는 다음과 같이 논할 수 있다
(이 증명은 매우 길다 — 초등적 증명처럼 국소 단면으로 작업하는
다른 논법도 있지만 역시 상당히 길다).
아벨 층의 범주에서 여과 여극한은 완전하므로 다음을 얻는다.
$$
H^d(\mathcal{P}) = \colim H^d(F_i\mathcal{P})
$$
각 $i\geq0$와 $d\in\mathbf{Z}$에 대해 (a) 다음 짧은 완전열과
$$
0 \to H^d(\mathcal{M}_{i + 1}[i]) \to
H^d(F_i\mathcal{P}) \to H^d(\mathcal{M}) \to 0
$$
두 번째 사상은 $F_i\mathcal{P}\to\mathcal{P}\to\mathcal{M}$에서
오는 것이며, (b) 다음 합성이
$$
H^d(\mathcal{M}_{i + 1}[i]) \to
H^d(F_i\mathcal{P}) \to
H^d(F_{i + 1}\mathcal{P})
$$
0이다. 이 주장이 증명을 끝내기에 충분함은 자명하다.

\medskip\noindent
주장의 증명. 모든 $i\geq0$에 대해
$\mathcal{M}_{i+1}[i]\to F_i\mathcal{P}$인 사상이 있다.
이는 $\mathcal{M}_{i+1}$을 $f_i$의 핵으로서 $\mathcal{P}_i$에
포함시키는 것에서 온다. $C_i$를 정의하는 $\mathcal{O}$-가군
복합체의 다음 짧은 완전열을 생각하자.
$$
0 \to \mathcal{M}_{i + 1}[i] \to F_i\mathcal{P} \to C_i \to 0
$$
$C_0=\mathcal{M}_0=\mathcal{M}$임을 주목하자.
또한 $C_i$는 다음 열을 갖는 이중 복합체
$C_i^{\bullet,\bullet}$에 대응하는 전체 복합체이다.
$$
\mathcal{M}_i = \mathcal{P}_i/\mathcal{M}_{i + 1},
\mathcal{P}_{i - 1}, \ldots, \mathcal{P}_0
$$
차수 $-i,-i+1,\ldots,0$에 놓인다. 이중 복합체의 사상
$C_i^{\bullet,\bullet}\to C_{i-1}^{\bullet,\bullet}$이 있는데,
차수 $-i$의 열에서는 $0$이고, 차수 $-i+1$에서는 전사
$\mathcal{P}_{i-1}\to\mathcal{M}_{i-1}$이며, 다른 열에서는 항등이다.
따라서 복합체의 사상
$$
C_i \longrightarrow C_{i - 1}
$$
이 사상들은 전사 유사동형이다. 핵이 차수 $-i,-i+1$에
$\mathcal{M}_i,\mathcal{M}_i$를 열로 갖고 이 두 열 사이의 사상이
항등인 이중 복합체의 전체 복합체이기 때문이다.
그 결과 얻는 다음 식별을 사용하면
$H^d(C_i) = H^d(C_{i - 1} = \ldots = H^d(\mathcal{M})$
이미 다음 긴 완전열을 얻는다.
% P03-STACKS-E1059: frozen source line 3168 has an unmatched parenthesis in the chain of H^d(C_i) identifications.
$$
H^d(\mathcal{M}_{i + 1}[i]) \to
H^d(F_i\mathcal{P}) \to H^d(\mathcal{M}) \to
H^{d + 1}(\mathcal{M}_{i + 1}[i])
$$
이는 위의 복합체의 짧은 완전열에서 따른다.
그러나 다음 가환 도표도 갖는다.
$$
\xymatrix{
\mathcal{M}_{i + 2}[i + 1] \ar[r]_a & T_{i + 1} \ar[r] &
F_{i + 1}\mathcal{P} \ar[r] & C_{i + 1} \ar[d] \\
& \mathcal{M}_{i + 1}[i] \ar[r] \ar[u]^b &
F_i\mathcal{P} \ar[u] \ar[r] &
C_i
}
$$
$T_{i+1}$은 차수 $-i-1$과 $-i$에 각각
$\mathcal{P}_{i+1},\mathcal{M}_{i+1}$을 열로 놓은 이중 복합체의
전체 복합체이다. 즉 $T_{i+1}$은 사상
$\mathcal{P}_{i+1}\to\mathcal{M}_{i+1}$의 원뿔의 이동이며,
$a$는 유사동형이고 $a^{-1}\circ b$는 다음 짧은 완전열에 대응하는
$D(\mathcal{O})$의 특별 삼각형에서 세 번째 사상의 이동임을 얻는다.
$$
0 \to \mathcal{M}_{i + 2} \to \mathcal{P}_{i + 1} \to \mathcal{M}_{i + 1} \to 0
$$
사상 $H^d(\mathcal{P}_{i+1})\to H^d(\mathcal{M}_{i+1})$은
다음 사상이 전사하도록 사상들을 택했으므로 전사이다.
$\Ker(\text{d}_{\mathcal{P}_{i + 1}}) \to
\Ker(\text{d}_{\mathcal{M}_{i + 1}})$는 전사이다.
따라서 $a^{-1}\circ b$는 코호몰로지 층에서 0이다.
이로써 주장의 (b)가 증명된다.
$T_{i+1}$은 전사 복합체 사상 $F_{i+1}\mathcal{P}\to C_i$의
핵이므로 긴 완전 코호몰로지 열의 사상을 얻는다.
$$
\xymatrix{
H^d(T_{i + 1}) \ar[r] &
H^d(F_{i + 1}\mathcal{P}) \ar[r] &
H^d(\mathcal{M}) \ar[r] &
H^{d + 1}(T_{i + 1}) \\
H^d(\mathcal{M}_{i + 1}[i]) \ar[r] \ar[u] &
H^d(F_i\mathcal{P}) \ar[r] \ar[u] &
H^d(\mathcal{M}) \ar[r] \ar[u] &
H^{d + 1}(\mathcal{M}_{i + 1}[i]) \ar[u]
}
$$
위의 논의에 의해 양 끝의 세로 사상은 0이다. 따라서
$H^d(F_{i+1}\mathcal{P})\to H^d(\mathcal{M})$은 전사이고
주장의 (a)가 따른다. 더 정확히는 $i>0$일 때 주장이 성립하며,
$i=0$인 경우는 독자에게 맡긴다(실제로 보조정리를 얻기 위해서는
모든 $i\gg0$에 대해 주장만 증명하면 충분하다).
\end{proof}

\begin{lemma}
\label{sdga-lemma-acyclic-good}
$(\mathcal{C}, \mathcal{O})$를 환 달린 사이트라 하자.
$\mathcal{A}$를 $(\mathcal{C}, \mathcal{O})$ 위의 미분 등급 대수의
층이라 하자. $\mathcal{P}$를 good 비순환 오른쪽 미분 등급
$\mathcal{A}$-가군이라 하자.
\begin{enumerate}
\item 임의의 미분 등급 왼쪽 $\mathcal{A}$-가군
$\mathcal{N}$에 대해 텐서곱
$\mathcal{P} \otimes_\mathcal{A} \mathcal{N}$은 비순환이다.
\item 임의의 사상 $(f, f^\sharp) : (\Sh(\mathcal{C}'), \mathcal{O}')
\to (\Sh(\mathcal{C}), \mathcal{O})$
인 환 달린 토포스의 사상, 미분 등급 $\mathcal{O}'$-대수
$\mathcal{A}'$, 미분 등급 $f^{-1}\mathcal{O}$-대수의 사상
$\varphi : f^{-1}\mathcal{A} \to \mathcal{A}'$에 대해
당김 $f^*\mathcal{P}$는 비순환이고 good이다.
\end{enumerate}
\end{lemma}

\begin{proof}
 (1)의 증명. Lemma \ref{sdga-lemma-resolve}에 의해 good
왼쪽 미분 등급 $\mathcal{Q}$와 유사동형
$\mathcal{Q}\to\mathcal{N}$을 택할 수 있다. 그러면
$\mathcal{P} \otimes_\mathcal{A} \mathcal{Q}$
는 $\mathcal{Q}$가 good이므로 비순환이다.
$\mathcal{N}'$을 사상 $\mathcal{Q}\to\mathcal{N}$의 원뿔이라 하자.
그러면 $\mathcal{P}\otimes_\mathcal{A}\mathcal{N}'$는
$\mathcal{P}$가 good이고 $\mathcal{N}'$가 비순환
(유사동형의 원뿔)이므로 비순환이다. 다음 특별 삼각형을 갖는다.
$$
\mathcal{Q} \to \mathcal{N} \to \mathcal{N}' \to \mathcal{Q}[1]
$$
이는 삼각 구조의 구성에 의해
$K(\textit{Mod}(\mathcal{A}, \text{d}))$ 안에 있다.
$\mathcal{P}\otimes_\mathcal{A}-$가 특별 삼각형을 특별 삼각형으로
보내므로 다음 특별 삼각형을 얻는다.
$$
\mathcal{P} \otimes_\mathcal{A} \mathcal{Q} \to
\mathcal{P} \otimes_\mathcal{A} \mathcal{N} \to
\mathcal{P} \otimes_\mathcal{A} \mathcal{N}' \to
\mathcal{P} \otimes_\mathcal{A} \mathcal{Q}[1]
$$
이는 $K(\textit{Mod}(\mathcal{O}))$ 안에 있다. 따라서 결론을 얻는다.

\medskip\noindent
 (2)의 증명. good 가군의 정의에 의해 $f^*\mathcal{P}$가 good임을
주목하자. 다음을 상기하자.
$f^*\mathcal{P} = f^{-1}\mathcal{P} \otimes_{f^{-1}\mathcal{A}} \mathcal{A}'$.
그러면 $f^{-1}$이 완전이므로 $f^{-1}\mathcal{P}$는 good 비순환
미분 등급 $f^{-1}\mathcal{A}$-가군이다. 따라서 (1)에 의해
$f^*\mathcal{P}$가 비순환임을 얻는다.
\end{proof}






\section{등급 가군의 미분 등급 포락}
\label{sdga-section-dg-hull}

\noindent
등급 가군 $\mathcal{N}$의 미분 등급 포락은 다음 보조정리의
함자 $G$를 적용한 결과이다.

\begin{lemma}
\label{sdga-lemma-dg-hull}
$(\mathcal{C}, \mathcal{O})$를 환 달린 사이트라 하자.
$\mathcal{A}$를 $(\mathcal{C}, \mathcal{O})$ 위의 미분 등급 대수의
층이라 하자. 망각 함자
$F : \textit{Mod}(\mathcal{A}, \text{d}) \to \textit{Mod}(\mathcal{A})$
는 왼쪽 수반
$G : \textit{Mod}(\mathcal{A}) \to
\textit{Mod}(\mathcal{A}, \text{d})$.
\end{lemma}

\begin{proof}
$G$의 존재를 증명하기 위해 수반 함자 정리를 사용할 수 있다.
Categories, Theorem \ref{categories-theorem-adjoint-functor}를 보라
(여기서는 $F$와 $G$의 역할을 바꾸었다). $F$의 완전성 조건은
Lemma \ref{sdga-lemma-dgm-abelian}에 의해 성립한다. 집합론적 조건은
다음과 같이 볼 수 있다. 등급 $\mathcal{A}$-가군 $\mathcal{N}$이
주어졌다고 하자. 임의의 사상
$$
\varphi : \mathcal{N} \longrightarrow F(\mathcal{M})
$$
에 대해 $\Im(\varphi)\subset F(\mathcal{M}')$를 만족하는
가장 작은 미분 등급 $\mathcal{A}$-부분가군
$\mathcal{M}'\subset\mathcal{M}$을 생각할 수 있다.
$\mathcal{M}'$가 다음 등급 $\mathcal{A}$-가군 사상의 상임은 자명하다.
$$
\mathcal{N} \oplus
\mathcal{N}[-1] \otimes_\mathcal{O} \mathcal{A}
\longrightarrow
\mathcal{M}
$$
이 사상은 다음과 같이 정의된다.
$$
(n, \sum n_i \otimes a_i) \longmapsto
\varphi(n) + \sum \text{d}(\varphi(n_i)) a_i
$$
이 사상의 상은 미분 등급 $\mathcal{M}$-부분가군임을 쉽게 알 수
있기 때문이다. 따라서 이러한 $\mathcal{M}'$의 가능한 동형류의
수는 유계이고, 결론을 얻는다.
\end{proof}

\noindent
$(\mathcal{C}, \mathcal{O})$를 환 달린 사이트라 하자.
$(\mathcal{C}, \mathcal{O})$ 위에서 $\mathcal{A}$를 미분 등급 대수의
층이라 하자. $\mathcal{M}$을 미분 등급 $\mathcal{A}$-가군이라 하고
다음과 같은 짧은 완전열이 주어졌다고 하자.
$$
0 \to \mathcal{N} \to F(\mathcal{M}) \to \mathcal{N}' \to 0
$$
이 $\textit{Mod}(\mathcal{A})$ 안에 있다. 그러면 다음과 같은
표준 등급 $\mathcal{A}$-가군 준동형을 얻는다.
$$
\overline{\text{d}} : \mathcal{N} \to \mathcal{N}'[1]
$$
다음과 같이 정의한다. $\mathcal{N}$의 국소 단면 $x$가 주어지면,
% P03-STACKS-E1062: frozen source line 3350 has the grammatical fragment "given ... denote".
$x$를 $\mathcal{M}$의 국소 단면으로 보았을 때의
$\text{d}_\mathcal{M}(x)$를 $\mathcal{N}'$로 보내 얻는 상을
$\overline{\text{d}}(x)$라 하자.

\begin{lemma}
\label{sdga-lemma-dg-hull-acyclic}
Lemma \ref{sdga-lemma-dg-hull}의 함자 $F, G$는 다음 성질을 갖는다.
등급 $\mathcal{A}$-가군 $\mathcal{N}$이 주어지면
\begin{enumerate}
\item 코단위 $\mathcal{N} \to F(G(\mathcal{N}))$는 단사이고,
\item 사상 $\overline{\text{d}} : \mathcal{N} \to
\Coker(\mathcal{N} \to F(G(\mathcal{N})))[1]$는 동형이고,
\item $G(\mathcal{N})$은 비순환 미분 등급 $\mathcal{A}$-가군이다.
\end{enumerate}
\end{lemma}

\begin{proof}
성질 (3)은 성질 (1)과 (2)의 결과임을 관찰하자.
실제로 $\text{d}(s) = 0$인 $F(G(\mathcal{N}))$의 0이 아닌 국소
단면 $s$를 생각하면, $s$는 $\mathcal{N} \to F(G(\mathcal{N}))$의
상에 있을 수 없다. 따라서 cokernel에서 $s$의 상 $\overline{s}$를
어떤 $\mathcal{N}$의 국소 단면 $s'$에 대해
$\overline{\text{d}}(s')$로 쓸 수 있다. 그러면
$s = \text{d}(s')$임을 알 수 있다. $s - \text{d}(s')$가 여전히
$\text{d}$의 핵에 있고 코단위의 상에 포함되기 때문이다.

\medskip\noindent
잠시 $\mathcal{A}_{gr}$ 및 $\mathcal{A}_{dg}$를 각각 층
$\mathcal{A}$를 자기 자신 위의 (오른쪽) 등급 가군으로 본 것,
및 자기 자신 위의 (오른쪽) 미분 등급 가군으로 본 것으로 쓰자.
이 보조정리에서 가장 중요한 경우는 $G(\mathcal{A}_{gr})$가
무엇인지 이해하는 것이다. 물론 $G(\mathcal{A}_{gr})$는 다음 함자를
나타내는 $\textit{Mod}(\mathcal{A}, \text{d})$의 대상이다.
$$
\mathcal{M} \longmapsto
\Hom_{\textit{Mod}(\mathcal{A})}(\mathcal{A}_{gr}, F(\mathcal{M})) =
\Gamma(\mathcal{C}, \mathcal{M})
$$
비고 \ref{sdga-remark-cone-identity}에 의해 이 함자는
% P03-STACKS-E1060: frozen source line 3392 omits "is" in "this functor represented by".
$\mathcal{A}_{dg}$의 항등사상 위의 원뿔인 $C$에 대해 $C[-1]$이
나타냄을 알 수 있다.
다음과 같은 짧은 완전열이 있다.
$$
0 \to \mathcal{A}_{dg}[-1] \to C[-1] \to \mathcal{A}_{dg} \to 0
$$
이는 $\textit{Mod}(\mathcal{A}, \text{d})$ 안에서
$\textit{Mod}(\mathcal{A})$의 코단위
$\mathcal{A}_{gr} \to F(C[-1])$에 의해 분할된다.
따라서 $G(\mathcal{A}_{gr})$는 성질 (1)과 (2)를 만족한다.

\medskip\noindent
$U$를 $\mathcal{C}$의 대상이라 하자. 국소화 사상을
$j_U : \mathcal{C}/U \to \mathcal{C}$로 쓰자.
$\mathcal{A}_U$를 $U$ 위로 제한한 $\mathcal{A}$라 하자.
$\mathcal{A}_U$를 등급 $\mathcal{A}_U$-가군으로 본 것을
$\mathcal{A}_{U, gr}$로 나타내겠다.
망각 함자 $F_U : \textit{Mod}(\mathcal{A}_U, \text{d}) \to
\textit{Mod}(\mathcal{A}_U)$를 쓰고 그 수반을 $G_U$로 쓰자.
그러면 다음 가환 도식을 얻는다.
$$
\vcenter{
\xymatrix{
\textit{Mod}(\mathcal{A}, \text{d}) \ar[d]_{j_U^*} \ar[r]_F &
\textit{Mod}(\mathcal{A}) \ar[d]^{j_U^*} \\
\textit{Mod}(\mathcal{A}_U, \text{d}) \ar[r]^{F_U} &
\textit{Mod}(\mathcal{A}_U)
}
}
\quad\text{and}\quad
\vcenter{
\xymatrix{
\textit{Mod}(\mathcal{A}_U, \text{d}) \ar[r]_{F_U} \ar[d]_{j_{U!}} &
\textit{Mod}(\mathcal{A}_U) \ar[d]^{j_{U!}} \\
\textit{Mod}(\mathcal{A}, \text{d}) \ar[r]^F &
\textit{Mod}(\mathcal{A})
}
}
$$
이는 Sections \ref{sdga-section-functoriality-graded},
\ref{sdga-section-functoriality-dg},
\ref{sdga-section-localize-graded}, \ref{sdga-section-localize-dg}에서
% P03-STACKS-E1061: frozen source line 3430 writes j^*_U although the diagrams use j_U^*.
$j^*_U$와 $j_{U!}$를 구성한 것에 따른 것이다.
수반의 유일성에 의해 $j_{U!} \circ G_U = G \circ j_{U!}$를 얻는다.
$j_{U!}$가 완전 함자이므로, 앞 부분의 증명에서 본 코단위
$\mathcal{A}_{U, gr} \to F_U(G_U(\mathcal{A}_{U, gr}))$
에 대한 성질 (1), (2)가 다음 코단위에 대한 성질 (1), (2)를
함의함을 알 수 있다.
$j_{U!}\mathcal{A}_{U, gr} \to F(G(j_{U!}\mathcal{A}_{U, gr})) =
j_{U!}F_U(G_U(\mathcal{A}_{U, gr}))$.

\medskip\noindent
Lemma \ref{sdga-lemma-gm-grothendieck-abelian}의 증명에서
$\textit{Mod}(\mathcal{A})$의 임의의 대상은
$j_{U!}\mathcal{A}_{U, gr}$의 복사본들의 직합의 몫임을 보았다.
$G$는 왼쪽 수반이므로 $G$가 직합과 가환한다. 따라서 성질 (1), (2)는
그 성질을 만족하는 대상들의 직합에 대해서도 성립한다.
그러므로 $\textit{Mod}(\mathcal{A})$의 모든 대상 $\mathcal{N}$은
다음 완전열에 들어간다.
$$
\mathcal{N}_1 \to \mathcal{N}_0 \to \mathcal{N} \to 0
$$
여기서 $\mathcal{N}_1$과 $\mathcal{N}_0$에 대해 성질 (1), (2)가
성립한다. $G$가 오른쪽 완전임을 사용하여 $\mathcal{N}$에 대한
성질 (1), (2)를 도출하는 것은 독자에게 맡긴다.
\end{proof}









\section{K-단사 미분 등급 가군}
\label{sdga-section-K-injective}

\noindent
이 절은 미분 등급 대수의 층 위 미분 등급 가군의 층이라는
설정에서 Injectives, Section \ref{injectives-section-K-injective}의
유사물이다.

\begin{lemma}
\label{sdga-lemma-characterize-injectives}
$(\mathcal{C}, \mathcal{O})$를 환 달린 사이트라 하자. $\mathcal{A}$를
$(\mathcal{C}, \mathcal{O})$ 위의 등급 대수의 층이라 하자.
집합 $T$가 존재하여 각 $t \in T$에 대해 등급 $\mathcal{A}$-가군의
단사 사상 $\mathcal{N}_t \to \mathcal{N}'_t$가 주어진다고 하자.
그러면 $\textit{Mod}(\mathcal{A})$의 대상 $\mathcal{I}$가 단사인 것은
모든 다음 가환 도식에 대해
$$
\xymatrix{
\mathcal{N}_t \ar[r] \ar[d] & \mathcal{I} \\
\mathcal{N}'_t \ar@{..>}[ru]
}
$$
도식을 가환하게 만드는 점선 화살표가
$\textit{Mod}(\mathcal{A})$에 존재하는 것과 동치이다.
\end{lemma}

\begin{proof}
이는 임의의 그로텐디크 아벨 범주에서 성립한다. Injectives,
Lemma \ref{injectives-lemma-characterize-injective}를 보라.
Lemma \ref{sdga-lemma-gm-grothendieck-abelian}에 의해 범주
$\textit{Mod}(\mathcal{A})$는 그로텐디크 아벨 범주이다.
\end{proof}

\begin{definition}
\label{sdga-definition-graded-injective}
$(\mathcal{C}, \mathcal{O})$를 환 달린 사이트라 하자.
$(\mathcal{A}, \text{d})$를 $(\mathcal{C}, \mathcal{O})$ 위의
미분 등급 대수의 층이라 하자. 미분 등급 $\mathcal{A}$-가군
% P03-STACKS-E1063: frozen source spells "diffential" at lines 3506, 3518, and 3564; Korean renders the intended "differential".
$\mathcal{I}$를 {\it 등급 단사}\footnote{이는 표준적이지 않을 수 있는
용어이다.}라고 하는 것은 $\mathcal{M}$을 등급 $\mathcal{A}$-가군으로
% P03-STACKS-E1064: frozen source line 3508 names \mathcal{M} although the definition's subject is \mathcal{I}.
보았을 때 등급 $\mathcal{A}$-가군의 범주
$\textit{Mod}(\mathcal{A})$에서 단사 대상인 경우이다.
\end{definition}

\begin{remark}
\label{sdga-remark-why-graded-injective}
$(\mathcal{C}, \mathcal{O})$를 환 달린 사이트라 하자.
$(\mathcal{A}, \text{d})$를 $(\mathcal{C}, \mathcal{O})$ 위의
미분 등급 대수의 층이라 하자. $\mathcal{I}$를 등급 단사
미분 등급 $\mathcal{A}$-가군이라 하자.
$$
0 \to \mathcal{M}_1 \to \mathcal{M}_2 \to \mathcal{M}_3 \to 0
$$
가 미분 등급 $\mathcal{A}$-가군의 짧은 완전열이라고 하자.
$\mathcal{I}$가 등급 단사이므로 다음 복합체의 짧은 완전열을 얻는다.
$$
0 \to
\Hom_{\textit{Mod}^{dg}(\mathcal{A}, \text{d})}(\mathcal{M}_3, \mathcal{I})
\to
\Hom_{\textit{Mod}^{dg}(\mathcal{A}, \text{d})}(\mathcal{M}_2, \mathcal{I})
\to
\Hom_{\textit{Mod}^{dg}(\mathcal{A}, \text{d})}(\mathcal{M}_1, \mathcal{I})
\to 0
$$
이는 $\Gamma(\mathcal{C}, \mathcal{O})$-가군의 짧은 완전열이다.
코호몰로지를 취하면 다음 긴 완전열을 얻는다.
$$
\xymatrix{
\Hom_{K(\textit{Mod}(\mathcal{A}, \text{d}))}(\mathcal{M}_3, \mathcal{I})
\ar[d] &
\Hom_{K(\textit{Mod}(\mathcal{A}, \text{d}))}(\mathcal{M}_3, \mathcal{I})[1]
\ar[d] \\
\Hom_{K(\textit{Mod}(\mathcal{A}, \text{d}))}(\mathcal{M}_2, \mathcal{I})
\ar[d] &
\Hom_{K(\textit{Mod}(\mathcal{A}, \text{d}))}(\mathcal{M}_2, \mathcal{I})[1]
\ar[d] \\
\Hom_{K(\textit{Mod}(\mathcal{A}, \text{d}))}(\mathcal{M}_1, \mathcal{I})
\ar[ruu]
&
\Hom_{K(\textit{Mod}(\mathcal{A}, \text{d}))}(\mathcal{M}_1, \mathcal{I})[1]
}
$$
호모토피 범주에서의 준동형 군들의 긴 완전열이다. 중요한 점은
우리의 짧은 완전열이 허용된다고 가정하지 않았음에도 이를 얻는다는
것이다 (따라서 일반적으로 이 짧은 완전열은 호모토피 범주에서
특별 삼각형을 정의하지 않는다).
\end{remark}

\begin{lemma}
\label{sdga-lemma-product-graded-injective}
$(\mathcal{C}, \mathcal{O})$를 환 달린 사이트라 하자.
$(\mathcal{A}, \text{d})$를 $(\mathcal{C}, \mathcal{O})$ 위의
미분 등급 대수의 층이라 하자. $T$를 집합이라 하고 각
$t \in T$에 대해 $\mathcal{I}_t$를 등급 단사 미분 등급
$\mathcal{A}$-가군이라 하자. 그러면
$\prod \mathcal{I}_t$는 등급 단사 미분 등급 $\mathcal{A}$-가군이다.
\end{lemma}

\begin{proof}
이는 단사 대상들의 곱이 단사이기 때문이다. Homology,
Lemma \ref{homology-lemma-product-injectives}를 보라.
또한 망각 함자를 통해
$\textit{Mod}(\mathcal{A}, \text{d})$에서의 곱은
$\textit{Mod}(\mathcal{A})$에서의 곱과 양립한다.
\end{proof}

\begin{lemma}
\label{sdga-lemma-characterize-graded-injectives-in-dg}
$(\mathcal{C}, \mathcal{O})$를 환 달린 사이트라 하자.
$(\mathcal{A}, \text{d})$를 $(\mathcal{C}, \mathcal{O})$ 위의
미분 등급 대수의 층이라 하자. 집합 $T$가 존재하여 각 $t \in T$에
대해 비순환 미분 등급 $\mathcal{A}$-가군 사이의 단사 사상
$\mathcal{M}_t \to \mathcal{M}'_t$가 주어진다고 하자.
그러면 $\textit{Mod}(\mathcal{A}, \text{d})$의 대상 $\mathcal{I}$에
대해 다음 조건들이 서로 동치이다.
\begin{enumerate}
\item $\mathcal{I}$가 등급 단사이다.
\item 모든 다음 가환 도식에 대해
$$
\xymatrix{
\mathcal{M}_t \ar[r] \ar[d] & \mathcal{I} \\
\mathcal{M}'_t \ar@{..>}[ru]
}
$$
도식을 가환하게 만드는 점선 화살표가
$\textit{Mod}(\mathcal{A}, \text{d})$에 존재한다.
\end{enumerate}
\end{lemma}

\begin{proof}
Lemma \ref{sdga-lemma-characterize-injectives}에서와 같이 $T$와
$\mathcal{N}_t \to \mathcal{N}'_t$를 택하자.
망각 함자를
$F : \textit{Mod}(\mathcal{A}, \text{d}) \to \textit{Mod}(\mathcal{A})$
로 쓰자. Lemma \ref{sdga-lemma-dg-hull}에서처럼 $G$를 $F$의 왼쪽
수반 함자라 하자. 다음과 같이 두자.
$$
\mathcal{M}_t = G(\mathcal{N}_t) \to
G(\mathcal{N}'_t) = \mathcal{M}'_t
$$
이는 Lemma \ref{sdga-lemma-dg-hull-acyclic}에 의해 비순환 미분 등급
$\mathcal{A}$-가군 사이의 단사 사상이다. $G$가 $F$의 왼쪽
수반이므로 다음 도식에 점선 화살표가 존재하는 것은
$$
\xymatrix{
\mathcal{M}_t \ar[r] \ar[d] & \mathcal{I} \\
\mathcal{M}'_t \ar@{..>}[ru]
}
$$
다음 도식에 점선 화살표가 존재하는 것과 동치이다.
$$
\xymatrix{
\mathcal{N}_t \ar[r] \ar[d] & F(\mathcal{I}) \\
\mathcal{N}'_t \ar@{..>}[ru]
}
$$
따라서 $\mathcal{N}_t \to \mathcal{N}_t'$ 사상들의 모음을
% P03-STACKS-E1065: frozen source line 3630 writes \mathcal{N}_t' instead of the established \mathcal{N}'_t notation.
선택한 것으로부터 결과가 따른다.
\end{proof}


\begin{lemma}
\label{sdga-lemma-small-acyclics}
$(\mathcal{C}, \mathcal{O})$를 환 달린 사이트라 하자.
$(\mathcal{A}, \text{d})$를 $(\mathcal{C}, \mathcal{O})$ 위의
미분 등급 대수의 층이라 하자.
% P03-STACKS-E1066: frozen source line 3638 says "for each s" without explicitly restricting s to the set S.
집합 $S$가 존재하여 각 $s$
에 대해 비순환 미분 등급 $\mathcal{A}$-가군 $\mathcal{M}_s$가
주어진다고 하자. 그러면 모든 0이 아닌 비순환 미분 등급
$\mathcal{A}$-가군 $\mathcal{M}$에 대해 어떤 $s \in S$와
$\textit{Mod}(\mathcal{A}, \text{d})$에서의 단사 사상
$\mathcal{M}_s \to \mathcal{M}$이 존재한다.
\end{lemma}

\begin{proof}
시작하기 전에 우리의 관례에 의해 사이트 $\mathcal{C}$가 대상과
사상의 집합 및 피복들의 집합 $\text{Cov}(\mathcal{C})$를 갖는다는
것을 상기하자.
$\mathcal{F}$가 미분 등급 $\mathcal{A}$-가군이면 $|\mathcal{F}|$를
다음 집합의 기수의 합으로 정의하자.
$$
\coprod\nolimits_{(U, n)} \mathcal{F}^n(U)
$$
$U$는 $\mathcal{C}$의 대상을, $n \in \mathbf{Z}$를 모두 돌며
취한다.
기수 $|\Ob(\mathcal{C})|$, $|\text{Arrows}(\mathcal{C})|$,
$|\text{Cov}(\mathcal{C})|$, $\sup |I|$에 대해
$\{U_i \to U\}_{i \in I} \in \text{Cov}(\mathcal{C})$인 경우, 그리고
$|\mathcal{A}|$보다 큰 무한 기수 $\kappa$를 택하자.

\medskip\noindent
미분 등급 $\mathcal{A}$-가군의 포함 $\mathcal{F} \subset \mathcal{M}$를
생각하자. 집합 $K$와 각 $k \in K$에 대해 $\mathcal{C}$의 대상
$U_k$, 정수 $n_k$, 단면 $x_k \in \mathcal{M}^{n_k}(U_k)$로
이루어진 삼중항 $(U_k, n_k, x_k)$가 주어졌다고 하자.
그러면 $\mathcal{F}$와 $k \in K$에 대한 단면 $x_k$들을 포함하는
가장 작은 미분 등급 $\mathcal{A}$-부분가군
$\mathcal{F}' \subset \mathcal{M}$을 생각할 수 있다.
다음을 통해
$$
(\mathcal{F}')^n(U) \subset \mathcal{M}^n(U)
$$
기술할 수 있다. 즉 이는 다음 조건을
만족하는 $x \in \mathcal{M}^n(U)$들의 집합이다. 어떤
$\{f_i : U_i \to U\}_{i \in I} \in \text{Cov}(\mathcal{C})$가
존재하여 각 $i \in I$에 대해 유한 집합 $T_i$와 사상
$g_t : U_i \to U_{k_t}$가 존재하고
$$
f_i^*x = y_i +
\sum\nolimits_{t \in T_i} a_{it}g_t^*x_{k_t} + b_{it}g_t^*\text{d}(x_{k_t})
$$
어떤 $y_i \in \mathcal{F}^n(U)$와
$a_{it} \in \mathcal{A}^{n - n_{k_t}}(U_i)$,
$b_{it} \in \mathcal{A}^{n - n_{k_t} - 1}(U_i)$가 존재한다.
(자세한 내용은 생략한다. 힌트: 이 단면들이 확실히
$\mathcal{F}'$에 속함을 보이고, 역으로 이 규칙이 미분 등급
$\mathcal{A}$-부분가군을 정의함을 보이면 된다.)
이 기술로부터
$|\mathcal{F}'| \leq \max(|\mathcal{F}|, |K|, \kappa)$.

\medskip\noindent
$\mathcal{M}$을 0이 아닌 비순환 미분 등급 $\mathcal{A}$-가군이라
하자. 그러면 어떤 정수 $n$과 $\mathcal{C}$의 대상 $U$ 위에서의
$\mathcal{M}^n$의 0이 아닌 단면 $x$를 찾을 수 있다. 다음을
두자.
$$
\mathcal{F}_0 \subset \mathcal{M}
$$
를 $x$를 포함하는 가장 작은 미분 등급 $\mathcal{A}$-부분가군이라
하자. 앞 문단에 의해
$|\mathcal{F}_0| \leq \kappa$. 귀납적으로 다음과 같이 하자.
$\mathcal{F}_0, \ldots, \mathcal{F}_n$가 주어졌다고 하자.
다음과 같이 $\mathcal{F}_{n + 1}$을 정의한다. 다음 집합을
생각하자.
$$
L = \{(U, n, x)\}
\{U_i \to U\}_{i \in I}, (x_i)_{i \in I})\}
$$
% P03-STACKS-E1067: frozen source lines 3706-3707 contain an unmatched/malformed set display for L.
이는 $U$가 $\mathcal{C}$의 대상이고 $n \in \mathbf{Z}$이며
$\text{d}(x) = 0$인 $x \in \mathcal{F}_n(U)$인 삼중항들의 집합이다.
$\mathcal{M}$이 비순환이므로 각 삼중항 $l = (U_l, n_l, x_l) \in L$에
대해 다음을 택할 수 있다.
% P03-STACKS-E1068: frozen source line 3713 has an unmatched opening parenthesis in the covering-family expression.
$\{(U_{l, i} \to U_l\}_{i \in I_l} \in \text{Cov}(\mathcal{C})$와
$x_{l, i} \in \mathcal{M}^{n_l - 1}(U_{l, i})$에 대해
$\text{d}(x_{l, i}) = x|_{U_{l, i}}$. 그러면
$$
K = \{(U_{l, i}, n_l - 1, x_{l, i}) \mid l \in L, i \in I_l\}
$$
로 두고, $\mathcal{F}_{n + 1}$을 $\mathcal{F}_n$과 단면 $x_{l, i}$들을
포함하는 $\mathcal{M}$의 가장 작은 미분 등급 $\mathcal{A}$-부분가군으로
정의한다. $|K| \leq \max(\kappa, |\mathcal{F}_n|)$이므로
귀납에 의해 $|\mathcal{F}_{n + 1}| \leq \kappa$를 얻는다.

\medskip\noindent
구성에 의해 포함 $\mathcal{F}_n \to \mathcal{F}_{n + 1}$은
코호몰로지 층에서 영사상을 유도한다. 따라서
$\mathcal{F} = \bigcup \mathcal{F}_n$은 0이 아닌 비순환 부분가군이고
$|\mathcal{F}| \leq \kappa$임을 알 수 있다. $|\mathcal{F}|$가
유계인 미분 등급 $\mathcal{A}$-가군 $\mathcal{F}$의 동형류는
집합 하나뿐이므로 결론을 얻는다.
\end{proof}

\begin{definition}
\label{sdga-definition-K-injective}
$(\mathcal{C}, \mathcal{O})$를 환 달린 사이트라 하자.
$(\mathcal{A}, \text{d})$를 $(\mathcal{C}, \mathcal{O})$ 위의
미분 등급 대수의 층이라 하자. 미분 등급 $\mathcal{A}$-가군
% P03-STACKS-E1069: frozen source line 3738 repeats the "diffential" spelling typo.
$\mathcal{I}$가 {\it K-단사}라는 것은 모든 비순환
미분 등급 $\mathcal{M}$에 대해 다음이 성립하는 것이다.
$$
\Hom_{K(\textit{Mod}(\mathcal{A}, \text{d}))}(\mathcal{M}, \mathcal{I}) = 0
$$
\end{definition}

\noindent
Derived Categories, Definition \ref{derived-definition-K-injective}와
유사함에 주목하자.

\begin{lemma}
\label{sdga-lemma-product-K-injective}
$(\mathcal{C}, \mathcal{O})$를 환 달린 사이트라 하자.
$(\mathcal{A}, \text{d})$를 $(\mathcal{C}, \mathcal{O})$ 위의
미분 등급 대수의 층이라 하자. $T$를 집합이라 하고 각
$t \in T$에 대해 $\mathcal{I}_t$를 K-단사
% P03-STACKS-E1070: frozen source line 3756 repeats the spelling "diffential".
미분 등급 $\mathcal{A}$-가군이라 하자. 그러면
$\prod \mathcal{I}_t$는 K-단사 미분 등급 $\mathcal{A}$-가군이다.
\end{lemma}

\begin{proof}
비순환 미분 등급 $\mathcal{A}$-가군 $\mathcal{K}$를 잡자.
그러면 다음이 성립한다.
$$
\Hom_{\textit{Mod}^{dg}(\mathcal{A}, \text{d})}(\mathcal{K},
\prod\nolimits_{t \in T} \mathcal{I}_t)
=
\prod\nolimits_{t \in T}
\Hom_{\textit{Mod}^{dg}(\mathcal{A}, \text{d})}(\mathcal{K}, \mathcal{I}_t)
$$
$\textit{Mod}(\mathcal{A}, \text{d})$에서 곱을 취하는 것이
등급 $\mathcal{A}$-가군으로 가는 망각 함자와 가환하기 때문이다.
곱을 취하는 것은 아벨 군의 범주에서 완전 함자이므로 결론을 얻는다.
\end{proof}

\begin{lemma}
\label{sdga-lemma-first-property-dg-injective}
$(\mathcal{C}, \mathcal{O})$를 환 달린 사이트라 하자.
$(\mathcal{A}, \text{d})$를 $(\mathcal{C}, \mathcal{O})$ 위의
미분 등급 대수의 층이라 하자. $\mathcal{I}$를
$\textit{Mod}(\mathcal{A}, \text{d})$의 K-단사이자 등급 단사인
대상이라 하자. $\textit{Mod}(\mathcal{A}, \text{d})$의 임의의
다음 가환 도식에 대해
$$
\xymatrix{
\mathcal{M} \ar[r]_a \ar[d]_b & \mathcal{I} \\
\mathcal{M}' \ar@{..>}[ru]
}
$$
여기서 $b$가 단사이고 $\mathcal{M}$이 비순환이면 도식을 가환하게
만드는 점선 화살표가 존재한다.
% P03-STACKS-E1071: frozen source line 3791 omits punctuation or a connective between the acyclicity condition and the conclusion.
\end{lemma}

\begin{proof}
$\mathcal{M}$이 비순환이고 $\mathcal{I}$가 K-단사이므로
등급 $\mathcal{A}$-가군 사상
$h : \mathcal{M} \to \mathcal{I}$의 차수가 $-1$인
가 존재하여 $a = \text{d}(h)$이다. $\mathcal{I}$가 등급 단사이고
$b$가 단사이므로 차수 $-1$인 등급 $\mathcal{A}$-가군 사상
$h' : \mathcal{M}' \to \mathcal{I}$가 존재하여
$h = h' \circ b$이다. 그러면 점선 화살표로
$a' = \text{d}(h')$를 취할 수 있다.
\end{proof}

\begin{lemma}
\label{sdga-lemma-second-property-dg-injective}
$(\mathcal{C}, \mathcal{O})$를 환 달린 사이트라 하자.
$(\mathcal{A}, \text{d})$를 $(\mathcal{C}, \mathcal{O})$ 위의
미분 등급 대수의 층이라 하자. $\mathcal{I}$를
$\textit{Mod}(\mathcal{A}, \text{d})$의 K-단사이자 등급 단사인
대상이라 하자. $\textit{Mod}(\mathcal{A}, \text{d})$의 임의의
다음 가환 도식에 대해
$$
\xymatrix{
\mathcal{M} \ar[r]_a \ar[d]_b & \mathcal{I} \\
\mathcal{M}' \ar@{..>}[ru]
}
$$
여기서 $b$가 유사동형이면 도식을 호모토피까지 가환하게 만드는
점선 화살표가 존재한다.
% P03-STACKS-E1072: frozen source lines 3820-3821 omit punctuation or a connective between the quasi-isomorphism condition and the conclusion.
\end{lemma}

\begin{proof}
 $\mathcal{M}'$를 $\mathcal{M}'$와 $\mathcal{M}$의 항등사상 위의
원뿔(이는 비순환이다)의 직합으로 바꾸면 $b$도 단사라고 가정할
수 있다. 그러면 $b$의 cokernel $\mathcal{Q}$는 비순환이다.
따라서
$$
\Hom_{K(\textit{Mod}(\mathcal{A}, \text{d}))}(\mathcal{Q}, \mathcal{I}) =
\Hom_{K(\textit{Mod}(\mathcal{A}, \text{d}))}(\mathcal{Q}, \mathcal{I})[1] = 0
$$
$\mathcal{I}$가 K-단사이기 때문이다. $\mathcal{I}$가 등급 단사이고
비고 \ref{sdga-remark-why-graded-injective}에 의해 다음을 얻는다.
$$
\Hom_{K(\textit{Mod}(\mathcal{A}, \text{d}))}(\mathcal{M}', \mathcal{I})
\longrightarrow
\Hom_{K(\textit{Mod}(\mathcal{A}, \text{d}))}(\mathcal{M}, \mathcal{I})
$$
는 전단사이고 증명이 끝난다.
\end{proof}

\begin{lemma}
\label{sdga-lemma-better-set-of-monos}
$(\mathcal{C}, \mathcal{O})$를 환 달린 사이트라 하자.
$(\mathcal{A}, \text{d})$를 $(\mathcal{C}, \mathcal{O})$ 위의
미분 등급 대수의 층이라 하자. 집합 $R$이 존재하여 각 $r \in R$에
대해 비순환 미분 등급 $\mathcal{A}$-가군 사이의 단사 사상
$\mathcal{M}_r \to \mathcal{M}'_r$가 주어진다고 하자.
그러면 $\textit{Mod}(\mathcal{A}, \text{d})$의 대상 $\mathcal{I}$에
대해 다음 조건들이 서로 동치이다.
\begin{enumerate}
\item $\mathcal{I}$가 K-단사이자 등급 단사이다.
\item 모든 다음 가환 도식에 대해
$$
\xymatrix{
\mathcal{M}_r \ar[r] \ar[d] & \mathcal{I} \\
\mathcal{M}'_r \ar@{..>}[ru]
}
$$
도식을 가환하게 만드는 점선 화살표가
$\textit{Mod}(\mathcal{A}, \text{d})$에 존재한다.
\end{enumerate}
\end{lemma}

\begin{proof}
Lemma \ref{sdga-lemma-characterize-graded-injectives-in-dg}에서와 같이
$T$와 $\mathcal{M}_t \to \mathcal{M}'_t$를 택하자.
Lemma \ref{sdga-lemma-small-acyclics}에서와 같이 $S$와
$\mathcal{M}_s$를 택하자.
비순환 미분 등급 $\mathcal{A}$-가군 사이에서 호모토픽하게 0인
단사 사상 $\mathcal{M}_s \to \mathcal{M}'_s$를 택하자.
이는 $\mathcal{M}'_s$를 항등사상 위의 원뿔로 택할 수 있으므로
가능하다. 이 경우에는 $\mathcal{M}'_s$의 항등사상 자체가
호모토픽하게 0이기도 하다. Differential Graded Algebra, Lemma
\ref{dga-lemma-id-cone-null}을 보라. 이는 Section
\ref{sdga-section-conclude-triangulated}의 논의에 의해 적용된다.
주어진 사상들과 함께 $R = T \coprod S$가 작동한다고
주장하자.

\medskip\noindent
함의 (1) $\Rightarrow$ (2)는 Lemma
\ref{sdga-lemma-first-property-dg-injective}에 의해 성립한다.

\medskip\noindent
 (2)를 가정하자. 먼저 Lemma
\ref{sdga-lemma-characterize-graded-injectives-in-dg}에 의해
$\mathcal{I}$가 등급 단사임을 알 수 있다. 다음으로
비순환 미분 등급 $\mathcal{A}$-가군 $\mathcal{M}$을 잡자.
다음을 보여야 한다.
$$
\Hom_{K(\textit{Mod}(\mathcal{A}, \text{d}))}(\mathcal{M}, \mathcal{I}) = 0
$$
증명은 Injectives, Lemma
\ref{injectives-lemma-characterize-K-injective}의 증명과 정확히
같다.

\medskip\noindent
순서수 $\alpha$에 대한 귀납으로 다음과 같이 비순환 미분 등급
부분가군 $\mathcal{K}_\alpha \subset \mathcal{M}$을 구성하자.
$\alpha = 0$이면 $\mathcal{K}_0 = 0$으로 둔다. $\alpha > 0$이면
다음과 같이 진행한다.
\begin{enumerate}
\item $\alpha = \beta + 1$이고 $\mathcal{K}_\beta = \mathcal{M}$이면
$\mathcal{K}_\alpha = \mathcal{K}_\beta$로 둔다.
\item $\alpha = \beta + 1$이고 $\mathcal{K}_\beta \not = \mathcal{M}$이면
$\mathcal{M}/\mathcal{K}_\beta$는 0이 아닌 비순환 미분 등급
$\mathcal{A}$-가군이다. Lemma \ref{sdga-lemma-small-acyclics}에 따라
어떤 $s \in S$에 대한 $\mathcal{M}_s$와 동형인 미분 등급
$\mathcal{A}$-부분가군
$\mathcal{N}_\alpha \subset \mathcal{M}/\mathcal{K}_\beta$를 택한다.
마지막으로 $\mathcal{K}_\alpha \subset \mathcal{M}$을
$\mathcal{N}_\alpha$의 원상으로 둔다.
\item $\alpha$가 극한 순서수이면
% P03-STACKS-E1073: frozen source line 3917 writes K_beta = colim K_alpha, reversing the expected limit-stage indices.
$\mathcal{K}_\beta = \colim \mathcal{K}_\alpha$로 둔다.
\end{enumerate}
적절히 큰 순서수 $\alpha$에 대해 $\mathcal{M} = \mathcal{K}_\alpha$임은
분명하다. 다음이 성립함을
$$
\Hom_{K(\textit{Mod}(\mathcal{A}, \text{d}))}(\mathcal{K}_\alpha, \mathcal{I})
$$
가 $\alpha$에 대한 초한 귀납으로 0임을 보이겠다.
$\alpha = 0$일 때는 $\mathcal{K}_0$가 0이므로 성립한다.
$\beta$에 대해 성립한다고 하고 $\alpha = \beta + 1$이라 하자.
위 목록의 경우 (1)에서는 결과가 자명하다. 경우 (2)에서는
다음 짧은 완전열이 있다.
$$
0 \to \mathcal{K}_\beta \to \mathcal{K}_\alpha \to \mathcal{N}_\alpha \to 0
$$
비고 \ref{sdga-remark-why-graded-injective}와 $\mathcal{I}$가 등급
단사임을 이미 보았다는 사실에 의해 다음 완전열을 얻는다.
$$
\Hom_{K(\textit{Mod}(\mathcal{A}, \text{d}))}(\mathcal{K}_\beta, \mathcal{I})
\to
\Hom_{K(\textit{Mod}(\mathcal{A}, \text{d}))}(\mathcal{K}_\alpha, \mathcal{I})
\to
\Hom_{K(\textit{Mod}(\mathcal{A}, \text{d}))}(\mathcal{N}_\alpha, \mathcal{I})
$$
귀납에 의해 왼쪽 항은 0이다. 가정 (2)에 의해 오른쪽 항도 0이다.
임의의 사상 $\mathcal{M}_s \to \mathcal{I}$는
$\mathcal{M}'_s$를 통해 인수분해되므로 호모토픽하게 0이기 때문이다.
따라서 가운데 군도 0이다. 마지막으로 $\alpha$가 극한 순서수라고
하자. 등급 $\mathcal{A}$-가군으로서
% P03-STACKS-E1074: frozen source line 3946 repeats K_alpha on both sides of the colimit instead of using the limit-stage index.
$\mathcal{K}_\alpha = \colim \mathcal{K}_\alpha$도 가지므로
다음이 성립한다.
$$
\Hom_{\textit{Mod}^{dg}(\mathcal{A}, \text{d})}
(\mathcal{K}_\alpha, \mathcal{I})
= \lim_{\beta < \alpha}
\Hom_{\textit{Mod}^{dg}(\mathcal{A}, \text{d})}
(\mathcal{K}_\beta, \mathcal{I})
$$
아벨 군의 복합체로서, 이 복합체들의 코호몰로지 군은
이동 사이의 $K(\textit{Mod}(\mathcal{A}, \text{d}))$에서의
사상을 계산한다. $\mathcal{I}$가 등급 단사이므로 비고
\ref{sdga-remark-why-graded-injective}에 의해 이 복합체 계의 전이 사상은
전사이다. 더 나아가 극한 순서수 $\beta \leq \alpha$에 대해서는
극한과 값이 일치한다. 따라서 Homology, Lemma
\ref{homology-lemma-ML-over-ordinals}를 적용하여 결론을 얻는다.
\end{proof}

\begin{lemma}
\label{sdga-lemma-functor-set-of-monos}
$(\mathcal{C}, \mathcal{O})$를 환 달린 사이트라 하자.
$(\mathcal{A}, \text{d})$를 $(\mathcal{C}, \mathcal{O})$ 위의 미분 등급
대수의 층이라 하자. 집합 $R$를 택하고, 각 $r \in R$에 대해
$\mathcal{M}_r \to \mathcal{M}'_r$인 비순환 미분 등급
$\mathcal{A}$-가군의 단사 사상이 주어졌다고 하자.
다음 조건을 만족하는 함자 $M : \textit{Mod}(\mathcal{A}, \text{d}) \to
\textit{Mod}(\mathcal{A}, \text{d})$와 자연 변환
$j : \text{id} \to M$가 존재한다.
\begin{enumerate}
\item $j_\mathcal{M} : \mathcal{M} \to M(\mathcal{M})$는 단사이고
유사동형이다.
\item 모든 실선 도식에 대해
$$
\xymatrix{
\mathcal{M}_r \ar[r] \ar[d] & \mathcal{M} \ar[d]^{j_\mathcal{M}} \\
\mathcal{M}'_r \ar@{..>}[r] & M(\mathcal{M})
}
$$
도식을 가환하게 하는 점선 화살표가
$\textit{Mod}(\mathcal{A}, \text{d})$에 존재한다.
\end{enumerate}
\end{lemma}

\begin{proof}
다음 도식의 밀어내기로 $M(\mathcal{M})$를 정의한다.
$$
\xymatrix{
\bigoplus_{(r, \varphi)} \mathcal{M}_r \ar[r] \ar[d] &
\mathcal{M} \ar[d] \\
\bigoplus_{(r, \varphi)} \mathcal{M}'_r \ar[r] &
M(\mathcal{M})
}
$$
여기서 직합은 모든 순서쌍 $(r, \varphi)$에 대해 취하며,
$r \in R$이고 $\varphi \in
\Hom_{\textit{Mod}(\mathcal{A}, \text{d})}(\mathcal{M}_r, \mathcal{M})$인 경우이다.
단사 사상의 밀어내기도 단사이므로 $\mathcal{M} \to M(\mathcal{M})$가
단사임을 알 수 있다. 왼쪽 수직 화살표의 cokernel이 비순환이므로,
$\mathcal{M} \to M(\mathcal{M})$의 그것과 동형인 cokernel도 비순환이다.
따라서 $\mathcal{M} \to M(\mathcal{M})$는 유사동형이다.
성질 (2)는 구성에 의해 성립한다. 이 절차를 함자로 만들 수 있다는
확인은 생략한다.
\end{proof}

\begin{theorem}
\label{sdga-theorem-qis-into-dg-injective}
$(\mathcal{C}, \mathcal{O})$를 환 달린 사이트라 하자.
$(\mathcal{A}, \text{d})$를 $(\mathcal{C}, \mathcal{O})$ 위의 미분 등급
대수의 층이라 하자. 모든 미분 등급 $\mathcal{A}$-가군 $\mathcal{M}$에 대해,
$\mathcal{I}$가 등급 단사이자 K-단사인 미분 등급
$\mathcal{A}$-가군이 되도록 하는 유사동형 $\mathcal{M} \to \mathcal{I}$가
존재한다. 또한 이 구성은 $\mathcal{M}$에 대해 함자적이다.
\end{theorem}

\begin{proof}
Lemma \ref{sdga-lemma-better-set-of-monos}에서 얻은 집합
$R$와 $\mathcal{M}_r \to \mathcal{M}'_r$에 속하는
$\textit{Mod}(\mathcal{A}, \text{d})$의 사상들을 택하자.
이 $R$와 $\mathcal{M}_r \to \mathcal{M}'_r$를 사용하여
Lemma \ref{sdga-lemma-functor-set-of-monos}에서 구성한 함자 $M$과
자연 변환 $\text{id} \to M$을 택하자.
초한 재귀를 사용하여 함자열
$M_\alpha$와 자연 변환 $M_\beta \to M_\alpha$를
$\alpha < \beta$일 때 다음과 같이 정의한다.
\begin{enumerate}
\item $M_0 = \text{id}$로 둔다.
\item $M_{\alpha + 1} = M \circ M_\alpha$로 두고,
$\beta < \alpha + 1$에 대한 자연 변환 $M_\beta \to M_{\alpha + 1}$는
이미 구성한 $M_\beta \to M_\alpha$와
$\text{id} \to M$에서 오는 사상 $M_\alpha \to M \circ M_\alpha$로부터
얻는다.
\item $\alpha$가 극한 순서수이면
$M_\alpha = \colim_{\beta < \alpha} M_\beta$로 두고, 여입사들을
% P03-STACKS-E1075: frozen source line 4045 reverses the ordinal inequality in the coprojection clause.
$\alpha < \beta$에 대한 변환 $M_\beta \to M_\alpha$로 둔다.
\end{enumerate}
모든 미분 등급 $\mathcal{A}$-가군에 대해 사상
$\mathcal{M} \to M_\beta(\mathcal{M}) \to M_\alpha(\mathcal{M})$가
단사 유사동형임을 관찰하자(여과 여극한이 완전하기 때문이다).

\medskip\noindent
$\textit{Mod}(\mathcal{A}, \text{d})$가 그로텐디크 아벨 범주임을
상기하자. 따라서 Injectives, Proposition
\ref{injectives-proposition-objects-are-small}을
($r \in R$인 모든 $\mathcal{M}_r$의 직합에 적용하여) 사용하면,
모든 $r \in R$에 대해 $\mathcal{M}_r$가 단사 사상에 관하여
$\alpha$-작아지는 극한 순서수 $\alpha$가 존재한다.
$\mathcal{M} \to M_\alpha(\mathcal{M})$가 $\mathcal{M}$을
등급 단사 K-단사 가군에 넣는 원하는 함자적 매장이라고 주장한다.

\medskip\noindent
실제로 임의의 사상 $\mathcal{M}_r \to M_\alpha(\mathcal{M})$는
어떤 $\beta < \alpha$에 대해 $M_\beta(\mathcal{M})$를 통해 인수분해된다.
그런데 $M$의 구성에 의해 이는
$\mathcal{M}_r \to M_{\beta + 1}(\mathcal{M}) = M(M_\beta(\mathcal{M}))$가
$\mathcal{M}'_r$를 통해 인수분해된다는 뜻이다.
$M_\beta(\mathcal{M}) \subset  M_{\beta + 1}(\mathcal{M})
\subset M_\alpha(\mathcal{M})$이므로 $M_\alpha(\mathcal{M})$으로의
원하는 인수분해를 얻는다.
% P03-STACKS-E1076: frozen source line 4069 misspells “factorization” as “factorizaton”.
Lemma \ref{sdga-lemma-better-set-of-monos}에서 $R$와
$\mathcal{M}_r \to \mathcal{M}'_r$를 선택한 방식에 의해 결론을 얻는다.
\end{proof}










\section{유도 범주}
\label{sdga-section-derived}

\noindent
이 절은 Differential Graded Algebra, Section
\ref{dga-section-derived}의 유사물이다.

\medskip\noindent
$(\mathcal{C}, \mathcal{O})$를 환 달린 사이트라 하자. 
$(\mathcal{A}, \text{d})$를 $(\mathcal{C}, \mathcal{O})$ 위의 미분 등급 대수
층이라 하자. $K(\textit{Mod}(\mathcal{A}, \text{d}))$에서 유사동형을
역으로 만들어 유도 범주 $D(\mathcal{A}, \text{d})$를 구성할 것이다.

\begin{lemma}
\label{sdga-lemma-cohomology-homological}
$(\mathcal{C}, \mathcal{O})$를 환 달린 사이트라 하자. 
$(\mathcal{A}, \text{d})$를 $(\mathcal{C}, \mathcal{O})$ 위의 미분 등급 대수
층이라 하자. 다음
$H^0 : \textit{Mod}(\mathcal{A}, \text{d}) \to \textit{Mod}(\mathcal{O})$
함자는 Section \ref{sdga-section-modules}의 함자를 통해 인수분해된다.
$$
H^0 : K(\textit{Mod}(\mathcal{A}, \text{d})) \to \textit{Mod}(\mathcal{O})
$$
이 함자는 Derived Categories, Definition
\ref{derived-definition-homological}의 의미에서 호몰로지적이다.
\end{lemma}

\begin{proof}
정의에서 즉시 다음 가환 도식이 존재함을 얻는다.
$$
\xymatrix{
\textit{Mod}(\mathcal{A}, \text{d}) \ar[r] \ar[d] &
K(\textit{Mod}(\mathcal{A}, \text{d})) \ar[d] \\
\text{Comp}(\mathcal{O}) \ar[r] &
K(\textit{Mod}(\mathcal{O}))
}
$$
$H^0(\mathcal{M})$는 $\mathcal{M}$의 바탕이 되는
$\mathcal{O}$-가군 복합체의 영차 코호몰로지 층으로 정의된다.
따라서 이 보조정리는 $\mathcal{O}$-가군 복합체의 경우에서 따른다.
이는 Derived Categories, Lemma
\ref{derived-lemma-cohomology-homological}의 특수한 경우이다.
\end{proof}

\begin{lemma}
\label{sdga-lemma-acyclics}
$(\mathcal{C}, \mathcal{O})$를 환 달린 사이트라 하자. 
$(\mathcal{A}, \text{d})$를 $(\mathcal{C}, \mathcal{O})$ 위의 미분 등급 대수
층이라 하자. 비순환 가군으로 이루어진 호모토피 범주
$K(\textit{Mod}(\mathcal{A}, \text{d}))$의 충만 부분범주 $\text{Ac}$는
$K(\textit{Mod}(\mathcal{A}, \text{d}))$의 엄밀 충만 포화
삼각 부분범주이다.
\end{lemma}

\begin{proof}
물론 $K(\textit{Mod}(\mathcal{A}, \text{d}))$의 대상 $\mathcal{M}$이
$\text{Ac}$에 속한다는 것은 모든 $i$에 대해
$H^i(\mathcal{M}) = H^0(\mathcal{M}[i])$가 0이라는 것과 동치이다.
따라서 이 보조정리는 이 사실과 Lemma
\ref{sdga-lemma-cohomology-homological}, Derived Categories, Lemma
\ref{derived-lemma-homological-functor-kernel}에서 따른다.
또한 Derived Categories, Definitions
\ref{derived-definition-saturated} 및
\ref{derived-definition-triangulated-subcategory}, 그리고
Lemma \ref{derived-lemma-triangulated-subcategory}를 보라.
\end{proof}

\begin{lemma}
\label{sdga-lemma-qis}
$(\mathcal{C}, \mathcal{O})$를 환 달린 사이트라 하자. 
$(\mathcal{A}, \text{d})$를 $(\mathcal{C}, \mathcal{O})$ 위의 미분 등급 대수
층이라 하자.
유사동형으로 이루어진 부분류
$\text{Qis} \subset \text{Arrows}(K(\textit{Mod}(\mathcal{A}, \text{d})))$를
생각하자. 이는 포화 곱셈계이며
$K(\textit{Mod}(\mathcal{A}, \text{d}))$ 위의 삼각 구조와 양립한다.
\end{lemma}

\begin{proof}
 $f , g : \mathcal{M} \to \mathcal{N}$가 호모토픽인
$\textit{Mod}(\mathcal{A}, \text{d})$의 사상이라고 하자.
그러면 $f$가 유사동형인 것은 $g$가 유사동형인 것과 동치이다.
실제로 Lemma \ref{sdga-lemma-cohomology-homological}에 의해 사상
$H^i(f) = H^0(f[i])$와 $H^i(g) = H^0(g[i])$는
서로 같다. 따라서 호모토피 범주
$K(\textit{Mod}(\mathcal{A}, \text{d}))$의 사상이 유사동형이라고
말하는 것은 모호하지 않다.
``곱셈계'', ``포화'', ``삼각 구조와 양립''의 정의는
Derived Categories, Definition \ref{derived-definition-localization}
및 Categories, Definitions
\ref{categories-definition-multiplicative-system}과
\ref{categories-definition-saturated-multiplicative-system}을 보라.

\medskip\noindent
이 보조정리를 실제로 증명하기 위해 삼각 범주 사이의 완전 함자의 합성
$$
K(\textit{Mod}(\mathcal{A}, \text{d}))
\longrightarrow
K(\textit{Mod}(\mathcal{O}))
\longrightarrow
D(\mathcal{O})
$$
를 생각하자. $K(\textit{Mod}(\mathcal{A}, \text{d}))$의 사상
$f : \mathcal{M} \to \mathcal{N}$가 $\text{Qis}$에 속하는 것은
그 사상이 $D(\mathcal{O})$에서 동형사상으로 보내지는 것과 동치이다.
따라서 이 보조정리는 Derived Categories, Lemma
\ref{derived-lemma-triangle-functor-localize}.
\end{proof}

\noindent
Lemma \ref{sdga-lemma-qis}의 상황에서는 Derived Categories, Proposition
\ref{derived-proposition-construct-localization}
을 적용하여 삼각 범주 사이의 완전 함자
$$
Q :
K(\textit{Mod}(\mathcal{A}, \text{d}))
\longrightarrow
\text{Qis}^{-1}K(\textit{Mod}(\mathcal{A}, \text{d}))
$$
그러나 $\textit{Mod}(\mathcal{A}, \text{d})$는 ``큰'' 범주,
즉 대상들이 진클래스를 이루는 범주이므로, $\mathcal{M}$과
$\mathcal{N}$이 주어졌을 때 $\text{Qis}^{-1}K(\textit{Mod}(\mathcal{A}, \text{d}))$
의 구성이 다음 사상의 {\it set}을 만든다는 것은 즉시 분명하지 않다.
$$
\Mor_{\text{Qis}^{-1}K(\textit{Mod}(\mathcal{A}, \text{d}))}
(\mathcal{M}, \mathcal{N})
$$
의 사상들이다. 그러나 K-단사 복합체를 구성했으므로 이는 실제로 성립한다.
즉 Theorem \ref{sdga-theorem-qis-into-dg-injective}에 의해 유사동형
$s_0 : \mathcal{N} \to \mathcal{I}$를 택할 수 있고, 여기서
$\mathcal{I}$는 등급 단사이자 K-단사인 미분 등급
$\mathcal{A}$-가군이다. 이제 위에 표시한 집합의 원소는 순서쌍
$(f : \mathcal{M} \to \mathcal{N}', s : \mathcal{N} \to \mathcal{N}')$
의 동치류임을 상기하자. 여기서 $f$는
$K(\textit{Mod}(\mathcal{A}, \text{d}))$의 임의의 사상이고 $s$는 유사동형이다.
% P03-STACKS-E1077: frozen source line 4224 misspells “quasi-isomorphism” as “quasi-isomorphsm”.
왼쪽 분수 계산에 대한 설명은 Categories, Section
\ref{categories-section-localization}을 보라.
Lemma \ref{sdga-lemma-second-property-dg-injective}에 의해
에 의해 다음 도식의 점선 화살표를 택할 수 있다.
$$
\xymatrix{
\mathcal{M} \ar[rd]^f & &
\mathcal{N} \ar[ld]_s \ar[rd]^{s_0} \\
& \mathcal{N}' \ar@{..>}[rr]^{s'} & & \mathcal{I}
}
$$
이 도식은 호모토피 범주에서 가환한다.
따라서 순서쌍 $(f, s)$는 순서쌍 $(s' \circ f, s_0)$와 동치이고,
동치류들의 모음은 집합을 이룬다.

\begin{definition}
\label{sdga-definition-derived-category}
$(\mathcal{C}, \mathcal{O})$를 환 달린 사이트라 하자. 
$(\mathcal{A}, \text{d})$를 $(\mathcal{C}, \mathcal{O})$ 위의 미분 등급 대수
층이라 하자. $\text{Qis}$를 Lemma \ref{sdga-lemma-qis}에서와 같이 두자.
{\it $(\mathcal{A}, \text{d})$의 유도 범주}는 다음 삼각
범주이다.
$$
D(\mathcal{A}, \text{d}) =
\text{Qis}^{-1}K(\textit{Mod}(\mathcal{A}, \text{d}))
$$
위에서 더 자세히 논의한 바와 같다.
\end{definition}

\noindent
이 구성에 관한 몇 가지 사실을 증명하자.

\begin{lemma}
\label{sdga-lemma-kernel-localization}
Definition \ref{sdga-definition-derived-category}에서 국소화 함자
$Q : K(\textit{Mod}(\mathcal{A}, \text{d})) \to D(\mathcal{A}, \text{d})$
의 핵은 Lemma \ref{sdga-lemma-acyclics}의 범주 $\text{Ac}$이다.
\end{lemma}

\begin{proof}
이는 Derived Categories, Lemma
\ref{derived-lemma-kernel-localization}과
$0 \to \mathcal{M}$이 유사동형인 것은 $\mathcal{M}$이 비순환인 것과
동치라는 사실에서 즉시 따른다.
\end{proof}

\begin{lemma}
\label{sdga-lemma-H0-over-D}
Definition \ref{sdga-definition-derived-category}에서 함자
$H^0 : K(\textit{Mod}(\mathcal{A}, \text{d})) \to
\textit{Mod}(\mathcal{O})$는 호몰로지 함자
$H^0 : D(\mathcal{A}, \text{d}) \to \textit{Mod}(\mathcal{O})$를 통해 인수분해된다.
\end{lemma}

\begin{proof}
Derived Categories, Lemma
\ref{derived-lemma-universal-property-localization}에서 즉시 따른다.
\end{proof}

\noindent
다음은 유도 범주에서 사상 집합을 계산하는 앞서 약속한 보조정리이다.

\begin{lemma}
\label{sdga-lemma-hom-derived}
$(\mathcal{C}, \mathcal{O})$를 환 달린 사이트라 하자. 
$(\mathcal{A}, \text{d})$를 $(\mathcal{C}, \mathcal{O})$ 위의 미분 등급 대수
층이라 하자.
$\mathcal{M}$과 $\mathcal{N}$을 미분 등급 $\mathcal{A}$-가군이라 하자.
$\mathcal{N} \to \mathcal{I}$가 유사동형이고 $\mathcal{I}$가 등급 단사이자
K-단사인 미분 등급 $\mathcal{A}$-가군이라고 하자. 그러면
$$
\Hom_{D(\mathcal{A}, \text{d})}(\mathcal{M}, \mathcal{N}) =
\Hom_{K(\textit{Mod}(\mathcal{A}, \text{d}))}(\mathcal{M}, \mathcal{I})
$$
\end{lemma}

\begin{proof}
$\mathcal{N} \to \mathcal{I}$가 유사동형이므로 다음이 성립한다.
$$
\Hom_{D(\mathcal{A}, \text{d})}(\mathcal{M}, \mathcal{N}) =
\Hom_{D(\mathcal{A}, \text{d})}(\mathcal{M}, \mathcal{I})
$$
Definition \ref{sdga-definition-derived-category} 앞의 논의에서
Lemma \ref{sdga-lemma-second-property-dg-injective}를 사용하여,
 $\mathcal{M} \to \mathcal{I}$가
$D(\mathcal{A}, \text{d})$에서 표현될 수 있고,
$f : \mathcal{M} \to \mathcal{I}$라는 사상으로
$K(\textit{Mod}(\mathcal{A}, \text{d}))$에서 나타남을 보았다.
이제 $f, f' :  \mathcal{M} \to \mathcal{I}$가
$K(\textit{Mod}(\mathcal{A}, \text{d}))$의 두 사상이라고 하자.
이들이 $D(\mathcal{A}, \text{d})$에서 같은 사상을 정의하는 것은
유사동형 $g : \mathcal{I} \to \mathcal{K}$가
 $K(\textit{Mod}(\mathcal{A}, \text{d}))$에 존재하여
$g \circ f = g \circ f'$가 되는 것과 동치이다. Categories, Lemma
\ref{categories-lemma-equality-morphisms-left-localization}을 보라.
그러나 Lemma \ref{sdga-lemma-second-property-dg-injective}에 의해
사상
$h : \mathcal{K} \to \mathcal{I}$
$h \circ g = \text{id}_\mathcal{I}$가
% P03-STACKS-E1078: frozen source line 4329 repeats “in” before the K-category clause.
$K(\textit{Mod}(\mathcal{A}, \text{d}))$에서 성립한다.
따라서 $g \circ f = g \circ f'$이면 $f = f'$이고
증명이 끝난다.
\end{proof}

\begin{lemma}
\label{sdga-lemma-derived-products}
$(\mathcal{C}, \mathcal{O})$를 환 달린 사이트라 하자. 
$(\mathcal{A}, \text{d})$를 $(\mathcal{C}, \mathcal{O})$ 위의 미분 등급 대수
층이라 하자. 그러면
\begin{enumerate}
\item $D(\mathcal{A}, \text{d})$는 직접합과 곱을 모두 갖는다,
\item 직접합은 미분 등급 $\mathcal{A}$-가군들의 직접합을 취하여 얻고,
\item 곱은 K-단사 미분 등급 가군들의 곱을 취하여 얻는다.
\end{enumerate}
\end{lemma}

\begin{proof}
이제 $\textit{Mod}(\mathcal{A}, \text{d})$가 임의의 직접합과 곱을 갖는
아벨 범주이고, 이것들이 $K(\textit{Mod}(\mathcal{A}, \text{d}))$에서
직접합과 곱을 준다는 사실을 사용한다.
Lemma \ref{sdga-lemma-dgm-abelian}과 Lemma \ref{sdga-lemma-homotopy-direct-sums}를 보라.

\medskip\noindent
미분 등급 $\mathcal{A}$-가군들의 족 $\mathcal{M}_j$를 잡자.
미분 등급 $\mathcal{A}$-가군으로서 직접합 $\mathcal{M} = \bigoplus \mathcal{M}_j$를
생각하자.
미분 등급 $\mathcal{A}$-가군 $\mathcal{N}$에 대하여
$\mathcal{N} \to \mathcal{I}$인 유사동형을 택하자. 여기서
$\mathcal{I}$는 미분 등급 $\mathcal{A}$-가군으로서 등급 단사이자
K-단사이다. Theorem \ref{sdga-theorem-qis-into-dg-injective}를 보라.
Lemma \ref{sdga-lemma-hom-derived}를 사용하면 다음을 얻는다.
\begin{align*}
\Hom_{D(\mathcal{A}, \text{d})}(\mathcal{M}, \mathcal{N})
& =
\Hom_{K(\mathcal{A}, \text{d})}(\mathcal{M}, \mathcal{I}) \\
& =
\prod \Hom_{K(\mathcal{A}, \text{d})}(\mathcal{M}_j, \mathcal{I}) \\
& =
\prod \Hom_{D(\mathcal{A}, \text{d})}(\mathcal{M}_j, \mathcal{I})
\end{align*}
% P03-STACKS-E1079: frozen source line 4374 writes D(A, \text{d}) without the calligraphic \mathcal{A} used throughout.
따라서 이 보조정리의 (2)번에서 말한 $D(A, \text{d})$의 직접합이
존재한다.

\medskip\noindent
미분 등급 $\mathcal{A}$-가군들의 족 $\mathcal{M}_j$를 잡자.
% P03-STACKS-E1080: frozen source line 4380 writes \mathcal{M} -> \mathcal{I}_j; the indexed family construction appears to require \mathcal{M}_j -> \mathcal{I}_j.
각 $j$에 대하여 $\mathcal{M} \to \mathcal{I}_j$인 유사동형을 택하자. 여기서
$\mathcal{I}_j$는 미분 등급 $\mathcal{A}$-가군으로서 등급 단사이자
K-단사이다.
미분 등급 $\mathcal{A}$-가군들의 곱 $\mathcal{I} = \prod \mathcal{I}_j$를
생각하자.
Lemmas \ref{sdga-lemma-product-K-injective}와
\ref{sdga-lemma-product-graded-injective}에 의해
$\mathcal{I}$는 미분 등급 $\mathcal{A}$-가군으로서 등급 단사이자
K-단사임을 알 수 있다.
미분 등급 $\mathcal{A}$-가군 $\mathcal{N}$에 대하여
Lemma \ref{sdga-lemma-hom-derived}를 사용하면 다음을 얻는다.
\begin{align*}
\Hom_{D(\mathcal{A}, \text{d})}(\mathcal{N}, \mathcal{I})
& =
\Hom_{K(\mathcal{A}, \text{d})}(\mathcal{N}, \mathcal{I}) \\
& =
\prod \Hom_{K(\mathcal{A}, \text{d})}(\mathcal{N}, \mathcal{I}_j) \\
& =
\prod \Hom_{D(\mathcal{A}, \text{d})}(\mathcal{N}, \mathcal{M}_j)
\end{align*}
따라서 이 보조정리의 (3)번에서 말한 $D(\mathcal{A}, \text{d})$의 곱이
존재한다.
\end{proof}







\section{정준 델타 함자}
\label{sdga-section-canonical-delta-functor}

\noindent
$(\mathcal{C}, \mathcal{O})$를 환 달린 사이트라 하자. 
$(\mathcal{A}, \text{d})$를 $(\mathcal{C}, \mathcal{O})$ 위의 미분 등급 대수
층이라 하자.
다음 함자를 생각하자.
$\textit{Mod}(\mathcal{A}, \text{d}) \to
K(\textit{Mod}(\mathcal{A}, \text{d}))$.
이 함자는 일반적으로 $\delta$-함자가 {\bf 아니다}.
그러나
% P03-STACKS-E1081: frozen source line 4422 writes D(A, \text{d}) instead of the established D(\mathcal{A}, \text{d}).
$\textit{Mod}(\mathcal{A}, \text{d}) \to D(A, \text{d})$가
$\delta$-함자임을 알 수 있다. 이를 보기 위해서는 다음 짧은 완전열에
연관된 사상 $\delta$를 정의해야 한다.
$$
0 \to \mathcal{K} \xrightarrow{a} \mathcal{L} \xrightarrow{b} \mathcal{M} \to 0
$$
아벨 범주 $\textit{Mod}(\mathcal{A}, \text{d})$에서 생각하자.
사상 $a$의 원뿔 $C(a)$와 그 정준 사상
$i : \mathcal{L} \to C(a)$ 및
$p : C(a) \to \mathcal{K}[1]$를 함께 생각하자. Definition
\ref{sdga-definition-cone}을 보라.
미분 등급 $\mathcal{A}$-가군의 준동형
$$
q : C(a) \longrightarrow \mathcal{M}
$$
을 Differential Graded Algebra, Lemma \ref{dga-lemma-cone}에 의해 얻는다.
이는 Section \ref{sdga-section-conclude-triangulated}의 논의에 따라
사용할 수 있으며, 다음 도식에 적용한 것이다.
$$
\xymatrix{
\mathcal{K} \ar[r]_a \ar[d] &
\mathcal{L} \ar[d]^b \\
0 \ar[r] &
\mathcal{M}
}
$$
사상 $q$가 유사동형이라는 것은, 예를 들어 $\mathcal{O}$-가군 복합체의
사상 범주에서도 그러하기 때문이다. Derived Categories, Section
\ref{derived-section-canonical-delta-functor}의 논의를 보라.
Differential Graded Algebra, Lemma \ref{dga-lemma-cone-homotopy}에 따르면
(Section \ref{sdga-section-conclude-triangulated}의 논의에 따라 이를 사용할 수 있다)
다음 삼각형
$$
(\mathcal{K}, \mathcal{L}, C(a), a, i, -p)
$$
는 $K(\textit{Mod}(\mathcal{A}, \text{d}))$에서 구별되는 삼각형이다.
국소화 함자
$K(\textit{Mod}(\mathcal{A}, \text{d})) \to D(\mathcal{A}, \text{d})$가
완전 함자이므로 $(\mathcal{K}, \mathcal{L}, C(a), a, i, -p)$가
$D(\mathcal{A}, \text{d})$에서 구별되는 삼각형임을 알 수 있다. $q$는
유사동형이고, $q$는 $D(\mathcal{A}, \text{d})$에서 동형사상이다.
따라서 다음을 얻는다.
$$
(\mathcal{K}, \mathcal{L}, \mathcal{M}, a, b, -p \circ q^{-1})
$$
는 $D(\mathcal{A}, \text{d})$의 구별되는 삼각형이다.
이는 다음 보조정리를 시사한다.

\begin{lemma}
\label{sdga-lemma-derived-canonical-delta-functor}
$(\mathcal{C}, \mathcal{O})$를 환 달린 사이트라 하자. 
$(\mathcal{A}, \text{d})$를 $(\mathcal{C}, \mathcal{O})$ 위의 미분 등급 대수
층이라 하자. 국소화 함자
$\textit{Mod}(\mathcal{A}, \text{d}) \to D(\mathcal{A}, \text{d})$
는 다음과 같이 자연스러운 $\delta$-함자 구조를 갖는다.
$$
\delta_{\mathcal{K} \to \mathcal{L} \to \mathcal{M}} = - p \circ q^{-1}
$$
여기서 $p$와 $q$는 위에서 설명한 것이다.
\end{lemma}

\begin{proof}
복합체의 짧은 완전열이 주어질 때 이 선택이 구별되는 삼각형을 준다는
것은 이미 보았다.
이 구성의 함자성을 보여야 한다. Derived Categories, Definition
\ref{derived-definition-delta-functor}를 보라.
이는 약간의 계산을 하면 Differential Graded Algebra, Lemma
\ref{dga-lemma-cone}에서 따른다(Section
\ref{sdga-section-conclude-triangulated}의 논의에 따라 이를 사용할 수 있다).
Derived Categories, Lemma \ref{derived-lemma-derived-canonical-delta-functor}와
비교하라.
\end{proof}

\begin{lemma}
\label{sdga-lemma-homotopy-colimit}
$(\mathcal{C}, \mathcal{O})$를 환 달린 사이트라 하자. 
$(\mathcal{A}, \text{d})$를 $(\mathcal{C}, \mathcal{O})$ 위의 미분 등급 대수
층이라 하자. 미분 등급 $\mathcal{A}$-가군의 계 $\mathcal{M}_n$을 잡자.
그러면 $D(\mathcal{A}, \text{d})$에서의 유도 여극한
$\text{hocolim} \mathcal{M}_n$은 미분 등급 가군
$\colim \mathcal{M}_n$으로 표현된다.
\end{lemma}

\begin{proof}
 $\mathcal{M} = \colim \mathcal{M}_n$으로 두자.
미분 등급 $\mathcal{A}$-가군의 다음 완전열이 있다.
$$
0 \to \bigoplus \mathcal{M}_n \to \bigoplus \mathcal{M}_n \to \mathcal{M} \to 0
$$
이는 Derived Categories, Lemma \ref{derived-lemma-compute-colimit}에 의해
성립한다(밑에 놓인 $\mathcal{O}$-가군 복합체에 적용한다).
직접합은 Lemma \ref{sdga-lemma-derived-products}에 의해
$D(\mathcal{A}, \text{d})$에서의 직접합이다.
따라서 결과는 Derived Categories, Definition
\ref{derived-definition-derived-colimit}의 유도 여극한 정의와,
복합체의 짧은 완전열이 구별되는 삼각형을 준다는 사실
(Lemma \ref{sdga-lemma-derived-canonical-delta-functor})에서 따른다.
\end{proof}






\section{유도 당김}
\label{sdga-section-derived-pullback}

\noindent
$(f, f^\sharp) : (\Sh(\mathcal{C}), \mathcal{O}_\mathcal{C})
\to (\Sh(\mathcal{D}), \mathcal{O}_\mathcal{D})$를 환 달린 위상 공간의
사상이라 하자. $\mathcal{A}$를 미분 등급
$\mathcal{O}_\mathcal{C}$-대수라 하고 $\mathcal{B}$를
미분 등급 $\mathcal{O}_\mathcal{D}$-대수라 하자.
다음 사상이 주어졌다고 하자.
$$
\varphi : f^{-1}\mathcal{B} \to \mathcal{A}
$$
이는 미분 등급 $f^{-1}\mathcal{O}_\mathcal{D}$-대수의 사상이다.
스칼라 제한과 확장의 수반성에 의해 이는 미분 등급
$\mathcal{O}_\mathcal{C}$-대수의 사상 $\varphi : f^*\mathcal{B} \to \mathcal{A}$와
같은 것이며, 동치로 $\varphi$를 다음 사상으로 볼 수 있다.
$$
\varphi : \mathcal{B} \to f_*\mathcal{A}
$$
이는 미분 등급 $\mathcal{O}_\mathcal{D}$-대수의 사상이다.
Remark \ref{sdga-remark-functoriality-dga}를 보라.

\medskip\noindent
위의 자료에 더하여 $\mathcal{A}'$를 두 번째 미분 등급
$\mathcal{O}_\mathcal{C}$-대수라 하고 $\mathcal{N}$을
미분 등급 $(\mathcal{A}, \mathcal{A}')$-쌍가군이라 하자.
이 설정에서 다음 함자를 생각할 수 있다.
$$
\textit{Mod}(\mathcal{B}, \text{d})
\longrightarrow
\textit{Mod}(\mathcal{A}', \text{d}),\quad
\mathcal{M} \longmapsto f^*\mathcal{M} \otimes_{\mathcal{A}} \mathcal{N}
$$
이는 다음 함자로 확장된다.
$$
\textit{Mod}^{dg}(\mathcal{B}, \text{d})
\longrightarrow
\textit{Mod}^{dg}(\mathcal{A}', \text{d}),\quad
\mathcal{M} \longmapsto f^*\mathcal{M} \otimes_{\mathcal{A}} \mathcal{N}
$$
Section \ref{sdga-section-functoriality-dg}와 Section
\ref{sdga-section-dg-bimodules}의 논의에 따라 미분 등급 범주 사이의 함자이다.
따라서 형식적으로 다음 완전 함자도 얻는다.
\begin{equation}
\label{sdga-equation-pullback}
K(\textit{Mod}(\mathcal{B}, \text{d}))
\longrightarrow
K(\textit{Mod}(\mathcal{A}', \text{d})),\quad
\mathcal{M} \longmapsto f^*\mathcal{M} \otimes_{\mathcal{A}} \mathcal{N}
\end{equation}
삼각 범주 사이의 함자이다.

\begin{lemma}
\label{sdga-lemma-derived-tensor-product}
위의 상황에서 함자 (\ref{sdga-equation-pullback})와 국소화 함자
$K(\textit{Mod}(\mathcal{A}', \text{d})) \to D(\mathcal{A}', \text{d})$
의 합성은 좌유도 확장
$D(\mathcal{B}, \text{d}) \to D(\mathcal{A}', \text{d})$을 갖는다.
이 함자의 좋은 오른쪽 미분 등급 $\mathcal{B}$-가군 $\mathcal{P}$에서의 값은
$f^*\mathcal{P} \otimes_\mathcal{A} \mathcal{N}$이다.
\end{lemma}

\begin{proof}
임의의 (오른쪽) 미분 등급 $\mathcal{B}$-가군 $\mathcal{M}$에 대하여,
$\mathcal{P}$가 좋은 미분 등급 $\mathcal{B}$-가군인 유사동형
$\mathcal{P} \to \mathcal{M}$이 존재함을 상기하자. Lemma
\ref{sdga-lemma-resolve}를 보라.
따라서 Derived Categories, Lemma
\ref{derived-lemma-find-existence-computes}에 의해, 좋은 미분 등급
$\mathcal{B}$-가군들의 유사동형 $\mathcal{P} \to \mathcal{P}'$이 주어졌을 때
유도된 사상
$$
f^*\mathcal{P} \otimes_\mathcal{A} \mathcal{N}
\longrightarrow
f^*\mathcal{P}' \otimes_\mathcal{A} \mathcal{N}
$$
이 유사동형임을 보이면 충분하다. $\mathcal{P} \to \mathcal{P}'$의 원뿔
$\mathcal{P}''$은 Lemma \ref{sdga-lemma-good-admissible-ses}에 의해
% P03-STACKS-E1082: frozen source line 4613 calls the cone a good differential graded \mathcal{A}-module although the preceding cone consists of \mathcal{B}-modules.
좋은 미분 등급 $\mathcal{A}$-가군이다.
구별되는 삼각형
$$
\mathcal{P} \to \mathcal{P}' \to \mathcal{P}'' \to \mathcal{P}[1]
$$
이 $K(\textit{Mod}(\mathcal{B}, \text{d}))$에 있으므로 구별되는 삼각형
$$
f^*\mathcal{P} \otimes_\mathcal{A} \mathcal{N} \to
f^*\mathcal{P}' \otimes_\mathcal{A} \mathcal{N} \to
f^*\mathcal{P}'' \otimes_\mathcal{A} \mathcal{N} \to
f^*\mathcal{P}[1] \otimes_\mathcal{A} \mathcal{N}
$$
을 $K(\textit{Mod}(\mathcal{A}', \text{d}))$에서 얻는다. Lemma
\ref{sdga-lemma-acyclic-good}에 의해 미분 등급 가군
$f^*\mathcal{P}'' \otimes_\mathcal{A} \mathcal{N}$은 비순환이고,
증명이 끝난다.
\end{proof}

\begin{definition}
\label{sdga-definition-pullback}
유도 텐서곱과 유도 당김.
\begin{enumerate}
\item $(\mathcal{C}, \mathcal{O})$를 환 달린 사이트라 하자. 
$\mathcal{A}$와 $\mathcal{B}$를 미분 등급 $\mathcal{O}$-대수라 하자.
$\mathcal{N}$을 미분 등급 $(\mathcal{A}, \mathcal{B})$-쌍가군이라 하자.
Lemma \ref{sdga-lemma-derived-tensor-product}에서 구성한 함자
$D(\mathcal{A}, \text{d}) \to D(\mathcal{B}, \text{d})$를
{\it 유도 텐서곱}이라 하고
$- \otimes_\mathcal{A}^\mathbf{L} \mathcal{N}$으로 나타낸다.
\item $(f, f^\sharp) : (\Sh(\mathcal{C}), \mathcal{O}_\mathcal{C})
\to (\Sh(\mathcal{D}), \mathcal{O}_\mathcal{D})$를 환 달린 위상 공간의
사상이라 하자. $\mathcal{A}$를 미분 등급
$\mathcal{O}_\mathcal{C}$-대수라 하고 $\mathcal{B}$를 미분 등급
$\mathcal{O}_\mathcal{D}$-대수라 하자. 미분 등급
$\mathcal{O}_\mathcal{D}$-대수의 준동형
$\varphi : \mathcal{B} \to f_*\mathcal{A}$를 잡자.
Lemma \ref{sdga-lemma-derived-tensor-product}에서 구성한 함자
$D(\mathcal{B}, \text{d}) \to D(\mathcal{A}, \text{d})$를
% P03-STACKS-E1083: frozen source line 4655 uses the ungrammatical singular 'and denote' after 'is called'; the intended wording is 'and denote it by'.
{\it 유도 당김}이라 하고 $Lf^*$로 나타낸다.
\end{enumerate}
\end{definition}

\noindent
이 용어를 사용하면 몇 가지 명백한 호환성을 표현할 수 있다.

\begin{lemma}
\label{sdga-lemma-compose-pullback-tensor}
Lemma \ref{sdga-lemma-derived-tensor-product}의 함자
$D(\mathcal{B}, \text{d}) \to D(\mathcal{A}', \text{d})$는
$\mathcal{M} \mapsto
Lf^*\mathcal{M} \otimes_\mathcal{A}^\mathbf{L} \mathcal{N}$.
\end{lemma}

\begin{proof}
이는 좋은 미분 등급 가군으로 대상을 나타내어 이 함자들을 계산할 수 있고,
$\mathcal{P}$가 좋은 미분 등급 $\mathcal{B}$-가군이면
$f^*\mathcal{P}$가 좋은 미분 등급 $\mathcal{A}$-가군이라는 사실에서
즉시 따른다.
\end{proof}

\begin{lemma}
\label{sdga-lemma-compose-pullback}
$(f, f^\sharp) : (\Sh(\mathcal{C}), \mathcal{O})
\to (\Sh(\mathcal{C}'), \mathcal{O}')$와
$(g, g^\sharp) : (\Sh(\mathcal{C}'), \mathcal{O}')
\to (\Sh(\mathcal{C}''), \mathcal{O}'')$를 환 달린 위상 공간의
사상이라 하자. $\mathcal{A}$, $\mathcal{A}'$, $\mathcal{A}''$를 각각
미분 등급 $\mathcal{O}$-대수, $\mathcal{O}'$-대수, $\mathcal{O}''$-대수라 하자.
$\varphi : \mathcal{A}' \to f_*\mathcal{A}$ 및
$\varphi' : \mathcal{A}'' \to g_*\mathcal{A}'$를 각각 미분 등급
$\mathcal{O}'$-대수와 $\mathcal{O}''$-대수의 준동형이라 하자.
그러면 $L(g \circ f)^* = Lf^* \circ Lg^* :
D(\mathcal{A}'', \text{d}) \to D(\mathcal{A}, \text{d})$이다.
\end{lemma}

\begin{proof}
이는 좋은 미분 등급 가군으로 대상을 나타내어 이 함자들을 계산할 수 있고,
% P03-STACKS-E1084: frozen source line 4699 says 'of \\mathcal{P} is a good'; the conditional should read 'if \\mathcal{P} is a good'.
$\mathcal{P}$가 좋은 미분 등급 $\mathcal{A}'$-가군이면
$f^*\mathcal{P}$가 좋은 미분 등급 $\mathcal{A}$-가군이라는 사실에서
즉시 따른다.
\end{proof}

\noindent
$(\mathcal{C}, \mathcal{O})$를 환 달린 사이트라 하자. 
$\mathcal{A}$와 $\mathcal{B}$를 미분 등급 $\mathcal{O}$-대수라 하자.
$\mathcal{N} \to \mathcal{N}'$을 미분 등급 $(\mathcal{A}, \mathcal{B})$-쌍가군의
준동형이라 하자. 그러면 정준 사상
$$
t :
\mathcal{M} \otimes_\mathcal{A}^\mathbf{L} \mathcal{N}
\longrightarrow
\mathcal{M} \otimes_\mathcal{A}^\mathbf{L} \mathcal{N}'
$$
$D(\mathcal{A}, \text{d})$에서 $\mathcal{M}$에 대해 함자적인 사상이며,
삼각 범주 사이의 완전 함자
$D(\mathcal{A}, \text{d}) \to D(\mathcal{B}, \text{d})$ 사이의 자연 변환을
정의한다. 좋은 미분 등급 $\mathcal{A}$-가군 $\mathcal{P}$에서 $t$의 값은
다음의 명백한 사상이다.
$$
\mathcal{P} \otimes_\mathcal{A}^\mathbf{L} \mathcal{N} =
\mathcal{P} \otimes_\mathcal{A} \mathcal{N}
\longrightarrow
\mathcal{P} \otimes_\mathcal{A} \mathcal{N}' =
\mathcal{P} \otimes_\mathcal{A}^\mathbf{L} \mathcal{N}'
$$

\begin{lemma}
\label{sdga-lemma-tensor-symmetry}
위의 상황에서 $\mathcal{N} \to \mathcal{N}'$이 코호몰로지 층의 동형사상이면,
$t$는 다음 함자 사이의 동형사상이다.
$(- \otimes_\mathcal{A}^\mathbf{L} \mathcal{N}) \to
(- \otimes_\mathcal{A}^\mathbf{L} \mathcal{N}')$.
\end{lemma}

\begin{proof}
임의의 좋은 미분 등급 $\mathcal{A}$-가군 $\mathcal{P}$에 대하여
$\mathcal{P} \otimes_\mathcal{A} \mathcal{N} \to
\mathcal{P} \otimes_\mathcal{A} \mathcal{N}'$가 코호몰로지 층의
동형사상임을 보이면 충분하다.
이를 위해 왼쪽 미분 등급 $\mathcal{A}$-가군으로서 사상
$\mathcal{N} \to \mathcal{N}'$의 원뿔 $\mathcal{N}''$를 잡자. Definition
\ref{sdga-definition-cone}을 보라. (정확히는 $\mathcal{N}''$도 쌍가군이지만
이를 사용할 필요는 없다.) 텐서 구성의 함자성
(이는 미분 등급 범주의 함자이다)에 의해
$\mathcal{P} \otimes_\mathcal{A} \mathcal{N}''$는 다음 사상의
($\mathcal{O}$-가군 복합체로서의) 원뿔이다.
$\mathcal{P} \otimes_\mathcal{A} \mathcal{N} \to
\mathcal{P} \otimes_\mathcal{A} \mathcal{N}'$.
따라서 $\mathcal{P} \otimes_\mathcal{A} \mathcal{N}''$가 비순환임을
보이면 충분하다. 이는 $\mathcal{P}$가 좋고, $\mathcal{N}''$가 유사동형의
원뿔이므로 비순환이라는 사실에서 따른다.
\end{proof}

\begin{lemma}
\label{sdga-lemma-good-on-other-side}
$(\mathcal{C}, \mathcal{O})$를 환 달린 사이트라 하자. 
$\mathcal{A}$와 $\mathcal{B}$를 미분 등급 $\mathcal{O}$-대수라 하자.
$\mathcal{N}$을 미분 등급 $(\mathcal{A}, \mathcal{B})$-쌍가군이라 하자.
$\mathcal{N}$이 왼쪽 미분 등급 $\mathcal{A}$-가군으로서 좋다면, 모든
미분 등급 $\mathcal{A}$-가군 $\mathcal{M}$에 대하여
$\mathcal{M} \otimes_\mathcal{A}^\mathbf{L} \mathcal{N} =
\mathcal{M} \otimes_\mathcal{A} \mathcal{N}$이다.
\end{lemma}

\begin{proof}
좋은 (오른쪽) 미분 등급 $\mathcal{A}$-가군 $\mathcal{P}$와 유사동형
$\mathcal{P} \to \mathcal{M}$을 택하자.
이 보조정리를 증명하려면 다음 사상이
$\mathcal{P} \otimes_\mathcal{A} \mathcal{N} \to
\mathcal{M} \otimes_\mathcal{A} \mathcal{N}$
가 유사동형임을 보여야 한다. 사상에 대한 원뿔 $C$는
$\mathcal{P} \to \mathcal{M}$에 대한 비순환 오른쪽 미분
등급 $\mathcal{A}$-가군이다. 따라서
$C \otimes_\mathcal{A} \mathcal{N}$은 $\mathcal{N}$이
왼쪽 미분 등급 $\mathcal{A}$-가군으로서 좋다고 가정했으므로 비순환이다.
$C \otimes_\mathcal{A} \mathcal{N}$은 다음 사상의 원뿔이다.
사상 $\mathcal{P} \otimes_\mathcal{A} \mathcal{N} \to
\mathcal{M} \otimes_\mathcal{A} \mathcal{N}$을
$\mathcal{O}$-가군 복합체로 보아 결론을 얻는다.
\end{proof}

\begin{lemma}
\label{sdga-lemma-compose-tensor}
$(\mathcal{C}, \mathcal{O})$를 환 달린 사이트라 하자. 
$\mathcal{A}$, $\mathcal{A}'$, $\mathcal{A}''$를 미분 등급
$\mathcal{O}$-대수라 하자. $\mathcal{N}$과 $\mathcal{N}'$를 각각
미분 등급 $(\mathcal{A}, \mathcal{A}')$-쌍가군과
$(\mathcal{A}', \mathcal{A}'')$-쌍가군이라 하자. 다음 정준 사상이
$$
\mathcal{N} \otimes_{\mathcal{A}'}^\mathbf{L} \mathcal{N}'
\longrightarrow
\mathcal{N} \otimes_{\mathcal{A}'} \mathcal{N}'
$$
 $D(\mathcal{A}'', \text{d})$에서 유사동형이라고 가정하자.
그러면 다음이 성립한다.
$$
(\mathcal{M}
\otimes_\mathcal{A}^\mathbf{L} \mathcal{N})
\otimes_{\mathcal{A}'}^\mathbf{L} \mathcal{N}'
=
\mathcal{M}
\otimes_\mathcal{A}^\mathbf{L}
(\mathcal{N} \otimes_{\mathcal{A}'} \mathcal{N}')
$$
 $D(\mathcal{A}, \text{d}) \to D(\mathcal{A}'', \text{d})$의 함자로서 성립한다.
\end{lemma}

\begin{proof}
좋은 미분 등급 $\mathcal{A}$-가군 $\mathcal{P}$와 유사동형
$\mathcal{P} \to \mathcal{M}$을 택하자. Lemma \ref{sdga-lemma-resolve}를 보라.
그러면
$$
\mathcal{M}
\otimes_\mathcal{A}^\mathbf{L}
(\mathcal{N} \otimes_{\mathcal{A}'} \mathcal{N}') =
\mathcal{P} \otimes_\mathcal{A} \mathcal{N}
\otimes_{\mathcal{A}'} \mathcal{N}'
$$
그리고 다음이 성립한다.
$$
(\mathcal{M}
\otimes_\mathcal{A}^\mathbf{L} \mathcal{N})
\otimes_{\mathcal{A}'}^\mathbf{L} \mathcal{N}' =
(\mathcal{P} \otimes_\mathcal{A} \mathcal{N})
\otimes_{\mathcal{A}'}^\mathbf{L} \mathcal{N}'
$$
따라서 다음 정준 사상이
$$
(\mathcal{P} \otimes_\mathcal{A} \mathcal{N})
\otimes_{\mathcal{A}'}^\mathbf{L} \mathcal{N}'
\longrightarrow
\mathcal{P} \otimes_\mathcal{A} \mathcal{N}
\otimes_{\mathcal{A}'} \mathcal{N}'
$$
유사동형임을 보여야 한다. $\mathcal{Q}$가 좋은 왼쪽 미분 등급
$\mathcal{A}'$-가군인 유사동형 $\mathcal{Q} \to \mathcal{N}'$을 택하자
(Lemma \ref{sdga-lemma-resolve}).
Lemma \ref{sdga-lemma-good-on-other-side}에 의해 위 사상은
$\mathcal{O}$-가군의 유도 범주에서 다음 사상이다.
$$
\mathcal{P} \otimes_\mathcal{A} \mathcal{N}
\otimes_{\mathcal{A}'} \mathcal{Q}
\longrightarrow
\mathcal{P} \otimes_\mathcal{A} \mathcal{N}
\otimes_{\mathcal{A}'} \mathcal{N}'
$$
가정에 의해 $\mathcal{N} \otimes_{\mathcal{A}'} \mathcal{Q} \to
\mathcal{N} \otimes_{\mathcal{A}'} \mathcal{N}'$는 유사동형이고,
% P03-STACKS-E1085: frozen source line 4858 says 'an quasi-isomorphism'; the article should be 'a quasi-isomorphism'.
는 유사동형이고, $\mathcal{P}$는 좋은 미분 등급
$\mathcal{A}$-가군이므로 이 사상은 Lemma
\ref{sdga-lemma-tensor-symmetry}에 의해 유사동형이다(왼쪽과 오른쪽은 각각
왼쪽은 $\mathcal{P} \otimes_\mathcal{A}^\mathbf{L} (\mathcal{N}
\otimes_{\mathcal{A}'} \mathcal{Q})$를 계산하고,
오른쪽은 $\mathcal{P} \otimes_\mathcal{A}^\mathbf{L} (\mathcal{N}
\otimes_{\mathcal{A}'} \mathcal{N}')$를 계산한다. 또는 보조정리의 증명에서
같은 논의를 반복해도 된다).
\end{proof}






\section{유도 직접상}
\label{sdga-section-derived-pushforward}

\noindent
% P03-STACKS-E1086: frozen source line 4876 says 'The existence of enough K-injective'; the plural noun should be 'K-injectives' (or 'enough K-injective modules').
충분한 K-단사 가군이 존재하므로 호모토피 범주의 임의의 완전 함자에 대해
오른쪽 유도 함자를 취할 수 있다.

\begin{lemma}
\label{sdga-lemma-right-derived}
환 달린 사이트 $(\mathcal{C}, \mathcal{O})$를 생각하자.
미분 등급 대수의 층 $(\mathcal{A}, \text{d})$가
$(\mathcal{C}, \mathcal{O})$ 위에 주어졌다고 하자. 그러면 임의의 완전 함자
$$
T : K(\textit{Mod}(\mathcal{A}, \text{d})) \longrightarrow \mathcal{D}
$$
삼각 범주 사이에는 오른쪽 유도 확장
$RT : D(\mathcal{A}, \text{d}) \to \mathcal{D}$가 존재하며,
등급 단사이고 K-단사인 미분 등급 $\mathcal{A}$-가군 $\mathcal{I}$에서의
그 값은 $T(\mathcal{I})$이다.
\end{lemma}

\begin{proof}
정리 \ref{sdga-theorem-qis-into-dg-injective}에 의해 임의의 (오른쪽)
미분 등급 $\mathcal{A}$-가군 $\mathcal{M}$에 대해
어떤 유사동형 $\mathcal{M} \to \mathcal{I}$가 존재한다.
여기서 $\mathcal{I}$는 등급 단사이고 K-단사인
미분 등급 $\mathcal{A}$-가군이다.
그러므로 유도 범주, 보조정리
\ref{derived-lemma-find-existence-computes}에 의해,
등급 단사이고 K-단사인 미분 등급 $\mathcal{A}$-가군들의 유사동형
$\mathcal{I} \to \mathcal{I}'$가 주어졌을 때
$T(\mathcal{I}) \to T(\mathcal{I}')$가 동형임을 보이면 충분하다.
이는 예를 들어 보조정리 \ref{sdga-lemma-hom-derived}에 의해
$\mathcal{I} \to \mathcal{I}'$가
$K(\textit{Mod}(\mathcal{A}, \text{d}))$에서 동형이기 때문이다
(또는 보조정리 \ref{sdga-lemma-second-property-dg-injective}로부터
이 사실을 얻을 수도 있다).
\end{proof}

\noindent
이 결과가 적용되는 함자는 이미 여러 개 보았다.
두 가지 예를 들면 다음과 같다.

\begin{definition}
\label{sdga-definition-pushforward}
유도 내부 Hom과 유도 직접상.
\begin{enumerate}
\item 환 달린 사이트 $(\mathcal{C}, \mathcal{O})$를 생각하자.
$\mathcal{A}$, $\mathcal{B}$를 미분 등급 $\mathcal{O}$-대수라 하자.
$\mathcal{N}$을 미분 등급
$(\mathcal{A}, \mathcal{B})$-쌍가군이라 하자. 내부 Hom 함자의
오른쪽 유도 확장
$$
R\SheafHom_\mathcal{B}(\mathcal{N}, -) :
D(\mathcal{B}, \text{d})
\longrightarrow
D(\mathcal{A}, \text{d})
$$
내부 Hom 함자 $\SheafHom_\mathcal{B}^{dg}(\mathcal{N}, -)$의
이 확장을 {\it 유도 내부 Hom}이라 한다.
\item $(f, f^\sharp) : (\Sh(\mathcal{C}), \mathcal{O}_\mathcal{C})
\to (\Sh(\mathcal{D}), \mathcal{O}_\mathcal{D})$가 환 달린 토포스의
사상이라고 하자. $\mathcal{A}$를 미분 등급
$\mathcal{O}_\mathcal{C}$-대수라 하고, $\mathcal{B}$를
미분 등급 $\mathcal{O}_\mathcal{D}$-대수라 하자. 다음과 같은
미분 등급 $\mathcal{O}_\mathcal{D}$-대수의 준동형
$\varphi : \mathcal{B} \to f_*\mathcal{A}$를 주어졌다고 하자.
직접상 $f_*$의 오른쪽 유도 확장
$$
Rf_* :
D(\mathcal{A}, \text{d})
\longrightarrow
D(\mathcal{B}, \text{d})
$$
을 {\it 유도 직접상}이라 한다.
\end{enumerate}
\end{definition}

\noindent
실제로
$Rf_* : D(\mathcal{A}, \text{d}) \to D(\mathcal{B}, \text{d})$
는 기저 $\mathcal{O}$-가군의 복합체에 대한 유도 직접상과
일치한다. 보조정리 \ref{sdga-lemma-pushforward-agrees}를 보라.

\medskip\noindent
이 함자들은 유도 당김과 유도 텐서곱에 수반된다.

\begin{lemma}
\label{sdga-lemma-derived-adjoint-tensor-hom}
환 달린 사이트 $(\mathcal{C}, \mathcal{O})$를 생각하자.
$\mathcal{A}$, $\mathcal{B}$를 미분 등급 $\mathcal{O}$-대수라 하자.
$\mathcal{N}$을 미분 등급 $(\mathcal{A}, \mathcal{B})$-쌍가군이라 하자.
그러면
$$
R\SheafHom_\mathcal{B}(\mathcal{N}, -) :
D(\mathcal{B}, \text{d})
\longrightarrow
D(\mathcal{A}, \text{d})
$$
는 다음 함자의 오른쪽 수반이다.
$$
- \otimes_\mathcal{A}^\mathbf{L} \mathcal{N} :
D(\mathcal{A}, \text{d})
\longrightarrow
D(\mathcal{B}, \text{d})
$$
\end{lemma}

\begin{proof}
이는 유도 범주, 보조정리
\ref{derived-lemma-pre-derived-adjoint-functors-general}과
보조정리 \ref{sdga-lemma-tensor-hom-adjunction-dg}에 따른다.
\end{proof}

\begin{lemma}
\label{sdga-lemma-derived-adjoint-push-pull}
$(f, f^\sharp) : (\Sh(\mathcal{C}), \mathcal{O}_\mathcal{C})
\to (\Sh(\mathcal{D}), \mathcal{O}_\mathcal{D})$가
환 달린 토포스의 사상이라고 하자. $\mathcal{A}$를 미분 등급
$\mathcal{O}_\mathcal{C}$-대수라 하고, $\mathcal{B}$를
미분 등급 $\mathcal{O}_\mathcal{D}$-대수라 하자. 다음과 같은
미분 등급 $\mathcal{O}_\mathcal{D}$-대수의 준동형
$\varphi : \mathcal{B} \to f_*\mathcal{A}$를 주어졌다고 하자.
그러면
$$
Rf_* : 
D(\mathcal{A}, \text{d})
\longrightarrow
D(\mathcal{B}, \text{d})
$$
는 다음 함자의 오른쪽 수반이다.
$$
Lf^* :
D(\mathcal{B}, \text{d})
\longrightarrow
D(\mathcal{A}, \text{d})
$$
\end{lemma}

\begin{proof}
이는 유도 범주, 보조정리
\ref{derived-lemma-pre-derived-adjoint-functors-general}과
보조정리 \ref{sdga-lemma-adjunction-push-pull-dg}에 따른다.
\end{proof}

\noindent
다음으로 절 \ref{sdga-section-derived-pullback}에서 고려한 상황에서
어떤 일이 일어나는지 논의한다.

\medskip\noindent
$(f, f^\sharp) : (\Sh(\mathcal{C}), \mathcal{O}_\mathcal{C})
\to (\Sh(\mathcal{D}), \mathcal{O}_\mathcal{D})$가
환 달린 토포스의 사상이라고 하자. $\mathcal{A}$를 미분
등급 $\mathcal{O}_\mathcal{C}$-대수라 하고, $\mathcal{B}$를
미분 등급 $\mathcal{O}_\mathcal{D}$-대수라 하자.
다음과 같은 사상
$$
\varphi : f^{-1}\mathcal{B} \to \mathcal{A}
$$
가 미분 등급 $f^{-1}\mathcal{O}_\mathcal{D}$-대수의 사상이라고 하자.
스칼라 제한과 스칼라 확장의 수반에 의해 이는
미분 등급 $\mathcal{O}_\mathcal{C}$-대수의 사상
$\varphi : f^*\mathcal{B} \to \mathcal{A}$와 같은 것이며,
동치로 $\varphi$를 다음 사상으로 볼 수 있다.
$$
\varphi : \mathcal{B} \to f_*\mathcal{A}
$$
미분 등급 $\mathcal{O}_\mathcal{D}$-대수의 사상으로서 말이다.
주의 \ref{sdga-remark-functoriality-dga}를 보라.

\medskip\noindent
위의 자료에 더하여 두 번째 미분 등급
$\mathcal{O}_\mathcal{C}$-대수 $\mathcal{A}'$와
미분 등급 $(\mathcal{A}, \mathcal{A}')$-쌍가군 $\mathcal{N}$을
택하자. 이 상황에서 다음 함자를 생각할 수 있다.
$$
\textit{Mod}(\mathcal{A}', \text{d})
\longrightarrow
\textit{Mod}(\mathcal{B}, \text{d}),\quad
\mathcal{M} \longmapsto
f_*\SheafHom^{dg}_{\mathcal{A}'}(\mathcal{N}, \mathcal{M})
$$
이를 다음 함자로 확장할 수 있음에 주의하자.
$$
\textit{Mod}^{dg}(\mathcal{A}', \text{d})
\longrightarrow
\textit{Mod}^{dg}(\mathcal{B}, \text{d}),\quad
\mathcal{M} \longmapsto
f_*\SheafHom^{dg}_{\mathcal{A}'}(\mathcal{N}, \mathcal{M})
$$
미분 등급 범주 사이의 함자이며, 이는 절
\ref{sdga-section-functoriality-dg}와 절 \ref{sdga-section-dg-bimodules}의
논의에 따른다. 따라서 형식적으로 다음 완전 함자도 얻는다.
\begin{equation}
\label{sdga-equation-pushforward}
K(\textit{Mod}(\mathcal{A}', \text{d}))
\longrightarrow
K(\textit{Mod}(\mathcal{B}, \text{d})),\quad
\mathcal{M} \longmapsto
f_*\SheafHom^{dg}_{\mathcal{A}'}(\mathcal{N}, \mathcal{M})
\end{equation}
삼각 범주 사이의 함자이다.

\begin{lemma}
\label{sdga-lemma-compose-pushforward-hom}
위의 상황에서 (\ref{sdga-equation-pushforward})의 오른쪽 유도 확장을
$RT : D(\mathcal{A}', \text{d}) \to
D(\mathcal{B}, \text{d})$로 나타내자. 그러면
$$
RT(\mathcal{M}) = Rf_* R\SheafHom(\mathcal{N}, \mathcal{M})
$$
$\mathcal{M}$에 대해 함자적으로 성립한다.
\end{lemma}

\begin{proof}
보조정리 \ref{sdga-lemma-tensor-hom-adjunction-dg}와
\ref{sdga-lemma-adjunction-push-pull-dg}에 의해 함자
(\ref{sdga-equation-pushforward})는 함자
(\ref{sdga-equation-pullback})의 오른쪽 수반이다. 유도 범주, 보조정리
\ref{derived-lemma-pre-derived-adjoint-functors-general}에 의해
함자 $RT$는 보조정리 \ref{sdga-lemma-derived-tensor-product}의
함자의 오른쪽 수반이며, 이 함자는 보조정리
\ref{sdga-lemma-compose-pullback-tensor}에 의해
$Lf^*(-) \otimes_\mathcal{A}^\mathbf{L} \mathcal{N}$과 같다.
보조정리 \ref{sdga-lemma-derived-adjoint-tensor-hom}과
\ref{sdga-lemma-derived-adjoint-push-pull}에 의해 함자
$Lf^*(-) \otimes_\mathcal{A}^\mathbf{L} \mathcal{N}$은
$Rf_* R\SheafHom(\mathcal{N}, -)$의 왼쪽 수반이다.
그러므로 수반의 유일성으로부터 결론을 얻는다.
\end{proof}

\begin{lemma}
\label{sdga-lemma-compose-pushforward}
$(f, f^\sharp) : (\Sh(\mathcal{C}), \mathcal{O})
\to (\Sh(\mathcal{C}'), \mathcal{O}')$와
$(g, g^\sharp) : (\Sh(\mathcal{C}'), \mathcal{O}')
\to (\Sh(\mathcal{C}''), \mathcal{O}'')$가
환 달린 토포스의 사상이라고 하자. $\mathcal{A}$, $\mathcal{A}'$ 및
$\mathcal{A}''$를 각각 미분 등급 $\mathcal{O}$-대수,
$\mathcal{O}'$-대수 및 $\mathcal{O}''$-대수라 하자. 다음과 같은
미분 등급 $\mathcal{O}'$-대수 및 $\mathcal{O}''$-대수의 준동형
$\varphi : \mathcal{A}' \to f_*\mathcal{A}$와
$\varphi' : \mathcal{A}'' \to g_*\mathcal{A}'$를 주어졌다고 하자.
그러면 $R(g \circ f)_* = Rg_* \circ Rf_* :
D(\mathcal{A}, \text{d}) \to D(\mathcal{A}'', \text{d})$이다.
\end{lemma}

\begin{proof}
이는 보조정리 \ref{sdga-lemma-compose-pullback}과
\ref{sdga-lemma-derived-adjoint-push-pull}, 그리고 수반의 유일성에 따른다.
\end{proof}

\begin{lemma}
\label{sdga-lemma-compose-hom}
환 달린 사이트 $(\mathcal{C}, \mathcal{O})$를 생각하자.
$\mathcal{A}$, $\mathcal{A}'$, $\mathcal{A}''$를 미분 등급
$\mathcal{O}$-대수라 하자. $\mathcal{N}$과 $\mathcal{N}'$를 각각
미분 등급 $(\mathcal{A}, \mathcal{A}')$-쌍가군과
$(\mathcal{A}', \mathcal{A}'')$-쌍가군이라 하자. 다음 정준 사상
$$
\mathcal{N} \otimes_{\mathcal{A}'}^\mathbf{L} \mathcal{N}'
\longrightarrow
\mathcal{N} \otimes_{\mathcal{A}'} \mathcal{N}'
$$
$D(\mathcal{A}'', \text{d})$에서 유사동형이라고 가정하자.
그러면
$$
R\SheafHom_{\mathcal{A}''}
(\mathcal{N} \otimes_{\mathcal{A}'} \mathcal{N}', -)
=
R\SheafHom_{\mathcal{A}'}(\mathcal{N},
R\SheafHom_{\mathcal{A}''}(\mathcal{N}', -))
$$
$D(\mathcal{A}'', \text{d}) \to D(\mathcal{A}, \text{d})$의 함자로서 성립한다.
\end{lemma}

\begin{proof}
이는 보조정리 \ref{sdga-lemma-compose-tensor}와
\ref{sdga-lemma-derived-adjoint-tensor-hom}, 그리고 수반의 유일성에 따른다.
\end{proof}

\begin{lemma}
\label{sdga-lemma-pushforward-agrees}
$(f, f^\sharp) : (\Sh(\mathcal{C}), \mathcal{O}_\mathcal{C})
\to (\Sh(\mathcal{D}), \mathcal{O}_\mathcal{D})$가
환 달린 토포스의 사상이라고 하자. $\mathcal{A}$를 미분
등급 $\mathcal{O}_\mathcal{C}$-대수라 하고, $\mathcal{B}$를
미분 등급 $\mathcal{O}_\mathcal{D}$-대수라 하자. 다음과 같은
미분 등급 $\mathcal{O}_\mathcal{D}$-대수의 준동형
$\varphi : \mathcal{B} \to f_*\mathcal{A}$를 주어졌다고 하자.
다음 도표는
$$
\xymatrix{
D(\mathcal{A}, \text{d}) \ar[d]_{Rf_*} \ar[rr]_{forget} & &
D(\mathcal{O}_\mathcal{C}) \ar[d]^{Rf_*} \\
D(\mathcal{B}, \text{d}) \ar[rr]^{forget} & &
D(\mathcal{O}_\mathcal{D})
}
$$
가환한다.
\end{lemma}

\begin{proof}
몇몇 범주를 동일시하는 것을 제외하면, 이 보조정리는
보조정리 \ref{sdga-lemma-compose-pushforward}로부터 즉시 따른다.

\medskip\noindent
$(\mathcal{O}_\mathcal{C}, 0)$을 $\mathcal{O}_\mathcal{C}$를
차수 $0$에 놓고 영 미분을 부여한 미분 등급
$\mathcal{O}_\mathcal{C}$-대수로 볼 수 있다. 다음이 성립함은 명백하다.
$$
\textit{Mod}(\mathcal{O}_\mathcal{C}, 0) =
\text{Comp}(\mathcal{O}_\mathcal{C})
\quad\text{and}\quad
D(\mathcal{O}_\mathcal{C}, 0) = D(\mathcal{O}_\mathcal{C})
$$
이 동일시에 의해 망각 함자
$\textit{Mod}(\mathcal{A}, \text{d}) \to
\text{Comp}(\mathcal{O}_\mathcal{C})$
는 절 \ref{sdga-section-functoriality-dg}에서 정의한
“직상” $\text{id}_{\mathcal{C}, *}$이다. 이는 항등 사상
$\text{id}_\mathcal{C} : (\mathcal{C}, \mathcal{O}_\mathcal{C}) \to
(\mathcal{C}, \mathcal{O}_\mathcal{C})$와 미분 등급
$\mathcal{O}_\mathcal{C}$-대수의 사상
$(\mathcal{O}_\mathcal{C}, 0) \to (\mathcal{A}, \text{d})$에
대응하는 환 달린 토포스의 사상에 의해 정의된다.
$\text{id}_{\mathcal{C}, *}$가 완전하므로 다음을 즉시 얻는다.
$$
R\text{id}_{\mathcal{C}, *} = forget :
D(\mathcal{A}, \text{d}) \longrightarrow
D(\mathcal{O}_\mathcal{C}, 0) = D(\mathcal{O}_\mathcal{C})
$$
정확히 같은 논의로
$$
R\text{id}_{\mathcal{D}, *} = forget :
D(\mathcal{B}, \text{d}) \longrightarrow
D(\mathcal{O}_\mathcal{D}, 0) = D(\mathcal{O}_\mathcal{D})
$$
더 나아가 사이트 위의 코호몰로지,
절 \ref{sites-cohomology-section-unbounded}에서의
$Rf_* : D(\mathcal{O}_\mathcal{C}) \to D(\mathcal{O}_\mathcal{D})$
의 구성은 정의
$Rf_* : D(\mathcal{O}_\mathcal{C}, 0) \to D(\mathcal{O}_\mathcal{D}, 0)$의
구성과 일치한다. 두 함자는 모두 기저 가군의 복합체에 대한
직상의 오른쪽 유도 확장으로 정의되기 때문이다.
정의 \ref{sdga-definition-pushforward}에 있다.
보조정리 \ref{sdga-lemma-compose-pushforward}에 의해
다음 두 함자는
$Rf_* \circ R\text{id}_{\mathcal{C}, *}$와
$R\text{id}_{\mathcal{D}, *} \circ Rf_*$는
$f_* \circ forget = forget \circ f_*$의 유도 함자이며,
따라서 수반의 유일성에 의해 서로 같다.
\end{proof}

\begin{lemma}
\label{sdga-lemma-cohomology-ext}
환 달린 사이트 $(\mathcal{C}, \mathcal{O})$를 생각하자.
$\mathcal{A}$를 미분 등급 $\mathcal{O}$-대수라 하고,
$\mathcal{M}$을 미분 등급 $\mathcal{A}$-가군이라 하자.
$n \in \mathbf{Z}$라 하자. 그러면
$$
H^n(\mathcal{C}, \mathcal{M}) =
\Hom_{D(\mathcal{A}, \text{d})}(\mathcal{A}, \mathcal{M}[n])
$$
여기서 좌변은 $\mathcal{M}$을 $\mathcal{O}$-가군의
복합체로 본 코호몰로지이다.
\end{lemma}

\begin{proof}
이 식을 증명하기 위해 다음을 관찰하자.
$$
R\Gamma(\mathcal{C}, \mathcal{M}) = \Gamma(\mathcal{C}, \mathcal{I})
$$
여기서 $\mathcal{M} \to \mathcal{I}$는 등급 단사이고 K-단사인
미분 등급 $\mathcal{A}$-가군 $\mathcal{I}$로 가는 유사동형이다
(보조정리 \ref{sdga-lemma-right-derived}와
\ref{sdga-lemma-pushforward-agrees}를 결합하라).
보조정리 \ref{sdga-lemma-hom-derived}에 의해 다음을 얻는다.
$$
\Hom_{D(\mathcal{A}, \text{d})}(\mathcal{A}, \mathcal{M}[n]) =
\Hom_{K(\textit{Mod}(\mathcal{A}, \text{d}))}(\mathcal{M}, \mathcal{I}[n]) =
H^0(\Gamma(\mathcal{C}, \mathcal{I}[n])) =
H^n(\Gamma(\mathcal{C}, \mathcal{I}))
$$
이 두 결과를 결합하면 원하는 등식을 얻는다.
\end{proof}








\section{유도 범주의 동치}
\label{sdga-section-equivalence}

\noindent
이 절은 미분 등급 대수, 절 \ref{dga-section-equivalence}에
대응한다.

\begin{lemma}
\label{sdga-lemma-qis-equivalence}
환 달린 사이트 $(\mathcal{C}, \mathcal{O})$를 생각하자.
$\varphi : \mathcal{A} \to \mathcal{B}$가 미분 등급
$\mathcal{O}$-대수의 준동형이고 코호몰로지 층에서
동형을 유도한다고 하자. 그러면
$$
D(\mathcal{A}, \text{d}) \longrightarrow D(\mathcal{B}, \text{d}), \quad
\mathcal{M}
\longmapsto
\mathcal{M} \otimes_\mathcal{A}^\mathbf{L} \mathcal{B}
$$
는 범주의 동치이다.
\end{lemma}

\begin{proof}
스칼라 제한 함자는
$$
\textit{Mod}^{dg}(\mathcal{B}, \text{d}) \to
\textit{Mod}^{dg}(\mathcal{A}, \text{d}),\quad
\mathcal{N} \mapsto res_\varphi \mathcal{N}
$$
이며 다음 함자의 오른쪽 수반이다.
$$
\textit{Mod}^{dg}(\mathcal{A}, \text{d}) \to
\textit{Mod}^{dg}(\mathcal{B}, \text{d}),\quad
\mathcal{M} \mapsto \mathcal{M} \otimes_\mathcal{A} \mathcal{B}
$$
절 \ref{sdga-section-dg-bimodules}를 보라. 제한 함자는
유사동형을 유사동형으로 보내므로, 자명하게 왼쪽 유도 확장을
(제한 함자 자체로 주어지는) 갖는다. 이는 유도 범주, 보조정리
\ref{derived-lemma-pre-derived-adjoint-functors-general}.
에 의해 $- \otimes_\mathcal{A}^\mathbf{L} \mathcal{B}$의
오른쪽 수반이 된다.
수반 사상
$$
\mathcal{M} \to
res_\varphi(\mathcal{M} \otimes_\mathcal{A}^\mathbf{L} \mathcal{B})
$$
는 $\mathcal{A} \to \mathcal{B}$가 (왼쪽) 미분 등급
$\mathcal{A}$-가군의 유사동형이라는 가정에 의해
$D(\mathcal{A}, \text{d})$에서 동형이다. 특히 이 보조정리의
함자는 완전 충실하다. 범주론, 보조정리
\ref{categories-lemma-adjoint-fully-faithful}를 보라.
제한 함자의 핵이
$D(\mathcal{B}, \text{d}) \to D(\mathcal{A}, \text{d})$
영이라는 것은 명백하다. 따라서 유도 범주, 보조정리
\ref{derived-lemma-fully-faithful-adjoint-kernel-zero}로부터
결론을 얻는다.
\end{proof}













\section{미분 등급 대수의 분해}
\label{sdga-section-resolution-dgas}

\noindent
이 절은 미분 등급 대수, 절 \ref{dga-section-resolution-dgas}에
대응한다.

\medskip\noindent
환 달린 사이트 $(\mathcal{C}, \mathcal{O})$를 생각하자.
주의 \ref{sdga-remark-sheaf-graded-sets}에서와 같이
$\mathcal{C}$ 위의 등급 집합의 층 $\mathcal{S}$를 생각하자.
$r$회 자기곱
$\mathcal{S} \times \ldots \times \mathcal{S}$
을 다음 규칙을 갖는 등급 집합의 층으로 보자.
$\deg(s_1 \cdot \ldots \cdot s_r) = \sum \deg(s_i)$.
여기서 국소 단면 $s_i \in \mathcal{S}(U)$,
$i = 1, \ldots, r$가 주어졌을 때
$s_1 \cdot \ldots \cdot s_r$를
$\mathcal{S} \times \ldots \times \mathcal{S}$의
$U$ 위 대응하는 단면을 나타내는 데 사용한다.
$\mathcal{O}\langle \mathcal{S} \rangle$를
$\mathcal{S}$ 위 자유 등급 $\mathcal{O}$-대수라 하자.
더 정확히는
$$
\mathcal{O}\langle \mathcal{S} \rangle =
\mathcal{O} \oplus
\bigoplus\nolimits_{r \geq 1}
\mathcal{O}[\mathcal{S} \times \ldots \times \mathcal{S}]
$$
주의 \ref{sdga-remark-sheaf-graded-sets}와 같은 표기를 사용한다.
이는 다음 이어붙이기에 의해 등급 $\mathcal{O}$-대수의 층이 된다.
$$
(s_1 \cdot \ldots \cdot s_r)  (s'_1 \cdot \ldots \cdot s'_{r'}) =
s_1 \cdot \ldots s_r \cdot s'_1 \cdot \ldots \cdot s'_{r'}
$$
$\mathcal{O}\langle \mathcal{S} \rangle$에
$\mathcal{S}$의 모든 국소 단면 $s$에 대해
$\text{d}(s) = 0$으로 놓고 라이프니츠 법칙을 사용해 유일하게
확장함으로써 미분을 부여할 수 있다. 다만 다른 미분들도
고려하는 것이 중요하다.

\medskip\noindent
실제로 다음과 같은 종류의 계가 주어졌다고 하자.
\begin{enumerate}
\item $i = 0, 1, 2, \ldots$에 대해 등급 집합의 층 $\mathcal{S}_i$,
\item $i = 0, 1, 2, \ldots$에 대해 다음 사상
$$
\delta_{i + 1} : \mathcal{S}_{i + 1}
\longrightarrow
\mathcal{A}_i =
\mathcal{O}\langle \mathcal{S}_0 \amalg \ldots \amalg \mathcal{S}_i\rangle
$$
의 상이 공역 위 귀납적으로 정의된 미분의 핵에 포함되는,
차수 $1$의 등급 집합의 층의 사상.
\end{enumerate}
더 정확히는 먼저
$\mathcal{A}_0 = \mathcal{O}\langle \mathcal{S}_0 \rangle$
로 놓고, $\mathcal{S}$의 임의의 국소 단면 $s$에 대해
$\text{d}(s) = 0$이 되도록 라이프니츠 법칙을 만족하는 유일한
미분을 부여한다. 귀납적으로 $\mathcal{A}_i$ 위의
미분 $\text{d}$가 주어졌다고 하자. 그러면 이를 라이프니츠 법칙을
만족하고 다음을 갖는 $\mathcal{A}_{i + 1}$ 위의 유일한
미분으로 확장한다.
$$
\text{d}(s) = \delta(s)
$$
여기서 $s$가
$\mathcal{S}_0 \amalg \ldots \amalg \mathcal{S}_{i + 1}$의
직합 성분 $\mathcal{S}_j$에 있으면 $\delta(s) = \delta_j(s)$로 한다.
이는 정확히 $\delta(s)$가 귀납적으로 정의된 미분의 핵에
속하기 때문에 가능하다.

\begin{lemma}
\label{sdga-lemma-special-good}
위 상황의 미분 등급 $\mathcal{O}$-대수
$$
\mathcal{A} = \colim \mathcal{A}_i
$$
는 다음 성질을 갖는다. 임의의 사상
$(f, f^\sharp) : (\Sh(\mathcal{C}'), \mathcal{O}')
\to (\Sh(\mathcal{C}), \mathcal{O})$
환 달린 토포스의 사상에 대해 당김 $f^*\mathcal{A}$는
등급 $\mathcal{O}'$-가군으로서 평탄하고
$\mathcal{O}'$-가군의 복합체로서 K-평탄하다.
\end{lemma}

\begin{proof}
 $f^*\mathcal{A} = \colim f^*\mathcal{A}_i$이고 다음임을 관찰하자.
$$
f^*\mathcal{A}_i = \mathcal{O}'\langle
f^{-1}\mathcal{S}_0 \amalg \ldots \amalg f^{-1}\mathcal{S}_i\rangle
$$
의 미분은 $f^{-1}\delta_{i + 1}$를 사용한 위의 귀납 절차로
주어진다. 따라서 $\mathcal{A}$가 등급 $\mathcal{O}$-가군으로서
평탄하고 $\mathcal{O}$-가군의 복합체로서 K-평탄함을 보이면 충분하다.
이를 위해서는 각 $\mathcal{A}_i$가 등급 $\mathcal{O}$-가군으로서
평탄하고 $\mathcal{O}$-가군의 복합체로서 K-평탄함을 보이면 충분하다.
보조정리 \ref{sdga-lemma-good-direct-sum}과 비교하라.


\medskip\noindent
 $i \geq 1$에 대해
$\mathcal{S} = \mathcal{S}_0 \amalg \ldots \amalg \mathcal{S}_i$로
쓰자. 그러면 등급 $\mathcal{O}$-대수로서
$\mathcal{A}_i = \mathcal{O}\langle \mathcal{S} \rangle$이다.
이 대수의 미분 등급 $\mathcal{O}$-부분가군에 의한 여과를 구성하자.

\medskip\noindent
$W = \mathbf{Z}_{\geq 0}^{i + 1}$를 사전식 순서와 함께 생각하자.
즉 $W$의 $w = (w_0, \ldots w_i)$와
$w' = (w'_0, \ldots, w'_i)$가 주어졌을 때 다음과 같이 말한다.
$$
w > w' \Leftrightarrow
\exists j,\ 0 \leq j \leq i :
w_i = w'_i,\ w_{i - 1} = w'_{i - 1},\ \ldots ,
\ w_{j + 1} = w'_{j + 1},\ w_j > w'_j
$$
이와 같이 계속한다. 단면
$s = s_1 \cdot \ldots \cdot s_r$이
$\mathcal{S} \times \ldots \times \mathcal{S}$
의 단면으로서 $U$ 위에 주어졌다고 하자. 어떤 유일한
$0 \leq j_a \leq i$에 대해 $s_a \in \mathcal{S}_{j_a}(U)$이면
{\it $s$의 가중치가 정의되었다}고 한다. 이 경우 가중치를
$$
w(s) = (w_0(s), \ldots, w_i(s)) \in W,\quad
w_j(s) = |\{a \mid j_a = j\}|
$$
 $\mathcal{S} \times \ldots \times \mathcal{S}$의 임의의 단면의
가중치는 국소적으로 정의된다. 다음과 같은 서로소 합 분해를
얻는 것은 쉽게 확인할 수 있다.
$$
\mathcal{S} \times \ldots \times \mathcal{S} =
\coprod\nolimits_{w \in W} \left(
\mathcal{S} \times \ldots \times \mathcal{S}\right)_w
$$
이는 주어진 가중치를 갖는 단면들의 부분층들로의 분해이다.
물론 주어진 $r$에 대해서는
$\sum_{0 \leq j \leq i} w_j = r$인 $w \in W$만 나타난다.
이에 대응하여 다음 분해를 얻는다.
$$
\mathcal{A}_i = \mathcal{O} \oplus
\bigoplus\nolimits_{r \geq 1}
\bigoplus\nolimits_{w \in W}
\mathcal{O}[\left(\mathcal{S} \times \ldots \times \mathcal{S}\right)_w]
$$
증명의 나머지는 다음 자명한 관찰에 의존한다.
$r$, $w$ 및
$\left(\mathcal{S} \times \ldots \times \mathcal{S}\right)_w$의
국소 단면 $s = s_1 \cdot \ldots \cdot s_r$가 주어지면
다음이 성립한다.
$$
\text{d}(s) \text{ is a local section of } \mathcal{O} \oplus
\bigoplus\nolimits_{r' \geq 1}
\bigoplus\nolimits_{w' \in W,\ w' < w}
\mathcal{O}[\left(\mathcal{S} \times \ldots \times \mathcal{S}\right)_{w'}]
$$
그 이유는 다음 식들 각각에서
$$
(-1)^{\deg(s_1) + \ldots + \deg(s_{a - 1})}
s_1 \cdot \ldots s_{a - 1} \cdot \delta(s_a) \cdot
s_{a + 1} \cdot \ldots \cdot s_r
$$
그 합이 원소 $\text{d}(s)$를 주며, 원소 $\delta(s_a)$는
국소적으로 원소
$s'_1 \cdot \ldots \cdot s'_{r'}$들의 $\mathcal{O}$-선형 결합이다.
여기서 어떤 $0 \leq j'_{a'} < j_a$에 대해
$s'_{a'}$는 $\mathcal{S}_{j'_a}$에 속하고, $j_a$는
$s_a$가 $\mathcal{S}_{j_a}$의 단면이 되도록 택한 것이다.

\medskip\noindent
이는 다음을 뜻한다. $w \in W$에 대해
$$
F_w \mathcal{A}_i = \mathcal{O} \oplus \bigoplus\nolimits_{r \geq 1}
\bigoplus\nolimits_{w' \in W,\ w' \leq w}
\mathcal{O}[\left(\mathcal{S} \times \ldots \times \mathcal{S}\right)_{w'}]
$$
위 관찰에 의해 이는 미분 등급 $\mathcal{O}$-부분가군이다.
다음의 허용 가능한 짧은 완전열을 얻는다.
$$
0 \to \colim_{w' < w} F_{w'}\mathcal{A}_i \to
F_w\mathcal{A}_i \to
\bigoplus\nolimits_{r \geq 1}
\mathcal{O}[\left(\mathcal{S} \times \ldots \times \mathcal{S}\right)_w]
\to 0
$$
미분 등급 $\mathcal{A}$-가군의 완전열이며, 오른쪽 항의 미분은
영이다.

\medskip\noindent
이제 순서 집합 $W$에 대한 초한 귀납법으로 증명을 마친다.
미분 등급 복합체 $F_0\mathcal{A}_0$는 직합 성분 $\mathcal{O}$이고,
이는 K-평탄하며 등급 평탄하다.
$w \in W$에 대해 $w' < w$인 경우
$F_{w'}\mathcal{A}_i$에 대한 결과가 성립한다고 하자.
그러면 보조정리 \ref{sdga-lemma-good-direct-sum},
\ref{sdga-lemma-good-admissible-ses}, 및
\ref{sdga-lemma-free-graded-module-good}에 의해 $w$에서도 결과를 얻는다.
마지막으로
$\mathcal{A}_i = \colim_{w \in W} F_w\mathcal{A}_i$이므로 결론을 얻는다.
\end{proof}

\begin{lemma}
\label{sdga-lemma-good-dga}
$(\mathcal{C}, \mathcal{O})$를 환 달린 사이트라고 하자.
$(\mathcal{B}, \text{d})$를 미분 등급 $\mathcal{O}$-대수라고 하자.
다음 조건을 만족하는 미분 등급 $\mathcal{O}$-대수의 유사동형이 존재한다.
$(\mathcal{A}, \text{d}) \to (\mathcal{B}, \text{d})$이고,
$\mathcal{A}$는 $\mathcal{O}$-가군의 복합체로서 등급 평탄하며 K-평탄하다.
또한 환 달린 위상공간의 임의의 사상에 의한 당김 뒤에도
같은 성질이 성립한다.
\end{lemma}

\begin{proof}
증명은 보조정리 \ref{sdga-lemma-resolve}의 첫 번째 증명과 정확히 같지만,
이번에는 자유 등급 가군 대신 자유 등급 대수를 사용한다.

\medskip\noindent
보조정리 \ref{sdga-lemma-special-good}에서와 같이 다음을 구성하여
$\mathcal{A} = \colim \mathcal{A}_i$를 만들 것이다.
$$
\mathcal{A}_0 \to \mathcal{A}_1 \to \mathcal{A}_2 \to \ldots \to \mathcal{B}
$$
 $\mathcal{S}_0$를 등급 집합의 층
(비고 \ref{sdga-remark-sheaf-graded-sets})으로 두고,
그 $n$차 부분이 $\Ker(\text{d}_\mathcal{B}^n)$이게 하자.
다음 미분 등급 가군 사상을 생각하자.
$$
\mathcal{A}_0 = \mathcal{O}\langle \mathcal{S}_0 \rangle
\longrightarrow
\mathcal{B}
$$
이 사상은 $\mathcal{S}_0$의 국소 단면 $s$를
이에 대응하는 $\mathcal{A}^{\deg(s)}$의 국소 단면으로 보낸다.
(이는 미분의 핵에 속하므로 실제로 미분 등급 대수의 사상이다.)
구성에 의해 코호몰로지 층 사이의 유도 사상
$H^n(\mathcal{A}_0) \to H^n(\mathcal{B})$는 전사이고,
따라서 이는 모든 $i$에 대해 계속 성립한다.

\medskip\noindent
구성의 귀납 단계. $\mathcal{A}_i \to \mathcal{B}$가 주어졌다고 하자.
차수 $n$ 부분이 다음인 등급 집합의 층을 $\mathcal{S}_{i + 1}$로 쓰자.
$$
\Ker(\text{d}_{\mathcal{A}_i}^{n + 1})
\times_{\mathcal{B}^{n + 1}, \text{d}}
\mathcal{B}^n
$$
이는 다음의 표준 사상을 갖는다.
$$
\delta_{i + 1} : \mathcal{S}_{i + 1} \longrightarrow
\mathcal{A}_i
$$
그 상은 구성에 의해 $\text{d}_{\mathcal{A}_i}$의 핵에 포함된다.
따라서 $\mathcal{A}_{i + 1} = \mathcal{O}\langle
\mathcal{S}_0 \amalg \ldots \mathcal{S}_{i + 1}\rangle$는
이 절의 앞부분에서 논의한 대로 $\mathcal{A}_i$의 미분을
확장하는 미분을 갖는다. $\mathcal{A}_{i + 1}$에서
$\mathcal{B}$로 가는 사상은 $\mathcal{A}_i$에서 주어진 사상으로
제한되고, $\mathcal{S}_{i + 1}$의 국소 단면 $s = (a, b)$를
$\mathcal{B}$의 $b$로 보내는 유일한 등급 대수 사상이다.
이는 $\text{d}(b)$가 $\mathcal{B}$에서 $a$의 상이기 때문에
정확히 미분과 양립한다.

\medskip\noindent
사상 $\mathcal{A} \to \mathcal{B}$는 유사동형이다.
실제로 $H^n(\mathcal{A}) = \colim H^n(\mathcal{A}_i)$이고,
각 $i$에 대해 사상 $H^n(\mathcal{A}_i) \to H^n(\mathcal{B})$는
전사이며 그 핵은 구성에 의해
$H^n(\mathcal{A}_i) \to H^n(\mathcal{A}_{i + 1})$에 의해 소멸한다.
마지막으로 $\mathcal{A}$의 평탄성 조건은
보조정리 \ref{sdga-lemma-special-good}에서 보였다.
\end{proof}






\section{기타 사항}
\label{sdga-section-misc}

\noindent
$(f, f^\sharp) : (\Sh(\mathcal{C}), \mathcal{O})
\to (\Sh(\mathcal{C}'), \mathcal{O}')$를 환 달린 토포스의 사상이라고 하자.
$\mathcal{A}$를 미분 등급 $\mathcal{O}$-대수의 층이라고 하자.
다음 합성을 사용한다.\footnote{더 정확히는 다음과 같이 쓰는 편이 낫다.
$F(\mathcal{A}) \otimes_\mathcal{O}^\mathbf{L} F(\mathcal{A})
\to F(\mathcal{A} \otimes_\mathcal{O} \mathcal{A}) \to F(\mathcal{A})$
여기서 $F$는 $\mathcal{O}$-가군 복합체로 가는 망각 함자를 나타낸다.
또한 $\mathcal{A} \otimes_\mathcal{O} \mathcal{A}$는
절 \ref{sdga-section-tensor-product-dg}의 텐서곱을 뜻하므로
$F(\mathcal{A} \otimes_\mathcal{O} \mathcal{A}) =
\text{Tot}(F(\mathcal{A}) \otimes_\mathcal{O} F(\mathcal{A}))$.
열의 첫 번째 화살표는 두 $\mathcal{O}$-가군 복합체의 유도 텐서곱에서
두 $\mathcal{O}$-가군 복합체의 통상적인 텐서곱으로 가는 표준 사상이다.}
$$
\mathcal{A} \otimes_\mathcal{O}^\mathbf{L} \mathcal{A}
\longrightarrow
\mathcal{A} \otimes_\mathcal{O} \mathcal{A}
\longrightarrow
\mathcal{A}
$$
그리고 상대 컵곱 (Cohomology on Sites, 비고
\ref{sites-cohomology-remark-cup-product} 및
절 \ref{sites-cohomology-section-cup-product} 참조)을 사용하면
곱셈을 얻는다.\footnote{여기와 이하에서
$Rf_* : D(\mathcal{O}) \to D(\mathcal{O}')$는 Cohomology on Sites, 절
\ref{sites-cohomology-section-unbounded} 이하에서 연구한 유도 함자이다.}
$$
\mu :
Rf_*\mathcal{A} \otimes_{\mathcal{O}'}^\mathbf{L} Rf_*\mathcal{A}
\longrightarrow
Rf_*\mathcal{A}
$$
 $D(\mathcal{O}')$에서. 이 곱셈은 다음 도표가
$$
\xymatrix{
Rf_*\mathcal{A} \otimes_{\mathcal{O}'}^\mathbf{L} Rf_*\mathcal{A}
\otimes_{\mathcal{O}'}^\mathbf{L} Rf_*\mathcal{A}
\ar[rr]_-{\mu \otimes 1} \ar[d]_{1 \otimes \mu} & &
Rf_*\mathcal{A} \otimes_{\mathcal{O}'}^\mathbf{L} Rf_*\mathcal{A}
\ar[d]^\mu \\
Rf_*\mathcal{A} \otimes_{\mathcal{O}'}^\mathbf{L} Rf_*\mathcal{A}
\ar[rr]^-\mu & &
Rf_*\mathcal{A}
}
$$
 $D(\mathcal{O}')$에서 가환한다. 이는
Cohomology on Sites, 보조정리
\ref{sites-cohomology-lemma-cup-product-associative}에서 따른다.
정확히 같은 방법으로, 오른쪽 미분 등급 $\mathcal{A}$-가군
$\mathcal{M}$이 주어지면 곱셈
$$
\mu_\mathcal{M} :
Rf_*\mathcal{M} \otimes_{\mathcal{O}'}^\mathbf{L} Rf_*\mathcal{A}
\longrightarrow
Rf_*\mathcal{M}
$$
 $D(\mathcal{O}')$에서 얻는다. 이 곱셈은 위의 $\mu$와 양립하며,
다음 도표가
$$
\xymatrix{
Rf_*\mathcal{M} \otimes_{\mathcal{O}'}^\mathbf{L} Rf_*\mathcal{A}
\otimes_{\mathcal{O}'}^\mathbf{L} Rf_*\mathcal{A}
\ar[rr]_-{\mu_\mathcal{M} \otimes 1} \ar[d]_{1 \otimes \mu} & &
Rf_*\mathcal{M} \otimes_{\mathcal{O}'}^\mathbf{L} Rf_*\mathcal{A}
\ar[d]^{\mu_\mathcal{M}} \\
Rf_*\mathcal{M} \otimes_{\mathcal{O}'}^\mathbf{L} Rf_*\mathcal{A}
\ar[rr]^-{\mu_\mathcal{M}} & &
Rf_*\mathcal{M}
}
$$
 $D(\mathcal{O}')$에서 가환하며, 이번에도
Cohomology on Sites, 보조정리
\ref{sites-cohomology-lemma-cup-product-associative}에서 따른다.

\medskip\noindent
위의 특별한 예로 $f$를 점 토포스 $\Sh(pt)$로 가는 사상으로
취하는 경우를 들 수 있다. 이때 $\mu$는 컵곱 사상
$$
R\Gamma(\mathcal{C}, \mathcal{A})
\otimes_{\Gamma(\mathcal{C}, \mathcal{O})}^\mathbf{L}
R\Gamma(\mathcal{C}, \mathcal{A})
\longrightarrow
R\Gamma(\mathcal{C}, \mathcal{A}),
\quad \eta \otimes \theta \mapsto \eta \cup \theta
$$
마찬가지로 $\mu_\mathcal{M}$은 컵곱 사상
$$
R\Gamma(\mathcal{C}, \mathcal{M})
\otimes_{\Gamma(\mathcal{C}, \mathcal{O})}^\mathbf{L}
R\Gamma(\mathcal{C}, \mathcal{A})
\longrightarrow
R\Gamma(\mathcal{C}, \mathcal{M}),
\quad \eta \otimes \theta \mapsto \eta \cup \theta
$$
일반적으로 다음 식별을 통해
$$
R\Gamma(\mathcal{C}, \mathcal{A}) =
R\Gamma(\mathcal{C}', Rf_*\mathcal{A})
\quad\text{and}\quad
R\Gamma(\mathcal{C}, \mathcal{M}) =
R\Gamma(\mathcal{C}', Rf_*\mathcal{M})
$$
Cohomology on Sites, 비고 \ref{sites-cohomology-remark-before-Leray}에 의해
$\mu_\mathcal{M}$은 코호몰로지에서 컵곱을 유도한다.
이를 보려면 Cohomology on Sites, 보조정리
\ref{sites-cohomology-lemma-compose-cup-product}를 사용하자.
여기서 두 번째 토포스 사상은 위와 같이 $\Sh(\mathcal{C}')$에서
점 토포스로 가는 사상이다.

\medskip\noindent
 $\mathcal{M}_1 \to \mathcal{M}_2$가 오른쪽 미분 등급
$\mathcal{A}$-가군의 준동형이면 다음 도표는
$$
\xymatrix{
Rf_*\mathcal{M}_1 \otimes_{\mathcal{O}'}^\mathbf{L} Rf_*\mathcal{A}
\ar[rr]_-{\mu_{\mathcal{M}_1}} \ar[d] & &
Rf_*\mathcal{M}_1 \ar[d] \\
Rf_*\mathcal{M}_2 \otimes_{\mathcal{O}'}^\mathbf{L} Rf_*\mathcal{A}
\ar[rr]^-{\mu_{\mathcal{M}_2}} & &
Rf_*\mathcal{M}_2
}
$$
 $D(\mathcal{O}')$에서 가환한다. 이는 상대 컵곱이
함자적이라는 사실에서 따른다. 다음 짧은 완전열이 주어졌다고 하자.
$$
0 \to \mathcal{M}_1 \xrightarrow{a} \mathcal{M}_2 \to \mathcal{M}_3 \to 0
$$
오른쪽 미분 등급 $\mathcal{A}$-가군의 완전열이다. 그러면 다음 도표가
$$
\xymatrix{
Rf_*\mathcal{M}_3 \otimes_{\mathcal{O}'}^\mathbf{L} Rf_*\mathcal{A}
\ar[rr]_-{\mu_{\mathcal{M}_3}} \ar[d]_{Rf_*\delta \otimes \text{id}} & & 
Rf_*\mathcal{M}_3 \ar[d]^{Rf_*\delta} \\
Rf_*\mathcal{M}_1[1] \otimes_{\mathcal{O}'}^\mathbf{L} Rf_*\mathcal{A}
\ar[rr]^-{\mu_{\mathcal{M}_1[1]}} & &
Rf_*\mathcal{M}_1[1]
}
$$
 $D(\mathcal{O}')$에서 가환한다. 여기서
$\delta : \mathcal{M}_3 \to \mathcal{M}_1[1]$는 주어진 짧은 완전열에서
오는 $D(\mathcal{O})$의 사상이다 (Derived Categories, 절
\ref{derived-section-canonical-delta-functor} 참조).
이 완전열이 등급 오른쪽 $\mathcal{A}$-가군의 완전열로서 분할되면
이는 자명하다. 이 경우 $\delta$를 오른쪽 $\mathcal{A}$-가군의 사상으로
나타낼 수 있으므로 위 논의를 적용할 수 있다.
일반적인 경우에는 $a$의 원뿔과 다음 도표를 사용한다.
$$
\xymatrix{
\mathcal{M}_1 \ar[r]_a \ar[d] &
\mathcal{M}_2 \ar[r]_i \ar[d] &
C(a) \ar[r]_{-p} \ar[d]^q &
\mathcal{M}_1[1] \ar[d] \\
\mathcal{M}_1 \ar[r] &
\mathcal{M}_2 \ar[r] &
\mathcal{M}_3 \ar[r]^\delta &
\mathcal{M}_1[1]
}
$$
오른쪽 정사각형은 Derived Categories, 보조정리
\ref{derived-lemma-derived-canonical-delta-functor}에서의 $\delta$의
정의에 의해 $D(\mathcal{O})$에서 가환한다.
이제 원뿔 $C(a)$는 오른쪽 미분 등급 $\mathcal{A}$-가군의 구조를 가지며,
$i$, $p$, $q$는 오른쪽 미분 등급 $\mathcal{A}$-가군의 준동형이다
(정의 \ref{sdga-definition-cone} 참조).
따라서 위 결과에 의해 $q$와 $-p$에 대응하는 도표들이 가환한다.
$q$가 $D(\mathcal{O})$에서 동형이므로, 원하는 대로 $\delta$에 대해서도
같은 결론을 얻는다.

\medskip\noindent
위 상황에서 오른쪽 미분 등급 $\mathcal{A}$-가군
$\mathcal{M}$과 다음 원소가 주어졌다고 하자.
$$
\xi \in H^n(\mathcal{C}, \mathcal{M})
$$
즉 $\xi$는 $\mathcal{O}$-가군의 복합체로 본
$\mathcal{M}$의 코호몰로지에서 차수 $n$인 코호몰로지류이다.
보조정리 \ref{sdga-lemma-cohomology-ext}에 의해 다음 사상들을 구성할 수 있다.
$$
x : \mathcal{A} \rightarrow \mathcal{M}'[n]
\quad\text{and}\quad
s : \mathcal{M} \to \mathcal{M}'
$$
오른쪽 미분 등급 $\mathcal{A}$-가군의 사상으로서 $s$는
유사동형이고, 유도 범주 $D(\mathcal{A}, \text{d})$에서의 사상
$s[n]^{-1} \circ x$를 통해
$1 \in H^0(\mathcal{C}, \mathcal{A})$의 상이 $\xi$가 되게 한다.
따라서 유도 범주 $D(\mathcal{O})$에서도 마찬가지이다.
그러므로 이에 대응하는 사상
$$
\xi' = (s[n])^{-1} \circ x : \mathcal{A} \longrightarrow \mathcal{M}[n]
$$
 $D(\mathcal{O})$에서 다음 두 성질로 유일하게 특징지어진다.
\begin{enumerate}
\item $\xi'$는 $D(\mathcal{A}, \text{d})$의 사상으로 올릴 수 있다.
\item $H^0(\mathcal{C}, \mathcal{M}[n]) = H^n(\mathcal{C}, \mathcal{M})$에서
$\xi = \xi'(1)$이다.
\end{enumerate}
앞에서 논의한 상대 컵곱과 $x$, $s$의 양립성을 사용하면,
모든\footnote{예를 들어 항등 사상.} 환 달린 토포스의 사상
$(f, f^\sharp) : (\Sh(\mathcal{C}), \mathcal{O})
\to (\Sh(\mathcal{C}'), \mathcal{O}')$에 대해 유도된 직상
$$
Rf_*\xi' : Rf_*\mathcal{A}  \longrightarrow Rf_*\mathcal{M}[n]
$$
 $\xi'$는 위에서 구성한 $\mu$ 및 $\mu_{\mathcal{M}[n]}$와
다음 의미에서 양립한다. 즉 도표가
$$
\xymatrix{
Rf_*\mathcal{A} \otimes_{\mathcal{O}'}^\mathbf{L} Rf_*\mathcal{A}
\ar[rr]_-\mu \ar[d]_{Rf_*\xi' \otimes \text{id}}  & &
Rf_*\mathcal{A} \ar[d]^{Rf_*\xi'} \\
Rf_*\mathcal{M}[n] \otimes_{\mathcal{O}'}^\mathbf{L} Rf_*\mathcal{A}
\ar[rr]^-{\mu_{\mathcal{M}[n]}} & &
Rf_*\mathcal{M}[n]
}
$$
 $D(\mathcal{O}')$에서 가환한다. 점 토포스로 가는 사상에 대해
이 양립성을 사용하면 특히 다음을 얻는다.
$$
\xymatrix{
R\Gamma(\mathcal{C}, \mathcal{A})
\otimes_{\Gamma(\mathcal{C}, \mathcal{O})}^\mathbf{L}
R\Gamma(\mathcal{C}, \mathcal{A})
\ar[d]_{\xi' \otimes \text{id}} \ar[r] &
R\Gamma(\mathcal{C}, \mathcal{A}) \ar[d]^{\xi'} \\
R\Gamma(\mathcal{C}, \mathcal{M}[n])
\otimes_{\Gamma(\mathcal{C}, \mathcal{O})}^\mathbf{L}
R\Gamma(\mathcal{C}, \mathcal{A})
\ar[r] &
R\Gamma(\mathcal{C}, \mathcal{M}[n])
}
$$
가환한다. $\xi'(1) = \xi$와 결합하면 코호몰로지에 유도되는 사상
$$
\xi' : R\Gamma(\mathcal{C}, \mathcal{A}) \to
R\Gamma(\mathcal{C}, \mathcal{M}[n]), \quad \eta \mapsto \xi \cup \eta
$$
은 표시된 대로 $\xi$에 의한 왼쪽 컵곱으로 주어진다.














\section{범주 위의 미분 등급 가군}
\label{sdga-section-modules-cohomology}

\noindent
이 절은 Cohomology on Sites의
절 \ref{sites-cohomology-section-modules-cohomology}의 계속이다.

\medskip\noindent
$\mathcal{C}$를 범주라고 하자. $\mathcal{C}$를 혼돈 위상을 갖는
사이트로 생각한다. $\mathcal{O}$를 $\mathcal{C}$ 위의 환의 층이라고 하자.
$(\mathcal{A}, \text{d})$를 미분 등급 $\mathcal{O}$-대수의 층이라고 하자.
즉 $\mathcal{O}$는 범주 $\mathcal{C}$ 위의 환의 프리셰프이고,
$(\mathcal{A}, \text{d})$는 $\mathcal{C}$ 위의 미분 등급
$\mathcal{O}$-대수의 프리셰프이다 (Categories,
정의 \ref{categories-definition-presheaf} 참조).

\begin{definition}
\label{sdga-definition-cartesian}
위 상황에서
{\it $\mathit{QC}(\mathcal{A}, \text{d})$}를
$D(\mathcal{A}, \text{d})$의 다음 성질을 갖는 대상 $M$들로 이루어진
충만 부분범주로 정의한다. 즉 $\mathcal{C}$의 모든 $U \to V$에 대해
표준 사상
$$
R\Gamma(V, M) \otimes_{\mathcal{A}(V)}^\mathbf{L} \mathcal{A}(U)
\longrightarrow
R\Gamma(U, M)
$$
은 $D(\mathcal{A}(U), \text{d})$에서 동형이다.
\end{definition}

\begin{lemma}
\label{sdga-lemma-cartesian-good-sub}
위 상황에서 부분범주 $\mathit{QC}(\mathcal{A}, \text{d})$는
$D(\mathcal{A}, \text{d})$의 엄밀히 충만하고 포화된 삼각 부분범주이며,
임의의 직합에 대해 보존된다.
\end{lemma}

\begin{proof}
 $\mathcal{C}$의 대상 $U$를 택하자. $\mathcal{C}$의 위상이
혼돈 위상이므로 함자 $\mathcal{F} \mapsto \mathcal{F}(U)$는
정확하고 직합과 가환한다. 따라서 정확한 함자
$M \mapsto R\Gamma(U, M)$는 $K$를 임의의 미분 등급
$\mathcal{A}$-가군 $\mathcal{M}$으로 나타내고 $\mathcal{M}(U)$를
취하여 계산된다. 그러므로 $R\Gamma(U, -)$는 직합과 가환한다
(보조정리 \ref{sdga-lemma-derived-products} 참조).
마찬가지로 $\mathcal{C}$의 사상 $U \to V$가 주어지면
유도 텐서곱 함자
$- \otimes_{\mathcal{O}(A)}^\mathbf{L} \mathcal{A}(U) :
D(\mathcal{A}(V)) \to D(\mathcal{A}(U))$
는 정확하고 직합과 가환한다. 보조정리는 이 관찰들로부터
곧바로 따른다. 세부사항은 생략한다.
\end{proof}

\begin{remark}
\label{sdga-remark-cartesian-brown}
앞과 같이 $\mathcal{C}$를 혼돈 위상을 갖는 사이트로 본 범주,
$\mathcal{O}$를 $\mathcal{C}$ 위의 환의 층,
$(\mathcal{A}, \text{d})$를 미분 등급 $\mathcal{O}$-대수의 층이라고 하자.
그러면 Cohomology on Sites, 명제
\ref{sites-cohomology-proposition-cartesian-brown}의 유사한 명제가
$\mathit{QC}(\mathcal{A}, \text{d})$에 대해 거의 정확히 같은 증명으로
성립한다.
\begin{enumerate}
\item 직합을 곱으로 보내는 임의의 반변 코호몰로지 함자
$H : \mathit{QC}(\mathcal{A}, \text{d}) \to \textit{Ab}$
는 표현 가능하다.
\item 임의의 정확한 함자
$F : \mathit{QC}(\mathcal{A}, \text{d}) \to \mathcal{D}$
삼각 범주 사이의 함자로서 직합을 직합으로 보내는 것은
정확한 오른쪽 수반을 갖는다.
\item 포함 함자
$\mathit{QC}(\mathcal{A}, \text{d}) \to D(\mathcal{A}, \text{d})$ has
는 정확한 오른쪽 수반을 갖는다.
\end{enumerate}
필요해지는 경우에는 여기서 이를 정확히 정식화하고 증명할 것이다.
\end{remark}

\noindent
 $u : \mathcal{C}' \to \mathcal{C}$를 범주 사이의 함자라고 하자.
$\mathcal{C}$와 $\mathcal{C}'$를 혼돈 위상을 갖는 사이트로 보면
$u$는 연속이고 쌍연속인 함자이다. 따라서 토포스의 사상
$g : \Sh(\mathcal{C}') \to \Sh(\mathcal{C})$를 얻는다
(Sites, 보조정리 \ref{sites-lemma-cocontinuous-morphism-topoi} 참조).
또한 $\mathcal{C}$ 위의 환의 층 $\mathcal{O}$와
$\mathcal{C}'$ 위의 환의 층 $\mathcal{O}'$, 그리고 사상
$g^\sharp : g^{-1}\mathcal{O} \to \mathcal{O}'$가 주어졌다고 하자.
이에 대응하는 환 달린 토포스의 사상을 간단히
$g : (\Sh(\mathcal{C}'), \mathcal{O}') \to (\Sh(\mathcal{C}), \mathcal{O})$,
Modules on Sites, 절 \ref{sites-modules-section-ringed-topoi} 참조).
마지막으로 $(\mathcal{A}, \text{d})$가 미분 등급
$\mathcal{O}$-대수의 층이고 $(\mathcal{A}', \text{d})$가 미분 등급
$\mathcal{O}'$-대수의 층이며, 미분 등급 $\mathcal{O}'$-대수의 사상
$\varphi : g^*\mathcal{A} \to \mathcal{A}'$가 주어졌다고 하자
(절 \ref{sdga-section-functoriality-dg} 참조).

\begin{lemma}
\label{sdga-lemma-cartesian-functorial}
 $g : (\Sh(\mathcal{C}'), \mathcal{O}') \to (\Sh(\mathcal{C}), \mathcal{O})$
와 $\varphi : g^*\mathcal{A} \to \mathcal{A}'$가 위와 같다고 하자.
그러면 함자
$Lg^* : D(\mathcal{A}, \text{d}) \to D(\mathcal{A}', \text{d})$
는 $\mathit{QC}(\mathcal{A}, \text{d})$를
$\mathit{QC}(\mathcal{A}', \text{d})$로 보낸다.
\end{lemma}

\begin{proof}
$U' \in \Ob(\mathcal{C}')$를 택하고 그 상을
$U = u(U')$가 $\mathcal{C}$에서의 상이라고 하자.
$pt$를 하나의 대상과 하나의 사상만을 갖는 범주라고 하자.
$(\Sh(pt), \mathcal{O}'(U'))$와 $(\Sh(pt), \mathcal{O}(U))$를
위와 같이 미분 등급 대수 $\mathcal{A}'(U)$와 $\mathcal{A}(U)$를 갖춘
환 달린 토포스로 두자.
이들의 미분 등급 가군 유도 범주는 각각
$D(\mathcal{A}'(U'), \text{d})$와 $D(\mathcal{A}(U), \text{d})$로
식별한다. 그러면 환 달린 토포스의 다음 가환 도표가 있다.
$$
\xymatrix{
(\Sh(pt), \mathcal{O}'(U')) \ar[rr]_{U'} \ar[d] & &
(\Sh(\mathcal{C}'), \mathcal{O}') \ar[d]^g \\
(\Sh(pt), \mathcal{O}(U)) \ar[rr]^U & &
(\Sh(\mathcal{C}), \mathcal{O})
}
$$
각 토포스에는 대응하는 미분 등급 대수가 주어져 있다.
아래 수평 사상에 따른 당김은 $D(\mathcal{A}, \text{d})$의 $M$을
$D(\mathcal{A}(U), \text{d})$의 대상으로 본 $R\Gamma(U, K)$로 보낸다.
왼쪽 수직 사상에 따른 당김은 $M$을
$M \otimes_{\mathcal{A}(U)}^\mathbf{L} \mathcal{A}'(U')$.
도표를 어느 방향으로 돌아도 같은 결과를 주므로
(보조정리 \ref{sdga-lemma-compose-pullback}) 다음을 얻는다.
$$
R\Gamma(U', Lg^*K) =
R\Gamma(U, K) \otimes_{\mathcal{A}(U)}^\mathbf{L} \mathcal{A}'(U')
$$
마지막으로 $f' : U' \to V'$를 $\mathcal{C}'$의 사상이라 하고,
그 상을 $\mathcal{C}$에서의
$f = u(f') : U = u(U') \to V = u(V')$로 쓰자.
$K$가 $\mathit{QC}(\mathcal{A}, \text{d})$에 속하면 다음을 얻는다.
\begin{align*}
R\Gamma(V', Lg^*K) \otimes_{\mathcal{A}'(V')}^\mathbf{L} \mathcal{A}'(U')
& =
R\Gamma(V, K) \otimes_{\mathcal{A}(V)}^\mathbf{L} \mathcal{A}'(V')
\otimes_{\mathcal{A}'(V')}^\mathbf{L} \mathcal{A}'(U') \\
& =
R\Gamma(V, K) \otimes_{\mathcal{A}(V)}^\mathbf{L} \mathcal{A}'(U') \\
& =
R\Gamma(V, K) \otimes_{\mathcal{A}(V)}^\mathbf{L} \mathcal{A}(U)
\otimes_{\mathcal{A}(U)}^\mathbf{L} \mathcal{A}'(U') \\
& =
R\Gamma(U, K) \otimes_{\mathcal{A}(U)}^\mathbf{L} \mathcal{A}'(U') \\
& =
R\Gamma(U', Lg^*K)
\end{align*}
원하는 결론이다. 여기서는 위의 관찰을 $U'$와 $V'$ 모두에 사용했다.
\end{proof}






\section{범주 위의 미분 등급 가군, 계속}
\label{sdga-section-modules-cohomology-bis}

\noindent
절 \ref{sdga-section-modules-cohomology}에서 도입한 준연접 가군의
개념에 대해 몇 가지 결과를 더 전개한다.

\begin{lemma}
\label{sdga-lemma-cartesian-cofinal-subcategory}
 $\mathcal{C}, \mathcal{O}, \mathcal{A}$가 절
\ref{sdga-section-modules-cohomology}에서와 같다고 하자. $\mathcal{C}'
\subset \mathcal{C}$를 다음 성질을 갖는 충만 부분범주라고 하자.
모든 $U \in \Ob(\mathcal{C})$에 대해 화살표 $U \to U'$들의 범주
$U/\mathcal{C}'$가 여과공이다.
$\mathcal{O}', \mathcal{A}'$를 $\mathcal{C}'$에 대한
$\mathcal{O}, \mathcal{A}$의 제한이라고 하자.
그러면 제한은 다음 동치를 유도한다.
$\mathit{QC}(\mathcal{A}, \text{d}) \to \mathit{QC}(\mathcal{A}', \text{d})$.
\end{lemma}

\begin{proof}
이 함자의 준역함자를 구성하자. 즉
$M'$를 $\mathit{QC}(\mathcal{A}', \text{d})$의 대상이라고 하자.
보조정리 \ref{sdga-lemma-resolve}에 따라 $M'$를 좋은 미분 등급 가군
$\mathcal{M}'$으로 나타낼 수 있다. 그러면 모든
$U' \in \Ob(\mathcal{C}')$에 대해 미분 등급
$\mathcal{A}'(U')$-가군 $\mathcal{M}'(U)$는 K-평탄하고 등급 평탄하다.
또한 $\mathcal{C}'$의 모든 사상 $U'_1 \to U'_2$에 대해 다음 사상은
$$
\mathcal{M}'(U'_2) \otimes_{\mathcal{A}'(U'_2)} \mathcal{A}'(U'_1)
\longrightarrow
\mathcal{M}'(U'_1)
$$
유사동형이다 (그 원천이 유도 텐서곱을 나타내기 때문이다).
다음 규칙으로 정의되는 미분 등급 $\mathcal{A}$-가군 $\mathcal{M}$을 생각하자.
$$
\mathcal{M}(U) =
\colim_{U \to U' \in U/\mathcal{C}'}
\mathcal{M}'(U') \otimes_{\mathcal{A}'(U')} \mathcal{A}(U)
$$
이는 보조정리의 가정에 의해 복합체들의 여과 직결한계이다.
$M'$가 $\mathit{QC}(\mathcal{A}', \text{d})$에 속하므로 계의
모든 전이 사상은 유사동형이다. 여과 직결한계는 정확하므로,
$U' \in \Ob(\mathcal{C}')$인 임의의 사상 $U \to U'$에 대해
$D(\mathcal{A}(U), \text{d})$에서 $\mathcal{M}(U)$는
$\mathcal{M}'(U') \otimes_{\mathcal{A}'(U')} \mathcal{A}(U)$와 동형이다.

\medskip\noindent
 $\mathcal{M}$이 $\mathit{QC}(\mathcal{A}, \text{d})$에 속함을 보이자.
즉 $\mathcal{C}$에서 $U \to V$가 주어지면
$V' \in \Ob(\mathcal{C}')$인 사상 $V \to V'$를 택한다.
위 결과에 의해 사상 $\mathcal{M}(V) \to \mathcal{M}(U)$는 다음 사상으로
식별된다.
$$
\mathcal{M}'(V') \otimes_{\mathcal{A}'(V')} \mathcal{A}(V)
\longrightarrow
\mathcal{M}'(V') \otimes_{\mathcal{A}'(V')} \mathcal{A}(U)
$$
 $\mathcal{M'}(V')$는 미분 등급 $\mathcal{A}'(V')$-가군으로서
K-평탄하므로 주장이 성립한다.

\medskip\noindent
자연 사상 $\mathcal{M}|_{\mathcal{C}'} \to \mathcal{M}'$는 위에서 즉시
따르는 대로 $D(\mathcal{A}', d)$에서 동형이다.

\medskip\noindent
반대로 $\mathit{QC}(\mathcal{A}, \text{d})$의 대상 $E$가 주어지면
이를 좋은 미분 등급 가군 $\mathcal{E}$로 나타내자.
$\mathcal{M}' = \mathcal{E}|_{\mathcal{C}'}$로 두면 (이는 또 다른 좋은
미분 등급 가군이다) 다음 사상이 존재한다.
$$
\mathcal{E} \to \mathcal{M}
$$
이는 $\mathcal{C}$의 $U$ 위에서 다음 사상으로 주어진다.
$$
\mathcal{E}(U)
\longrightarrow
\colim_{U \to U' \in U/\mathcal{C}'}
\mathcal{E}(U') \otimes_{\mathcal{A}'(U')} \mathcal{A}(U)
$$
같은 이유로 이는 유사동형이다. 따라서 제한과 위 구성은 원하는 대로
서로 준역함자이다.
\end{proof}

\begin{lemma}
\label{sdga-lemma-cartesian-qis-equivalence}
 $\mathcal{C}, \mathcal{O}$가 절
\ref{sdga-section-modules-cohomology}에서와 같다고 하자. $\varphi :
\mathcal{A} \to \mathcal{B}$를 코호몰로지 층에서 동형을 유도하는
미분 등급 $\mathcal{O}$-대수의 준동형이라고 하자. 그러면
다음 동치
$D(\mathcal{A}, \text{d}) \to D(\mathcal{B}, \text{d})$
보조정리 \ref{sdga-lemma-qis-equivalence}가 다음 동치를 유도한다.
$\mathit{QC}(\mathcal{A}, \text{d}) \to
\mathit{QC}(\mathcal{B}, \text{d})$.
\end{lemma}

\begin{proof}
다음을 보이면 충분하다. $\mathcal{C}$의 사상 $U \to V$와
$D(\mathcal{A}, \text{d})$의 $M$이 주어졌을 때 다음 두 조건은 동치이다.
\begin{enumerate}
\item $R\Gamma(V, M) \otimes_{\mathcal{A}(V)}^\mathbf{L} \mathcal{A}(U)
\to \Gamma(U, M)$는 $D(\mathcal{A}(U), \text{d})$에서 동형이다.
\item $R\Gamma(V, M \otimes_\mathcal{A}^\mathbf{L} \mathcal{B})
\otimes_{\mathcal{B}(V)}^\mathbf{L} \mathcal{B}(U)
\to \Gamma(U, M \otimes_\mathcal{A}^\mathbf{L} \mathcal{B})$
는 $D(\mathcal{B}(U), \text{d})$에서 동형이다.
\end{enumerate}
 $\mathcal{C}$의 위상이 혼돈 위상이므로 이는 단순히
$\mathcal{A}(U) \to \mathcal{B}(U)$와
$\mathcal{A}(V) \to \mathcal{B}(V)$가 유사동형이라는 사실로 귀결된다.
세부사항은 생략한다.
\end{proof}






\section{미분 등급 대수의 역계}
\label{sdga-section-qc-inverse-system}

\noindent
이 절에서는 절 \ref{sdga-section-modules-cohomology}에서 논의한 상황의
다음 특수한 경우를 생각한다.
\begin{enumerate}
\item $\mathcal{C}$는 범주 $\mathbf{N}$이며,
$i \leq j$일 때 그리고 그때에만 $i \to j$라는 유일한 사상이 있다.
\item $\mathcal{O}$는 주어진 환 $R$을 값으로 갖는 상수 (프리)층이다.
\end{enumerate}
이 설정에서 미분 등급 $\mathcal{O}$-대수의 층 $\mathcal{A}$는
미분 등급 $R$-대수의 역계 $(A_n)$와 같다.
미분 등급 $\mathcal{A}$-가군의 층 $\mathcal{M}$은,
$M_n$이 미분 등급 $A_n$-가군이고 전이 사상
$M_{n + 1} \to M_n$이 $A_{n + 1}$-가군 사상인 역계 $(M_n)$와 같다.
이 절에서는 표기에서 미분을 생략하고,
정의 \ref{sdga-definition-cartesian}에서
$\mathit{QC}(\mathcal{A}, \text{d})$로 표시한 범주를
$\mathit{QC}(\mathcal{A})$로 쓰겠다.

\medskip\noindent
 $\mathcal{B} = (B_n)$가 두 번째 미분 등급 $R$-대수 역계라고 하자.
프로 대상의 사상 $\varphi : (A_n) \to (B_n)$이 주어지면
$\mathit{QC}(\mathcal{A})$에서 $\mathit{QC}(\mathcal{B})$로 가는
정확한 함자를 구성할 것이다. 구체적으로 Categories, 예
\ref{categories-example-pro-morphism-inverse-systems}
사상 $\varphi$는 정수의 수열
$\ldots \geq m(3) \geq m(2) \geq m(1)$과 다음 가환 도표로 주어진다.
$$
\xymatrix{
\ldots \ar[r] &
A_{m(3)} \ar[d]^{\varphi_3} \ar[r] &
A_{m(2)} \ar[d]^{\varphi_2} \ar[r] &
A_{m(1)} \ar[d]^{\varphi_1} \\
\ldots \ar[r] &
B_3 \ar[r] &
B_2 \ar[r] &
B_1
}
$$
미분 등급 $R$-대수의 도표이다. $\mathit{QC}(\mathcal{A})$의
대상을 나타내는 좋은 미분 등급 $\mathcal{A}$-가군의 층
$\mathcal{M} = (M_n)$이 주어지면 다음과 같이 둘 수 있다.
$$
N_n = M_{m(n)} \otimes_{A_{m(n)}} B_n
$$
이 역계는 $\mathit{QC}(\mathcal{B})$의 대상을 정한다.
$A_{m(n)}$-가군 $M_{m(n)}$이 K-평탄하기 때문이다. 세부사항은 생략한다.
구성에서 한 선택과 무관하게 결과 함자가 정해짐을 보이는 일은
독자에게 맡긴다.

\begin{lemma}
\label{sdga-lemma-pro-isomorphic}
위 상황에서 $\mathcal{A} = (A_n)$와
$\mathcal{B} = (B_n)$가 미분 등급 $R$-대수의 역계라고 하자.
$\varphi : (A_n) \to (B_n)$가 프로 대상의 동형이면, 위에서 구성한
함자 $\mathit{QC}(\mathcal{A}) \to \mathit{QC}(\mathcal{B})$는 동치이다.
\end{lemma}

\begin{proof}
 $\psi : (B_n) \to (A_n)$를 $\varphi$의 역인 프로 대상의 사상이라고 하자.
앞의 Categories, 예
\ref{categories-example-pro-morphism-inverse-systems}
에 따라 $\varphi$가 위와 같은 사상들의 계로 주어지고 $\psi$가
$n(1) < n(2) < \ldots$와 다음 가환 도표로 주어진다고 가정할 수 있다.
$$
\xymatrix{
\ldots \ar[r] &
B_{n(3)} \ar[d]^{\psi_3} \ar[r] &
B_{n(2)} \ar[d]^{\psi_2} \ar[r] &
B_{n(1)} \ar[d]^{\psi_1} \\
\ldots \ar[r] &
A_3 \ar[r] &
A_2 \ar[r] &
A_1
}
$$
미분 등급 $R$-대수의 도표이다. $\varphi \circ \psi = \text{id}$이므로,
함수 $n(\cdot)$와 $m(\cdot)$의 값을 필요하면 증가시켜
$B_{n(m(i))} \to A_{m(i)} \to B_i$가 전이 사상이 되게 가정할 수 있다.
따라서 다음 함자들의 합성
$$
\mathit{QC}(\mathcal{B}) \to
\mathit{QC}(\mathcal{A}) \to
\mathit{QC}(\mathcal{B})
$$
은 좋은 미분 등급 $\mathcal{B}$-가군의 층
$\mathcal{N} = (N_n)$을 다음 값을 갖는 역계
$\mathcal{N}' = (N'_i)$로 보낸다.
$$
N'_i = N_{n(m(i))} \otimes_{B_{n(m(i))}} B_i
$$
이는 $\mathcal{N}$이 $\mathit{QC}(\mathcal{B})$의 대상이고
모든 $j$에 대해 $N_j$가 K-평탄 미분 등급 가군이기 때문에
표준적으로 $\mathcal{N}$과 유사동형이다.
반대 방향의 합성에도 같은 사실이 성립하므로 결론을 얻는다.
\end{proof}

\noindent
$\mathcal{C} = \mathbf{N}$이고 $\mathcal{O}$가 환 $R$을 값으로 갖는
상수층이며, $\mathcal{A}$가 미분 등급 $R$-대수의 역계 $(A_n)$으로
주어진다고 하자. 역계 $(N_n)$과 $(N'_n)$으로 주어진 두 왼쪽
미분 등급 $\mathcal{A}$-가군 $\mathcal{N}$과 $\mathcal{N}'$을 생각하자.
따라서 각 $N_n$과 $N'_n$은 왼쪽 미분 등급 $A_n$-가군이다.
다음 정수 수열과 사상들이 존재할 때 $(N_n)$과 $(N'_n)$이
{it 유도 범주에서 프로 동형}이라고 잠시 부르자.
$$
1 = n_0 < n_1 < n_2 < n_3 < \ldots
$$
그리고 사상
$$
N_{n_{2i}} \to N'_{n_{2i - 1}}
\quad\text{in}\quad
D(A_{n_{2i}}^{opp}, \text{d})
$$
및
$$
N'_{n_{2i + 1}} \to N'_{n_{2i}}
\quad\text{in}\quad
D(A_{n_{2i + 1}}^{opp}, \text{d})
$$
이때 합성 $N_{n_{2i}} \to N_{n_{2i - 2}}$와
$N'_{n_{2i + 1}} \to N'_{2i - 1}$은 각각의 계의 전이 사상으로 주어진다.

\begin{lemma}
\label{sdga-lemma-pro-isomorphic-systems-tensor}
$(N_n)$과 $(N'_n)$이 위에서 정의한 의미로 유도 범주에서 프로 동형이면,
$D(\mathbf{N}, \mathcal{A})$의 모든 대상 $(M_n)$에 대해 다음이 성립한다.
$$
R\lim (M_n \otimes_{A_n}^\mathbf{L} N_n) =
R\lim (M_n \otimes_{A_n}^\mathbf{L} N'_n)
$$
 $D(R)$에서.
\end{lemma}

\begin{proof}
가정에 의해 $D(R)$의 역계
$(M_n \otimes_{A_n}^\mathbf{L} N_n)$은 통상적인 의미에서
$D(R)$의 역계 $(M_n \otimes_{A_n}^\mathbf{L} N'_n)$과 프로 동형이다.
따라서 결과는 More on Algebra, 보조정리
\ref{more-algebra-lemma-pro-isomorphism-Rlim}에서 따른다.
\end{proof}

\begin{lemma}
\label{sdga-lemma-different-tensors}
\begin{reference}
이는 \cite[Lemma 3.5.4]{BS}의 변형이다.
\end{reference}
$R$을 환이라 하고 $f_1, \ldots, f_r \in R$라 하자.
$K_n$을 $f_1^n, \ldots, f_r^n$에 대한 코쥘 복합체를
미분 등급 $R$-대수로 본 것이라 하자.
$(M_n)$을 $D(\mathbf{N}, (K_n))$의 대상이라고 하자.
그러면 임의의 $t \geq 1$에 대해 다음이 성립한다.
$$
R\lim (M_n \otimes_R^\mathbf{L} K_t) =
R\lim (M_n \otimes_{K_n}^\mathbf{L} K_t)
$$
 $D(R)$에서.
\end{lemma}

\begin{proof}
 $t \geq 1$을 고정하자. $n \geq t$에 대해 미분 등급 $R$-대수
$K_t$를 왼쪽 미분 등급 $K_n$-가군으로 본 것을 ${}_nK_t$로 쓰자.
다음을 관찰하자.
$$
M_n \otimes_R^\mathbf{L} K_t =
M_n \otimes_{K_n}^\mathbf{L} (K_n \otimes_R^\mathbf{L} K_t) =
M_n \otimes_{K_n}^\mathbf{L} (K_n \otimes_R K_t)
$$
따라서 보조정리 \ref{sdga-lemma-pro-isomorphic-systems-tensor}에 의해
$({}_nK_t)$와 $(K_n \otimes_R K_t)$가 유도 범주에서 프로 동형임을
보이면 충분하다. 곱셈 사상
$$
K_n \otimes_R K_t \longrightarrow {}_nK_t
$$
은 왼쪽 미분 등급 $K_n$-가군의 사상이다. 따라서 증명을 마치려면
모든 $n \geq 1$에 대해 $N > n$과 다음 사상이 존재하여
$$
{}_NK_t \longrightarrow {}_NK_n \otimes_R K_t
$$
 $D(K_N^{opp}, \text{d})$에서 이 사상과 곱셈 사상의 합성이
어느 방향에서든 전이 사상이 됨을 보이면 충분하다.
이는 Divided Power Algebra, 보조정리
\ref{dpa-lemma-construct-some-maps}에서 명시적 구성으로 행해진다.
\end{proof}

\begin{proposition}
\label{sdga-proposition-derived-complete}
$R$을 뇌터 환이라 하고 $I \subset R$을 아이디얼이라 하자.
다음 세 범주는 표준적으로 동치이다.
\begin{enumerate}
\item $\mathcal{A}$를 $R$-대수의 역계 $A_n = R/I^n$에 대응하는
$\mathbf{N}$ 위의 $R$-대수의 층이라고 하자. 범주
$\mathit{QC}(\mathcal{A})$.
\item $I$의 생성원 $f_1, \ldots, f_r$을 택하자.
$\mathcal{B}$를 $f_1^n, \ldots, f_r^n$에 대한 코쥘 대수의 역계에
대응하는 $\mathbf{N}$ 위의 미분 등급 $R$-대수의 층이라고 하자.
범주 $\mathit{QC}(\mathcal{B})$.
\item 유도 완비 대상들로 이루어진 충만 부분범주
$D_{comp}(R, I) \subset D(R)$ (More on Algebra, 정의
\ref{more-algebra-definition-derived-complete} 및 그 뒤의 글 참조).
\end{enumerate}
\end{proposition}

\begin{proof}
환 달린 토포스의 자명한 사상
$f : (\Sh(\mathbf{N}), \mathcal{A}) \to (\Sh(pt), R)$
을 생각하고 수반 함자 $Lf^*$와 $Rf_*$를 생각하자.
첫 번째 함자는 다음 함자로 제한된다.
$$
F : D_{comp}(R, I) \longrightarrow \mathit{QC}(\mathcal{A})
$$
이는 K-평탄 복합체 $K^\bullet$로 나타낸 $D_{comp}(R, I)$의 대상
$K$를 $\mathit{QC}(\mathcal{A})$의 대상
$(K^\bullet \otimes_R R/I^n)$으로 보낸다. 두 번째 함자는 다음 함자로
제한된다.
$$
G : \mathit{QC}(\mathcal{A}) \longrightarrow D_{comp}(R, I)
$$
이는 $\mathit{QC}(\mathcal{A})$의 대상 $(M_n^\bullet)$을
$R\lim M_n^\bullet$로 보낸다. 그 결과는 예를 들어 More on Algebra,
보조정리 \ref{more-algebra-lemma-naive-derived-completion}에 의해 유도 완비이다.
또한 More on Algebra, 명제
\ref{more-algebra-proposition-noetherian-naive-completion-is-completion}
에서 $G \circ F = \text{id}$가 따른다. 따라서 $F$와 $G$가 서로
준역인 동치임을 보이려면 $G$의 핵이 영임을 보이면 충분하다
(Derived Categories, 보조정리
\ref{derived-lemma-fully-faithful-adjoint-kernel-zero} 참조).
그러나 이를 직접 보이는 것은 쉽지 않아 보인다!

\medskip\noindent
이 문단에서는 $\mathit{QC}(\mathcal{A})$와
$\mathit{QC}(\mathcal{B})$가 동치임을 보인다.
$\mathcal{B} = (B_n)$로 쓰자. 여기서 $B_n$은 코쥘 복합체를
차수 $-r, -r + 1, \ldots, 0$의 코체인 복합체로 본 것이다.
Divided Power Algebra, 비고 \ref{dpa-remark-pro-system-koszul}에 의해
(단, 체인 복합체를 코체인 복합체로 바꾸어)
$1 < n_1 < n_2 < \ldots$와 미분 등급 $R$-대수의 사상
$B_{n_i} \to E_i \to R/(f_1^{n_i}, \ldots, f_r^{n_i})$ 및
$E_i \to B_{n_{i - 1}}$를 찾아서 다음이 성립하게 할 수 있다.
$$
\xymatrix{
B_{n_1} \ar[d] &
B_{n_2} \ar[d] \ar[l] &
B_{n_3} \ar[d] \ar[l] &
\ldots \ar[l] \\
E_1 \ar[d] &
E_2 \ar[l] \ar[d] &
E_3 \ar[l] \ar[d] &
\ldots \ar[l] \\
B_1 &
B_{n_1} \ar[l] &
B_{n_2} \ar[l] &
\ldots \ar[l]
}
$$
이는 미분 등급 $R$-대수의 가환 도표이고,
$E_i \to R/(f_1^{n_i}, \ldots, f_r^{n_i})$는 유사동형이다.
따라서 다음을 얻는다.
\begin{enumerate}
\item $\mathit{QC}(\mathcal{B})$와 $\mathit{QC}((E_i))$ 사이에 동치가 있다.
\item $\mathit{QC}((E_i))$와
$\mathit{QC}((R/(f_1^{n_i}, \ldots, f_r^{n_i})))$ 사이에 동치가 있다.
\item $\mathit{QC}((R/(f_1^{n_i}, \ldots, f_r^{n_i})))$와
$\mathit{QC}(\mathcal{A})$ 사이에 동치가 있다.
\end{enumerate}
실제로 (1)에 대해서는 위 도표에 보조정리 \ref{sdga-lemma-pro-isomorphic}을
적용할 수 있으며, 이 도표는 $(E_i)$와 $(B_n)$이 프로 동형임을 보인다.
(2)에 대해서는 유사동형의 역계
$E_i \to R/(f_1^{n_i}, \ldots, f_r^{n_i})$에 보조정리
\ref{sdga-lemma-cartesian-qis-equivalence}를 적용할 수 있다.
(3)에 대해서는 보조정리 \ref{sdga-lemma-pro-isomorphic}과 역계
$(R/I^n)$ 및 $(R/(f_1^{n_i}, \ldots, f_r^{n_i})$가
프로 동형이라는 초등적 사실을 적용할 수 있다.

\medskip\noindent
증명의 첫 문단에서와 정확히 같이 수반 함자를 정의할 수 있다.\footnote{
이 함자들이 위 동치들을 통해 앞서 정의한 $F$ 및 $G$와 양립함을
보일 수 있다.}
$$
F' : D_{comp}(R, I) \longrightarrow \mathit{QC}(\mathcal{B})
\quad\text{and}\quad
G' : \mathit{QC}(\mathcal{B}) \longrightarrow D_{comp}(R, I).
$$
첫 번째 함자는 K-평탄 복합체 $K^\bullet$로 나타낸
$D_{comp}(R, I)$의 대상 $K$를 $\mathit{QC}(\mathcal{B})$의 대상
$(K^\bullet \otimes_R B_n)$으로 보낸다. 두 번째 함자는
$\mathit{QC}(\mathcal{B})$의 대상 $(M_n)$을 $R\lim M_n$으로 보낸다.
위와 같이 논증하면 $G'$의 핵이 영임을 보이면 충분하다.
그러므로 $\mathcal{M} = (M_n)$을 $G'$의 핵에 있는
$\mathit{QC}(\mathcal{B})$의 대상을 나타내는, $\mathcal{B}$ 위의
좋은 미분 등급 가군의 층이라고 하자. 그러면
$$
0 = R\lim M_n \Rightarrow
0 = (R\lim M_n) \otimes_R^\mathbf{L} B_t =
R\lim (M_n \otimes_R^\mathbf{L} B_t)
$$
보조정리 \ref{sdga-lemma-different-tensors}에 의해
$R\lim (M_n \otimes_R^\mathbf{L} B_t) =
R\lim (M_n \otimes_{B_n}^\mathbf{L} B_t)$.
$(M_n)$이 $\mathit{QC}(\mathcal{B})$의 대상이므로 역계
$M_n \otimes_{B_n}^\mathbf{L} B_t$는 결국 $M_t$라는 값으로 상수이다.
따라서 원하는 대로 $M_t = 0$이다.
\end{proof}

\begin{remark}
\label{sdga-remark-proregular}
$R$을 환이라 하고, $f_1, \ldots, f_r \in R$을 아이디얼 $I$를 생성하는
원소들의 수열이라고 하자. $K_n$을 $f_1^n, \ldots, f_r^n$에 대한
코쥘 복합체를 미분 등급 $R$-대수로 본 것이라 하자.
모든 $n$에 대해 $m > n$이 존재하여 $K_m \to K_n$이 차수 $0$을
제외한 코호몰로지에서 영사상을 유도하면 $f_1, \ldots, f_r$을
{\it 약한 프로정칙 수열}이라고 한다. 그러면 $R$이 뇌터 환이 아니어도
명제 \ref{sdga-proposition-derived-complete}의 증명에 있는 논증들이 계속 성립한다.
특히 $I$가 약한 프로정칙 수열로 생성될 수 있으면
$\mathit{QC}(\{R/I^n\})$는 $R$-선형 삼각 범주로서 유도 완비 대상들의
범주 $D_{comp}(R, I)$와 동치이다.
필요해지는 경우에는 여기서 이를 정확히 진술하고 증명할 것이다.
\end{remark}
